%~Mouliné par MaN_auto v.0.32.0 2024-07-01 09:54:38
\documentclass[AHL,Unicode,longabstracts,published]{cedram}

\usepackage{graphicx}
\usepackage{tikz-cd}
\usepackage[all]{xy}
\usepackage{comment}



\newcommand{\GG}{\mathcal{G}}
\newcommand{\FF}{\mathcal{F}}
\newcommand{\Tau}{\mathcal{T}}
\newcommand{\LL}{\mathcal{L}}
\newcommand{\clR}{\mathcal{R}}
\newcommand{\clG}{\mathcal{G}}
\newcommand{\MM}{\mathcal{M}}
\newcommand{\B}{\mathcal{B}}

\newcommand{\Z}{\mathbb{Z}}
\newcommand{\R}{\mathbb{R}}
\newcommand{\T}{\mathbb{T}}
\newcommand{\Prob}{\mathbb{P}}

\newcommand{\GL}{\mathrm{GL}}
\newcommand{\id}{\mathrm{id}}
\newcommand{\Diff}{\mathrm{Diff}}
\newcommand{\width}{\mathrm{width}}
\newcommand{\length}{\mathrm{length}}
\newcommand{\rspan}{\mathrm{span}}
\newcommand{\rst}{\mathrm{st}}
\newcommand{\bh}{\mathfrak{h}}
\newcommand{\bu}{\mathfrak{u}}
\newcommand{\bv}{\mathfrak{v}}


\newcommand{\ie}{i.e.}



%%%%%%%%%%%%%%%%%%%%%%%%%%%%%%%%%%%%%%%%%%%%%%%%%%%%%%%%%%%%%%%%%%%%%%%%%%%%%%%%%%
\graphicspath{{./figures/}}

\newcommand*{\mk}{\mkern -1mu}
\newcommand*{\Mk}{\mkern -2mu}
\newcommand*{\mK}{\mkern 1mu}
\newcommand*{\MK}{\mkern 2mu}

%\hypersetup{urlcolor=purple, linkcolor=blue, citecolor=red}

\newcommand*{\romanenumi}{\renewcommand*{\theenumi}{\roman{enumi}}}
\newcommand*{\Romanenumi}{\renewcommand*{\theenumi}{\Roman{enumi}}}
\newcommand*{\alphenumi}{\renewcommand*{\theenumi}{\alph{enumi}}}
\newcommand*{\Alphenumi}{\renewcommand*{\theenumi}{\Alph{enumi}}}
\let\oldtilde\tilde
\renewcommand*{\tilde}[1]{\mathchoice{\widetilde{#1}}{\widetilde{#1}}{\oldtilde{#1}}{\oldtilde{#1}}}
\let\oldhat\hat
\renewcommand*{\hat}[1]{\mathchoice{\widehat{#1}}{\widehat{#1}}{\oldhat{#1}}{\oldhat{#1}}}
%%%%%%%%%%%%%%%%%%%%%%%%%%%%%%%%%%%%%%%%%%%%%%%%%%%%%%%%%%%%%%%%%%%%%%%%%%%%%%%%%%%%%%%%%%

\title[Deforming Finsler two-tori]{Deforming a Finsler metric on the two-torus to a flat Finsler metric with conjugate geodesic flows}

\alttitle{Déformation d'une métrique de Finsler sur le tore bidimensionnel en une métrique de Finsler plate avec des flots géodésiques conjugués}

\subjclass[2020] {37J39, 53C22, 53C65}

\keywords{Finsler metrics, dynamical systems, geodesic flow, conjugate flows, conjugate points, integral geometry, Crofton formula, Heber foliation, curve shortening flow}



\author[\initial{E.} \lastname{Nakhlé}]{\firstname{Elie} \lastname{Nakhlé}}
\address{Univ Paris Est Creteil,\\
CNRS, LAMA,\\
F-94010 Creteil (France)\\
Univ Gustave Eiffel, LAMA,\\
F-77447 Marne-la-Vallée (France)}
\email{elie.nakhle@u-pec.fr}
\thanks{Partially supported by the ANR project Min-Max (ANR-19-CE40-0014).}

\author[\initial{S.} \lastname{Sabourau}]{\firstname{Stéphane} \lastname{Sabourau}}
\address{Univ Paris Est Creteil,\\ 
CNRS, LAMA,\\
F-94010 Creteil (France)\\
Univ Gustave Eiffel, LAMA,\\
F-77447 Marne-la-Vallée (France)}
\email{stephane.sabourau@u-pec.fr}


\begin{abstract}
We show that the space of (reversible) Finsler metrics on the two-torus~$\mathbb{T}^2$ whose geodesic flow is conjugate to the geodesic flow of a flat Finsler metric strongly deformation retracts to the space of flat Finsler metrics with respect to the uniform convergence topology. Along the proof, we also show that two Finsler metrics on~$\mathbb{T}^2$ without conjugate points, whose Heber foliations are smooth and with the same marked length spectrum, have conjugate geodesic flows.
\end{abstract}

\begin{altabstract}
Nous montrons que l'espace des métriques de Finsler (réversibles) sur le tore bidimensionnel~$\mathbb{T}^2$, dont le flot géodésique est conjugué au flot géodésique d'une métrique de Finsler plate se rétracte par déformation forte sur l'espace des métriques de Finsler plates par rapport à la topologie de la convergence uniforme. Au cours de la preuve, nous montrons également que deux métriques de Finsler
sur~$\mathbb{T}^2$ sans points conjugués, dont les feuilletages d'Heber sont lisses et ont le même spectre marqué des longueurs, ont des flots géodésiques conjugués.
\end{altabstract}


\datereceived{2023-08-30}
\dateaccepted{2024-05-28}

\editors{X. Caruso and V. Colin}
%%%%%%%%%%%%%%%%%%%%%%%%%%%%%%%%%%%%%%%%%%%%%%%%%%%%%%%%%%%%%%%%%%%%%%%%%%%%%%%%%
\dateposted{2024-09-05}
\begin{document}
\maketitle

\section{Introduction}

The goal of this article is to study the space of Finsler two-tori whose geodesic flow is (smoothly) conjugate to the geodesic flow of a flat Finsler two-torus. (By definition, all Finsler metrics are reversible and quadratically convex; see Definition~\ref{def:finsler}). In particular, we will determine the topology of this space up to homotopy with respect to the uniform convergence topology; see Corollary~\ref{coro:main}.


This problem can be seen as the counterpart of a famous question about Zoll metrics (\ie, metrics all of whose geodesics are simple closed curves of the same length) asking whether the space of Zoll metrics, say in the two-sphere, is path-connected; see~\cite[Question~200]{B03}. (Though this question is wide open, it has a positive answer for Zoll Finsler metrics on the projective plane; see~\cite{Sa19}.) Since the geodesic flow of a Zoll metric on the two-sphere is conjugate to that of the canonical metric, see~\cite[\S~4F]{Be78} or~\cite{ABHS17}, this question amounts to asking whether the space of metrics whose geodesic flow is conjugate to that of the round two-sphere is path-connected. Clearly, the two-torus does not admit any Zoll metric, but the latter formulation of the problem makes sense if one replaces the round two-sphere with a flat two-torus.


Complete Finsler manifolds without conjugate points have for characteristic property that any pair of points in their universal cover can be joined by a unique geodesic; see Definition~\ref{def:ncp}. There are strong ties between Finsler metrics whose geodesic flow is conjugate to the geodesic flow of a flat Finsler two-torus and metrics without conjugate points. Indeed, every Finsler metric~$F$ on the two-torus whose geodesic flow is conjugate to that of a flat Finsler metric~$F_\diamond$ has no conjugate points; see~\cite{cro90} or Lemma~\ref{lem:ncp}. For a Riemannian metric~$F$, this implies that $F$ is flat by Hopf's theorem~\cite{Ho48} (see~\cite{BI94} for a generalization to any dimension), in which case, the metrics~$F$ and~$F_\diamond$ coincide. The problem is therefore interesting only for non-Riemannian Finsler metrics. It is an open question whether the geodesic flow of every Finsler two-torus without conjugate points is conjugate to that of some flat Finsler metric; see~\cite{BC21,CK94}. The issue is related to the regularity of the so-called Heber foliation (a continuous foliation of~$T^*\T^2$ by Lipschitz, Lagrangian, flow-invariant graphs given by the covectors of the same norm generating geodesics with the same asymptotic direction); see Section~\ref{sec:foliation}. Actually, it can be proved that two Finsler two-tori without conjugate points having the same marked length spectrum and smooth Heber foliations have conjugate geodesic flows; see Section~\ref{def:conjugacy}. To further illustrate our poor understanding of Finsler metrics without conjugate points, let us also mention that it is still an open question whether a geodesic with irrational direction on a Finsler two-torus without conjugate points is dense; see~\cite{BC21}.


Before stating our main result, it should be noted that Finsler two-tori without conjugate points are incredibly flexible: given any point on a Finsler surface, there exists a neighborhood of this point which isometrically embeds into a Finsler two-torus without conjugate points; see~\cite{Ch19}. In particular, modulo isometries and rescaling, Finsler two-tori without conjugate points form a infinite-dimensional space. Similarly, the space of Finsler two-tori with geodesic flow conjugate to that of a flat Finsler metric (modulo isometries) has infinite dimension; see~\cite[Appendix]{Sa19} for instance.

Our main theorem is the following.

\begin{theo}
Let $M=(\T^2,F)$ be a Finsler two-torus whose geodesic flow is conjugate to the geodesic flow of a flat Finsler two-torus~$M_\diamond = (\T^2,F_\diamond)$. Then there exists a canonical deformation~$(F_t)_{t\,\geq\,0}$ of Finsler metrics on~$\T^2$ with $F_0=F$ such that
\begin{enumerate}
\item\label{theo1.1.1} the geodesic flow of~$F_t$ is conjugate to the geodesic flow of~$F_\diamond$;
\item\label{theo1.1.2} the metric~$F_t$ converges to~$F_\diamond$ for the uniform convergence topology, up to isometry, as $t$ goes to infinity.
\end{enumerate}
\end{theo}

In this theorem, we consider the uniform convergence of metric spaces, where a sequence~$(d_n)$ of metrics on a given set~$X$ converges to a metric~$d$ on~$X$ if $d_n \to d$ uniformly on~$X \times X$ as $n$ goes to infinity. We refer to the topology it induces on the space of metrics on~$X$ as the \emph{uniform convergence topology}.

The following result is a consequence of the main theorem.

\begin{coro}\label{coro:main}
The space of Finsler metrics on the two-torus, modulo isometries, whose geodesic flow is conjugate to that of a flat Finsler metric strongly deformation retracts to the space of flat Finsler metrics, modulo isometries (which is contractible).

In addition, the strong deformation retraction is induced by the deformation of the geodesic foliation by the curve shortening flow on the Euclidean plane.
\end{coro}

Our results are of the same flavor as those of~\cite{Sa19}, where it is proved that the space of Zoll Finsler metrics on the projective plane~$\R \Prob^{2}$ whose geodesic length is equal to $\pi$ strongly deformation retracts to the canonical round metric. In this case too, the deformation retraction is induced by the curve shortening flow on the canonical round projective plane.



%\forget
\begin{comment}
We will however face several new difficulties due to the fact that not all geodesics are closed. Among them, the space of geodesics on the torus is not a manifold. This forces us to work on the universal cover of the torus, leading to several regularity and compactness issues, including convergence problems, averaging difficulties, wildness of group actions, etc. 

The first step is to deform the geodesics of the universal cover~$\bar{M}$ of a Finsler two-torus~$M$ without conjugate points into straight lines, in a $\Z^2$-equivariant way, pairwise preserving their intersection properties. This is done by applying the Euclidean curve shortening flow to the geodesic curves of~$\bar{M}$; see Section~\ref{sec:convergenceCSF}. The continuity of the curve shortening flow at the limit is an issue which is studied in Sections~\ref{sec:convergenceCSF} and~\ref{sec:uniform}. The deformation of the geodesic foliation induced by the curve shortening flow gives rise to a family of diffeomorphisms on the unit tangent bundle of the torus and yields a deformation of the geodesic flow of the Finsler metric into that of a flat Finsler metric; see Sections~\ref{sec:diffeo} and~\ref{sec:deformation}. Now, the idea is to apply the construction of Finsler metrics with prescribed geodesics on the plane through Crofton's formula due to \'Alvarez Paiva and Berck~\cite{AB}. More precisely, there is a one-to-one correspondence between Finsler metrics without conjugate points and its space of geodesics endowed with a natural smooth measure; see Sections~\ref{sec:crofton} and~\ref{sec:prescribe}. In our case, we show that, at the limit, the curve shortening flow preserves the natural measure on the space of deformed geodesics when the geodesic flow of the Finsler metric is conjugate to that of a flat Finsler metric; see Section~\ref{sec:conjugate}.


For the proof of the main theorem, we introduced the curve shortening flow for planar graphs and the equivalent graphical curve shortening flow given by the equation $\frac{\partial\gamma}{\partial t}=\kappa\nu$, where $\gamma(.,t)$ in $\R^{2}$ is a family of curves in $\R^{2}$ and $\kappa$ stands for the curvature of $\gamma$ and $\nu$ is the unit normal vector to $\gamma$. The notion of asymptotic direction of a geodesic of a two-torus will be used a lot in this paper regarding the geodesic foliations where in a Finsler two-torus without conjugate points, the geodesics of rational and irrational direction foliate the two-torus; see Theorem~\ref{theo:foliation}. In the universal cover of a Finsler two-torus $\T^{2}$ without conjugate points, we proved some interesting properties of the curve shortening flow, such as it preserves any asymptotic direction of a minimizing geodesic and that the limit of the curve shortening flow of any graph in $\R^{2}$ exists and will give a geodesic of a Finsler two-torus. In addition, another invariance of the curve shortening flow is the natural measures on the space of geodesics of $M$. The construction of the deformation retraction relies on the construction of Finsler metrics without conjugate points through Crofton formula due to \'Alvarez Paiva and Berck~\cite{AB}, that is, there is a one-to-one correspondence between Finsler metrics without conjugate points and Crofton metrics (with respect to the area form constructed on the space of geodesics).
\end{comment}
%\forgotten


In the proof of the main theorem, we will also establish the following result regarding the deformation of the geodesic flows of a two-torus without conjugate points.

\begin{theo}\label{theo:intro2}
Let $M=(\T^2,F)$ be a Finsler two-torus without conjugate points. Denote by $\rho:\R \times U_0\R^2 \to U_0\R^2$ the action induced by the geodesic flow on the unit tangent bundle~$U\bar{M} \simeq U_0\R^2$ of~$\bar{M}$. Then, there exists a deformation
\[
\rho_t:\R \times U_0\R^2 \to U_0\R^2
\]
of smooth, free, proper, $\Z^2$-equivariant actions that starts at~$\rho_0=\rho$ and converges to the action $\rho_\infty:\R \times U_0\R^2 \to U_0\R^2$ induced by the geodesic flow on the unit tangent bundle~$U_0\R^2$ of the Euclidean plane. Here, the convergence is in the compact-open $C^k$-topology for any given~\mbox{$k \geq 0$}.

Furthermore, for every~$t \in [0,\infty]$, every $\rho_t$-orbit projects to an embedding of~$\R$ into~$\R^2$ under the canonical projection $U_0\R^2 \to \R^2$.
\end{theo}

By construction, the deformation $\rho_t:\R \times U_0\R^2 \to U_0\R^2$ in Theorem~\ref{theo:intro2} is induced by the deformation of the geodesics of~$\bar{M}$ under the Euclidean curve shortening flow. Since the Euclidean curve shortening flow preserves the asymptotic directions of these geodesics, we deduce the following result regarding the deformation of the Heber foliation.

\begin{coro}
Let $M=(\T^2,F)$ be a Finsler two-torus without conjugate points. Then, the Heber foliation of~$T^*\bar{M}$ is deformed into the canonical Heber foliation of~$T^*\R^2$ (induced by straight lines) under the deformation \mbox{$\rho_t:\R \times U_0\R^2 \to U_0\R^2$}, via the Legendre transform.
\end{coro}

The proof strategy of our main theorem relies on arguments from dynamical systems, geometric flow, and integral geometry, and follows the approach developed in~\cite{Sa19} to study the space of Zoll Finsler metrics on the projective plane. We work on the universal cover~$\bar{M}$ of Finsler two-tori~$M$ without conjugate points, where all geodesics are minimizing (but not closed), and perform a series of $\Z^2$-equivariant constructions. The space of (oriented) geodesics of~$\bar{M}$ is a manifold diffeomorphic to~$S^1 \times \R$ equipped with a natural smooth measure induced by the canonical symplectic form on the cotangent bundle~$T^*\bar{M}$ of~$\bar{M}$. Now, the idea is to make use of the one-to-one correspondence between Finsler metrics without conjugate points and their space of geodesics endowed with their natural smooth measure through Crofton's formula. This correspondence, due to \'Alvarez Paiva and Berck~\cite{AB}, allows us to construct Finsler metrics with prescribed geodesics on the plane. We then deform the geodesics of~$\bar{M}$ into straight lines while preserving their asymptotic direction and intersection properties through the curve-shortening flow for the Euclidean metric (with an analysis of the underlying parabolic PDE). This deformation of the geodesics of~$\bar{M}$ along with the deformation of the natural measure on their moduli space induces a canonical deformation~$(\bar{F}_t)$ of the initial Finsler metric~$\bar{F}$ on~$\bar{M}$. Using a fine analysis of the continuity of the curve-shortening flow at the limit for irrational directions, we show that the deformation process preserves the natural measure on the space of deformed geodesics when the geodesic flow of the Finsler metric~$F$ is conjugate to that of a flat Finsler metric~$F_\diamond$. This ensures that the Finsler metric deformation~$(\bar{F}_t)$ converges to~$\bar{F}_\diamond$. By construction, the Finsler metrics~$(\bar{F}_t)$ of the deformation have no conjugate points. The final step is to show that these metrics have conjugate geodesic flows modulo the action of~$\Z^2$. This is precisely where we need to assume the geodesic flow of our initial metric is conjugate to that of a flat Finsler metric. We first observe that the Finsler metrics~$(\bar{F}_t)$ of the deformation have the same length spectrum. Then we show that their Heber foliations are smooth. To conclude, we apply the following dynamical result proved in Section~\ref{sec:smoothfoliation}.

\begin{theo}
Let $M_1=(\T^2,F_1)$ and $M_2=(\T^2,F_2)$ be two Finsler tori without conjugate points whose Heber foliations are smooth. Assume that both metrics have the same marked length spectrum. Then the geodesic flows of~$M_1$ and~$M_2$ are conjugate.
\end{theo}


To implement the proof scheme described above, numerous difficulties need to be faced: lack of regularity, lack of compactness, convergence issues, lack of averaging procedure, non-continuous nature of group actions, etc. Consequently, the proof relies on a wide range of techniques in dynamical systems, integral geometry and partial differential equations.

\subsection*{Acknowledgments} 
The authors thank Roman Karasev for several comments, which helped improve the exposition.



\section{Asymptotic directions, Busemann functions and the Heber foliation} \label{sec:foliation}

In this section, we introduce the asymptotic direction of a minimizing geodesic, define Busemann functions and present some general results regarding the Heber foliation of Finsler two-tori without conjugate points.

\begin{defi}\label{def:finsler}
A (reversible) \emph{Finsler metric} on a manifold~$M$ is a continuous function $F:TM \to [0,\infty)$ on the tangent bundle~$TM$ of~$M$ satisfying the following properties (here, $F_x:=F_{|T_x M}$ for short):
\begin{enumerate}
\item\label{item:smooth} Smoothness: $F$ is smooth outside the zero section; 
\item\label{item:homogene} Homogeneity: $F_x(tv) = |t| \, F_x(v)$ for every $v \in T_x M$ and $t \in \R$; 
\item\label{item:convex} Quadratic convexity: for every $x \in M$, the square of the function~$F_x$ has positive definite second derivatives on $T_x M \setminus \{ 0\}$, that is, for every $p\in T_{x}M\setminus\{0\}$, and $u,v \in T_x M$, the symmetric bilinear form
\[
g_p(u,v) = \frac{1}{2} \, \frac{\partial^2}{\partial s \partial t} F_x^2(p+tu+sv)_{|t=s=0}
\]
is an inner product. 
\end{enumerate}


The $F$-length of a smooth curve~$\gamma:[a,b] \to M$ is defined by
\[
\length_F(\gamma) = \int_a^b F(\gamma'(t)) \, dt
\]
and the distance $d_F(p,q)$ between two points $p,q \in M$ is the infimum $F$-length of the curves joining~$p$ and~$q$.

Consider the \emph{Legendre transform}
\begin{equation}\label{eq:Legendre}
\LL:TM\rightarrow T^{*}M
\end{equation}
of the Lagrangian~$\frac{1}{2}F^{2}$ between the tangent and cotangent bundles of~$M$; see~\cite{Be78}. Since $F$ is quadratically convex, the Legendre transform~$\LL$ is a diffeomorphism between $TM \setminus\{0\}$ and $T^{*}M\setminus\{0\}$. By the homogeneity of~$F$, it preserves the norm on each fiber of the bundle vectors~$TM$ and~$T^{*}M$. In addition, it induces a diffeomorphism between the unit tangent and cotangent bundles~$UM$ and~$U^{*}M$ of~$M$. Geometrically, this diffeomorphism is defined as follows: For every vector $v\in U_{x}M$, the image~$\LL(v)$ is the unique co-vector of~$U_{x}^{*}M$ such that $\LL(v)(v)=1$.

The quadratically convex condition (as opposed to a mere convex condition) also allows us to define the geodesic and cogeodesic flows on~$UM$ and~$U^*M$; see~\cite{Be78}. Note that both flows are conjugate by the Legendre transform.


The \emph{tautological one-form} on $T^{*}M$ is defined by
\begin{equation}\label{eq:alpha}
\alpha_{\xi}(X)=\xi\left(dp_{\xi}(X)\right)
\end{equation}
for every $\xi\in T^{*}M$ and $X\in T_{\xi}T^{*}M$, where $p:T^*M \to M$ is the canonical projection. Similarly, the \emph{canonical symplectic form} on~$T^*M$ is defined as
\begin{equation}\label{eq:symplec}
\omega=d\alpha.
\end{equation}
Note that neither~$\alpha$ nor~$\omega$ depend on the Finsler metric on~$M$. Still, by the Liouville theorem, see~\cite[Remark~2.12]{Be78}, both the one-form~$\alpha$ and the symplectic form~$\omega$ are invariant under the cogeodesic flow of~$M$.
\end{defi}

%\forget
\begin{comment}
Conjugate points on Finsler manifolds can be defined as follows; see~\cite[\S~5.4]{BCS}.

\begin{defi}
Let $M$ be a complete Finsler manifold. Two points~$x,y \in M$ are \emph{conjugate} along a geodesic arc~$\gamma$ connecting them if there exists a nonzero Jacobi field~$Y$ along~$\gamma$ vanishing at~$x$ and~$y$. Recall that a \emph{Jacobi field}~$Y$ along a geodesic arc~$\gamma$ of~$M$ is a vector field along $\gamma$ such that
\[
Y''+R(\gamma',Y)\gamma'=0,
\]
where $R$ denotes the curvature tensor.
\end{defi}

We will make use of the following fundamental property regarding Finsler manifolds without conjugate points; see~REF???

\begin{theo}
Let $M$ be a Finsler manifold without conjugate points. Denote by~$\bar{M}$ its universal Finsler cover. Then any pair of points in~$\bar{M}$ are joined by a unique geodesic arc.
\end{theo}
\end{comment}
%\forgotten

Finsler metrics without conjugate points can be defined as follows.

\begin{defi}\label{def:ncp}
A complete Finsler manifold~$M$ has no conjugate points if the exponential map of its universal Finsler cover~$\bar{M}$ is a diffeomorphism at every point, that is, if $\exp_x:T_x \bar{M} \to \bar{M}$ is a diffeomorphism for every $x \in \bar{M}$. In this case, any pair of points in~$\bar{M}$ can be joined by a unique geodesic.
\end{defi}


We will also need the following definition about Finsler geodesics.

\begin{defi}
Let $M$ be a complete Finsler manifold with universal Finsler cover~$\bar{M}$. A geodesic of~$\bar{M}$ is \emph{minimizing} if it minimizes the length between any pair of its points. (Such geodesics are also referred to as \emph{$A$-geodesics}.)
\end{defi}


Before giving a characterization of Finsler tori without conjugate points in terms of integrable cogeodesic flow, we need to introduce the following definitions.

\begin{defi}
A \emph{Lipschitz graph} of~$T^*M$, where $M$ is a manifold, is a Lipschitz section of the canonical projection $T^*M \to M$. In particular, it is a Lipschitz submanifold of~$T^*M$. By Rademacher's theorem, a Lipschitz submanifold of~$T^*M$ admits a tangent subspace almost everywhere with respect to the Lebesgue measure on the submanifold.

A Lipschitz graph of~$T^*M$ is \emph{Lagrangian} if its (almost everywhere defined) tangent subspaces are Lagrangian with respect to the canonical symplectic form~$\omega$ on~$T^*M$; see~\eqref{eq:symplec}.

A subset of~$T^*M$ is \emph{flow-invariant} if it is invariant under the cogeodesic flow of~$M$, where $M$ is a complete Finsler manifold.
\end{defi}

The following result has been established in~\cite{MS11} in the context of Tonelli Hamiltonians and generalized in higher dimension in~\cite{AABZ15}. See also~\cite{Ba88,BP86,He32,Ma91} and~\cite{Sch15}.

\begin{theo}\label{theo:foliation}
Let $M=(\T^2,F)$ be a Finsler two-torus. Then the torus~$M$ has no conjugate points if and only if there exists a continuous foliation of~$T^* M$ by Lipschitz, Lagrangian, flow-invariant graphs.
\end{theo}

In the rest of this section, we will present a construction of the continuous foliation of~$T^* M$ by Lipschitz, Lagrangian, flow-invariant graphs given by Theorem~\ref{theo:foliation} when $M$ is a Finsler two-torus without conjugate points. This torus foliation is referred to as the \emph{Heber foliation} of~$T^*M$. \\

The notion of asymptotic direction of a Finsler geodesic will play a key role regarding geodesic foliations of Finsler two-tori without conjugate points.

\begin{defi}\label{def:asymptotic}
The \emph{asymptotic direction} of a curve~$\gamma:\R \to \R^2$ going to infinity (\ie, \mbox{$\lVert \gamma(t) \rVert \to +\infty$} as $t \to +\infty$) is defined as
\begin{equation}\label{eq:delta}
\theta(\gamma) = \lim_{t\,\to\,+\infty} \frac{\gamma(t)}{\lVert \gamma(t) \rVert} \in S^1
\end{equation}
where $\lVert \cdot \rVert$ represents the Euclidean norm of~$\R^2$ (if the limit exists). The asymptotic direction of a curve in~$\T^2$ is defined as the asymptotic direction of any of its lifts.

The asymptotic direction of a curve of~$\R^2$ or~$\T^2$ is \emph{rational} if it is proportional to a vector with rational coefficients and \emph{irrational} otherwise.
\end{defi}

The following result about the existence of asymptotic directions for minimizing geodesics in the universal cover of a Finsler two-torus is due to Hedlund~\cite{He32} in the Riemannian case; see~\cite{Ba88,BP86,Ma91,Sch15} for further generalizations encompassing the Finsler case.


\begin{theo}\label{theo:hedlund}
Let $M=(\T^2,F)$ be a Finsler two-torus with universal Finsler cover $\bar{M}=(\R^2,\bar{F})$. Then, there exists $w>0$ such that for every minimizing geodesic $\gamma:\R \to \bar{M}$, the asymptotic direction~$\theta(\gamma)$ of~$\gamma$ is well-defined (that is, the limit in~\eqref{eq:delta} exists) and the geodesic~$\gamma$ lies in a strip bounded by two parallel lines in~$\R^2$ of width at most~$w$. Furthermore, $\theta(\bar{\gamma})=-\theta(\gamma)$, where $\bar{\gamma}:\R \to \bar{M}$ is the minimizing geodesic obtained by reversing the orientation, that is, $\bar{\gamma}(t)=\gamma(-t)$.

Conversely, for every~$\theta_0 \in S^1$, there exists a minimizing geodesic~\mbox{$\gamma:\R \to \bar{M}$} such that \mbox{$\theta(\gamma) = \theta_0$} and \mbox{$\theta(\bar{\gamma}) = -\theta_0$}. Moreover, for every $x \in \bar{M}$ and every~$\theta_0 \in S^1$, there exists a minimizing geodesic ray based at~$x$ with asymptotic direction~$\theta_0$.
\end{theo}



We will also need the following definition of Busemann functions; see~\cite{CK95,CS86,esc77,Heb94} and~\cite{gro99} for further details.

\begin{definition}\label{def:Bus}
Let $M=(\T^2,F)$ be a Finsler two-torus without conjugate points. Denote by~$\bar{M}=(\R^2,\bar{F})$ the universal Finsler cover of~$M$. The \emph{Busemann function} based at the origin~\mbox{$o \in \R^2$} and pointing in the direction of~$\theta$ is defined as
\[
B_\theta(x) = \lim_{t\,\to\,+\infty} d_{\bar{M}}(x,c_\theta(t)) -t
\]
where $c_\theta$ is the arc length parametrized geodesic ray arising from~$o$ with asymptotic direction~$\theta$. This limit is well-defined since the function $t \mapsto d_{\bar{M}}(x,c_\theta(t)) -t$ is monotone nonincreasing and bounded from below by the triangle inequality.

The collection of Busemann functions~$B_\theta$ with $\theta \in S^1$ gives rise to the \emph{Busemann map}
\begin{align*}
\begin{split}
B:S^1 \times \bar{M} & \rightarrow \R \\
(\theta,x) & \mapsto B_\theta(x).
\end{split}
\end{align*}
\end{definition}

Let us recall some regularity properties of Busemann functions for Finsler tori without conjugate points; see~\cite{CK95,CS86,esc77,Heb94}.

\begin{prop}\label{prop:heber}
Let $M=(\T^2,F)$ be a Finsler two-torus without conjugate points. Denote by~$\bar{M}=(\R^2,\bar{F})$ its universal Finsler cover. Then
\begin{enumerate}
\item\label{heber1} The Busemann map $B:S^1 \times \bar{M} \rightarrow \R$ is continuous.
\item\label{heber2} Every Busemann function $B_\theta : \bar{M} \to \R$ is $C^{1,1}$ (\ie, $B_{\theta}$ is $C^{1}$ and its differential $dB_{\theta}$ is locally Lipschitz on $M$). 
\item\label{heber3} The differential $dB_\theta:T\bar{M} \to \R$ is $\Z^2$-invariant and has unit norm at every point of~$\bar{M}$ (\ie, $\lVert d_{x}B_{\theta} \rVert=1$ for every $x\in\bar{M}$). 
\item\label{heber4} Via the Legendre transform~$\LL:T\bar{M} \to T^*\bar{M}$, see~\eqref{eq:Legendre}, the differential~$-dB_\theta$ at~$x \in \bar{M}$ corresponds to the unit tangent vector at~$x$ generating a geodesic ray with asymptotic direction~$\theta$.
\end{enumerate}
\end{prop}

We can now describe the Heber foliation of a Finsler two-torus without conjugate points.

\begin{defi}\label{def:heber}
Let $M=(\T^2,F)$ be a Finsler two-torus without conjugate points. The unit tangent bundle~$U\bar{M}$ of the universal cover $\bar{M}$ of $M$ identifies with~$\R^2 \times S^1$ as follows
\begin{align}\label{eq:UM}
\begin{split}
U\bar{M} & \rightarrow \R^2 \times S^1 \\
(x,v) & \mapsto (x,\theta)
\end{split}
\end{align}
where $\theta$ is the asymptotic direction of the geodesic~$\gamma_v$ of~$\bar{M}$ induced by~$v$. By Proposition~\ref{prop:heber}.\eqref{heber4}, this map is a homeomorphism whose inverse map given by
\begin{equation}\label{eq:v}
v=-\LL^{-1}(dB_\theta(x)).
\end{equation}
By definition, the \emph{Heber homeomorphism} is the inverse map of~\eqref{eq:UM}.

By $\Z^2$-invariance of the differential of the Busemann functions, see Proposition~\ref{prop:heber} \eqref{heber3}, the map
\begin{align*}
\begin{split}
\bar{M} & \rightarrow T^*\bar{M} \\
x & \mapsto \varrho \, d B_\theta(x)
\end{split}
\end{align*}
passes to the quotient under the $\Z^2$-action on~$\bar{M}$ and~$T^*\bar{M}$, and induces a map
\begin{equation}\label{eq:Hleaves}
M \to T^*M
\end{equation}
for $\theta \in S^1$ and $\varrho \in \R_+$.

The \emph{Heber foliation} associated to~$M$ is a foliation of~$T^*M$ whose leaves are the images of the maps~\eqref{eq:Hleaves}. Note that the leaves~\eqref{eq:Hleaves} of the Heber foliation are Lipschitz graphs of~$T^*M$; see Proposition~\ref{prop:heber}$\MK$\eqref{heber2}. The Heber foliation is said to be \emph{smooth} if the Heber homeomorphism $\R^2 \times S^1 \to U \bar{M}$ is a smooth diffeomorphism.

%\forget
\begin{comment}
The \emph{Heber foliation} of~$M$ is a foliation of~$T^*M$ which can be described through its leaves given by
\begin{align}\label{eq:Hleaves}
\begin{split}
M & \rightarrow T^*M \\
x & \mapsto \varrho \, d B_\theta(x)
\end{split}
\end{align}
where $\theta \in S^1$ and $\varrho \in \R_+$. Here, despite the risk of confusion, we still denote by~$d B_\theta(x)$ its projection to~$T^*M$ under the $\Z^2$-action on~$T^* \bar{M}$, where the $\Z^2$-invariance of the differential is given by the item~\eqref{heber2} of Proposition~\ref{prop:heber}. Note that the leaves~\eqref{eq:Hleaves} are only Lipschitz; see Proposition~\ref{prop:heber}$\MK$\eqref{heber2}.
%\forgotten
\end{comment}


\end{defi}

%\forget
\begin{comment}
A graph is defined as a section of $\pi:T^{*}M\to M$. We shall need the following definition about Lipschitz lagrangian graphs.

\begin{defi}\label{Lip:Lag}
The graph of a continuous $1$-form $\eta:M\to\R$ is said to be Lipschitz Lagrangian if and only if for every Lipschitz loop $\gamma$ drawn on $M$ which is homotopic to a point, we have $\int_{\gamma}\eta=0$.
\end{defi}
\end{comment}
%\forgotten


The main properties of the Heber foliation are given by the following proposition; see~\cite{AABZ15,Ba88,BP86,He32,Ma91,MS11} and~\cite{Sch15}.

\begin{prop}\label{prop:Hfoliation}
Let $M=(\T^2,F)$ be a Finsler two-torus without conjugate points. Then its Heber foliation is a continuous foliation of~$T^* M$ by Lipschitz, Lagrangian, flow-invariant graphs. Each of these graphs, except for the zero section, can be described as the space of covectors of the same norm generating a geodesic of~$M$ with the same asymptotic direction in~$S^1$.

Moreover, the geodesics of~$M$ do not have self-intersection and can be classified according to their asymptotic direction as follows:
\begin{enumerate}
\item\label{geod1} the geodesics of~$M$ with a given rational asymptotic direction are closed of the same length and smoothly foliate~$M$; 
\item\label{geod2} the geodesics of~$M$ with a given irrational asymptotic direction are non-closed and foliate~$M$. 
\end{enumerate}
\end{prop}


\begin{rema}
It is an open question whether the Heber foliation associated to a Finsler two-torus~$M$ without conjugate points is always smooth. Observe however that the closed geodesics of~$M$ homotopic to a simple closed curve define a smooth geodesic foliation of~$M$ (and have the same length); see the case~\eqref{geod1} of Proposition~\ref{prop:Hfoliation} (or~\cite[Section~1.2]{AABZ15}). The situation is unclear in the case~\eqref{geod2}: the foliation of~$M$ by geodesics with a given irrational asymptotic direction may not be smooth.
\end{rema}



%\forget
\begin{comment}
The following description of geodesics of the two-torus can be found in~\cite{Ba88,BP86,He32,Ma91,Sch15}.

\begin{prop}
Let $M=(\T^2,F)$ be a Finsler two-torus. Then, the closed minimal geodesics of~$M$ with a given rational asymptotic direction foliate a closed subset of~$M$, while nonclosed minimal geodesics of~$M$ with the same rational asymptotic direction lie in the complement of this closed subset.
\end{prop}

\begin{rema}
For Finsler metrics on~$\T^2$ without conjugate points, closed minimal geodesics of~$\T^2$ with a given rational asymptotic direction foliate the whole torus~$\T^2$. In particular, every minimal geodesic of~$\T^2$ with a given rational asymptotic direction is closed.
\end{rema}

\begin{prop}\label{prop:foliation-baby}
Let $\alpha$ be a simple closed geodesic on a Finsler two-torus~$M$ without conjugate points. Then the closed geodesics homotopic to~$\alpha$ have the same length and define a smooth geodesic foliation of~$M$.
\end{prop}
\end{comment}
%\forgotten


\section{Standard identification of Finsler two-tori without conjugate points}

In this technical section, we present a ``standard'' identification of Finsler two-tori without conjugate points, where the horizontal and vertical closed curves are geodesics.

\begin{prop}\label{prop:identification}
Every Finsler two-torus~$M=(\T^2,F)$ without conjugate points can be identified to~$S^1 \times S^1$ so that the horizontal and vertical curves~$S^1 \times \{t\}$ and~\mbox{$\{s\} \times S^1$} are geodesics with $s,t \in S^1$.

In this case, any nonvertical geodesic in the universal cover~$\R^2$ of~$\T^2$ is the graph of a monotonic smooth function $\bu:\R \to \R$ over its horizontal axis.
\end{prop}

\begin{proof}
Two closed geodesics of an orientable closed Finsler surface of minimal length in their homotopy classes have a minimal number of intersection points among homotopic loops; see~\cite{FHS82}.

Let $\alpha$ and~$\beta$ be two simple closed geodesics of~$M$ parameterized proportionally by arclength intersecting once. We know that the closed geodesics homotopic to~$\alpha$ have the same length and form a smooth geodesic foliation~$\alpha_t$ of~$M$; see Proposition~\ref{prop:Hfoliation}. Observe that the geodesics~$\alpha_t$ transversely intersect the geodesic~$\beta$ at a single point. Thus, we can choose the parameter~$t$ in the geodesic foliation~$(\alpha_t)$ so that $\alpha_t$ is the unique closed geodesic homotopic to~$\alpha$ with~$\alpha_t(0) = \beta(t)$. Similarly, there exists a smooth geodesic foliation~$\beta_s$ of~$M$ so that $\beta_s$ is the unique closed geodesic homotopic to~$\beta$ (of the same length) with~\mbox{$\beta_s(0)=\alpha(s)$}. Note that both $\alpha_t$ and~$\beta_s$ have minimal length in their homotopy classes and smoothly depend on~$t$ and~$s$. It follows from the observation at the beginning of the proof that the closed geodesics~$\alpha_t$ and~$\beta_s$ intersect once.


Since $(\alpha_t)$ and~$(\beta_s)$ are two transverse foliation of~$\T^2$ whose leaves intersect once, the map
\[
\begin{array}{rccc}
\phi : & S^{1}\times S^{1} & \to & \T^{2} \\
& (s,t) & \mapsto & \alpha_{t}\cap\beta_{s} \\
\end{array}
\]
is a bijection. Let us show that $\phi$ is a diffeomorphism. Since both $\alpha_t$ and~$\beta_s$ smoothly depend on~$t$ and~$s$, the map~$\phi$ is smooth. The curves~$\phi(\cdot,t)=\alpha_t$ define a geodesic variation. Since the Finsler metric on~$M$ has no conjugate points, the Jacobi vector field~$\frac{\partial \phi}{\partial t}$ it generates along the closed geodesic~$\alpha_t$ does not vanish; see~\cite[\S~5.4]{BCS}. Thus, $\frac{\partial \phi}{\partial t}$ is a nonvanishing vector field parallel to~$\beta'_s$. Similarly, $\frac{\partial \phi}{\partial s}$ is a nonvanishing vector field parallel to~$\alpha'_t$.

Since $\alpha'_t$ and~$\beta'_s$ are noncolinear, the same goes for~$\frac{\partial \phi}{\partial s}$ and~$\frac{\partial \phi}{\partial t}$. It follows that the bijective map $\phi:S^1 \times S^1 \to \T^2$ is a local diffeomorphism. With this identification, the horizontal and vertical curves~$S^1 \times \{t\}$ and~$\{s\} \times S^1$ coincide with the closed geodesics~$\alpha_t$ and~$\beta_s$.

By the description of geodesics on a two-torus given by Theorem~\ref{theo:hedlund}, every nonvertical geodesic~$\gamma$ on the universal cover of a Finsler two-torus~$\T^2$ without conjugate points has a nonvertical direction. Thus, the geodesic~$\gamma$ transversely intersects the vertical geodesic lines of~$\R^2$ at least once and so exactly once. Therefore, the geodesic~$\gamma$ is the graph of a smooth function $\bu:\R \to \R$ over the horizontal axis of~$\R^2$. Similarly, the geodesic~$\gamma$ intersects every horizontal line exactly once unless it has a horizontal direction in which case it coincides with a horizontal line. It follows that the function $\bu:\R \to \R$ is monotonic.
\end{proof}



\section{Convergence of the curve shortening flow} \label{sec:convergenceCSF}

We introduce the curve shortening flow in the Euclidean plane and show the existence of a limit when applied to the graph of a function representing a (minimizing) geodesic on the universal cover of a Finsler two-torus without conjugate points.


\begin{defi}\label{def:CSF}
Consider the curve shortening flow for (planar) graphs given by the following quasilinear parabolic equation
\begin{equation}\label{eq:GCSF}
\bu_t=\frac{\bu_{xx}}{1+\bu_{x}^{2}}
\end{equation}
with a $C^\infty$ initial condition~$\bu(.,0)=\bu_0$. Here, the function $\bu:\R \times [0,+\infty) \to \R$ takes~$(x,t)$ to~$\bu(x,t)$. By~\cite{EH91} (see also~\cite{CZ98}), for every $C^\infty$ initial condition~$\bu_0$, the curve shortening flow~\eqref{eq:GCSF} has a unique solution defined for every $t\in[0,+\infty)$. Note that this flow extends to curves of the square flat two-torus~$\T^2=\R^2/\Z^2$ whose lifts~$\gamma$ to the universal cover~$\R^2$ are smooth graphs.

The curve shortening flow for curves in $\R^{2}$ is defined using the following equation
\begin{equation}\label{eq:CSF}
\frac{\partial \gamma}{\partial t} = \kappa \, \nu
\end{equation}
where $\gamma(.,t)$ is a family of curves in~$\R^2$ with curvature~$\kappa$ and unit normal vector~$\nu$.

It is known that for an initial curve~$\gamma(.,0)$ which is the graph of a function~$\bu_0$, the evolution equations~\eqref{eq:GCSF} and~\eqref{eq:CSF} are equivalent: the evolution curve~$\gamma(.,t)$ is given by the graph of~$\bu(.,t)$; see~\cite{EH91} or~\cite{CZ98} for instance.
\end{defi}


For a smooth, complete, properly embedded curve~$\gamma_0=\gamma(.,0)$ dividing the plane into two regions of infinite area (which is the case when the initial condition~$\gamma_0$ is the graph of a smooth function~$\bu_0$), the curve shortening flow~\eqref{eq:CSF} has a unique solution; see~\cite{CZ98}. As previously, this flow extends to curves of the flat two-torus~$\T^2=\R^2/\Z^2$ whose lifts to the universal cover~$\R^2$ satisfy the previous existence and uniqueness condition. 


The curve shortening flow satisfies the following crucial property; see~\cite{CZ98}.

\begin{theo}\label{theo:CSF}
Two disjoint smooth, complete, properly embedded curves~$\gamma_1$ and~$\gamma_2$ dividing the plane into two regions of infinite area remain disjoint and embedded through the curve shortening flow.
\end{theo}



The following result applies to all geodesics in the universal cover of a Finsler two-torus without conjugate points.

\begin{prop}\label{prop:direction}
The asymptotic direction of a minimizing geodesic in the universal cover~$\bar{M}=(\R^2,\bar{F})$ of a Finsler two-torus~$M=(\T^2,F)$ is preserved under the curve shortening flow.
\end{prop}

\begin{proof}
Every minimizing geodesic~$\gamma$ in~$\bar{M} \simeq \R^2$ lies in an open strip~$S$ bounded by two straight lines of~$\R^2$ with the same asymptotic direction as~$\gamma$; see Theorem~\ref{theo:hedlund}. Since these two lines are fixed under the curve shortening flow of~$\R^2$ and disjoint curves remain disjoint, see Theorem~\ref{theo:CSF}, the curve~$\gamma$ evolves within the strip~$S$ under the curve shortening flow, keeping the same asymptotic direction.
\end{proof}

In general, the curve shortening flow of a smooth function~$\bu_0:\R \to \R$ does not necessarily converge; see~\cite{EH89}. Conditions under which such a convergence occurs can be found in~\cite{NT07} and~\cite{WW}. In our case, we show the following result.

\begin{theo}\label{theo:affine}
Let $M=(\T^2,F)$ be a Finsler two-torus without conjugate points. Denote by~$\bar{M}=(\R^2,\bar{F})$ its universal Finsler cover. Let $\bu_0:\R \to \R$ be a smooth function whose graph~$\GG_0$ is a (minimizing) Finsler geodesic of~$\bar{M}$. Then, the solution~$\bu(.,t)$ of the curve shortening flow~\eqref{eq:GCSF} with initial condition~$\bu_0$ converges to an affine function~$\bu_\infty$ (with the same asymptotic direction as~$\bu_0$). Here, the convergence is in the $C^k$-topology for any given~$k \geq 0$.
\end{theo}

We will consider two cases in the proof of Theorem~\ref{theo:affine}: first, the rational case where the asymptotic direction of the minimizing geodesic represented by the graph of~$\bu_0$ is rational (proved at the end of this section, even for metrics with conjugate points), then the irrational case (proved in the next section). In the latter case, we will actually need a stronger result, namely Theorem~\ref{theo:unifwidth}, giving a uniform convergence of the curve shortening flow for geodesics whose asymptotic direction lies in a small enough neighborhood of a fixed irrational direction.

\begin{rema}
We emphasize that the geometric nature of the curves and functions considered in Theorem~\ref{theo:affine} is crucial. Indeed, the conclusion fails for general functions as exemplified by the grim reaper curve $y = - \log \cos x$, which is a translating soliton of the curve shortening flow; see~\cite{H12} for a classification of self-similar solutions to the curve shortening flow in the plane.
\end{rema}


We will first prove Theorem~\ref{theo:affine} in the rational case, without assuming that $M$ has no conjugate points (as long as the geodesic represented by~$\clG_0$ is minimizing). 

Let us start with a first lemma about the growth of~$\bu_0$, which will also be used in Section~\ref{sec:diffeo}.

\begin{lemm}\label{lem:u0growth}
The graph~$\GG_0$ of~$\bu_0$ has linear growth. More precisely, there exists $C_M>0$ (which does not depend on~$\bu_0$) such that $|\bu'_0| \leq C_M$.
\end{lemm}

\begin{proof}
Let $a=\min_{[-1,1]} \bu_0$ and $b=\max_{[-1,1]} \bu_0$. Choose~$L>0$ such that every vertical open interval of Euclidean length~$L$ in~$\R^2$ has a translate by an element of the lattice~$\Z^2$ which intersects the rectangle~$[-1,1] \times [a-1,b+1]$ along a segment $\{x\} \times [a-1,b+1]$ with $x\in[-\frac{1}{2},\frac{1}{2}]$, and in particular, transversally intersects the part of the graph~$\GG_0$ over~$[-\frac{1}{2},\frac{1}{2}]$. For example, we can take $L=b-a+4$.

By cocompactness of the Finsler metric on~$\R^2$ (due to its $\Z^2$-periodicity), for every~$\varepsilon >0$, there exists $\delta>0$ such that every pair of $\delta$-close tangent vectors to~$\R^2$ generate $\varepsilon$-close length~$L$ geodesic rays.

By contradiction, suppose that $|\bu'_0|$ becomes arbitrarily large, that is, some tangent vector to the graph~$\GG_0$ becomes $\delta$-close to a vertical vector. Since the vertical lines of~$\R^2$ are Finsler geodesics, this implies that the geodesic $\GG_0$ becomes $\varepsilon$-close to some vertical segment of Euclidean length~$L$. By our choice of~$L$ and for $\varepsilon$ small enough, this implies that a translate of~$\GG_0$ by a (nontrivial) element of~$\Z^2$ intersects the rectangle $[-1,1]\times[a,b]$ along an arc joining $[-1,1]\times\{a\}$ and $[-1,1]\times\{b\}$ and in particular, transversally intersects~$\GG_0$. Thus, the geodesic of~$M$ corresponding to the projection of~$\GG_0$ to the torus has a transverse selfintersection. Hence a contradiction with Proposition~\ref{prop:Hfoliation}.
\end{proof}


The following result showing that the curve shortening flow flattens out the graph of~$\bu_0$ is a direct consequence of the upper bound on the curvature obtained in~\cite[Proposition~4.4]{EH89} combined with Lemma~\ref{lem:u0growth}.


\begin{lemm}\label{lem:curv}
The curvature of the graph~$\GG_t$ of~$\bu(.,t)$ converges to zero in the $C^k$-topology as $t$~goes to infinity (uniformly with respect to~$\bu_0$). More precisely, there exists $c_k=c_k(M) >0$ (which does not depend on~$\bu_0$) such that
\[
\lVert \kappa_t \rVert_{C^k}^2 \leq \frac{c_k}{t^{k+1}}
\]
for every $t>0$, where $\kappa_t$ is the curvature of~$\GG_t$.
\end{lemm}

\begin{proof}
A uniform upper bound on the second fundamental form of entire (hypersurface) graphs in~$\R^n$ moving by the mean curvature flow can be found in~\cite[Proposition~4.4]{EH89}. In the case of the graph~$\GG_t$, this bound can be written
\[
\lVert \kappa_t \rVert_{C^k}^2 \leq \frac{c_k}{t^{k+1}}
\]
where $c_k=c_k(M)$ is a constant depending only on~$k$ and the constant~$C_M$ given by~Lemma~\ref{lem:u0growth}.
\end{proof}



%\forget
\begin{comment}
\begin{proof}
By~\cite[Proposition~4.4]{EH89}, the result follows if we show that the graph~$\GG_0$ of~$\bu_0$ has linear growth, that is, $|\bu'_0|$ is bounded.

Let $a=\min_{[-1,1]} \bu_0$ and $b=\max_{[-1,1]} \bu_0$. Choose~$L>0$ such that every vertical open interval of Euclidean length~$L$ in~$\R^2$ has a translate by an element of the lattice~$\Z^2$ which intersects the rectangle~$[-1,1] \times [a-1,b+1]$ along a segment $\{x\} \times [a-1,b+1]$ with $x\in[-\frac{1}{2},\frac{1}{2}]$, and in particular, transversally intersects the part of the graph~$\GG_0$ over~$[-\frac{1}{2},\frac{1}{2}]$. For example, we can take $L=b-a+4$.

By contradiction, suppose that $|\bu'_0|$ becomes arbitrarily large, that is, $\GG_0$ becomes almost vertical. Since $\GG_0$ and the vertical lines of~$\R^2$ are Finsler geodesics, a subarc of~$\GG_0$ becomes arbitrarily close to some vertical segment of Euclidean length~$L$. By our choice of~$L$, this implies that a translate of~$\GG_0$ by a (nontrivial) element of~$\Z^2$ intersects the rectangle $[-1,1]\times[a,b]$ along an arc joining $[-1,1]\times\{a\}$ and $[-1,1]\times\{b\}$ and in particular, transversally intersects~$\GG_0$. Thus, the minimal geodesic of~$M$ corresponding to the projection of~$\GG_0$ to the torus has a transverse selfintersection. Hence a contradiction with Theorem~\ref{theo:foliation}.
\end{proof}
\end{comment}
%\forgotten

Our second result is about the width of the graph~$\GG_t$ defined as follows.

\begin{defi}
By Theorem~\ref{theo:hedlund}, the graph~$\GG_0$ lies in a strip of~$\R^2$ bounded by two parallel straight lines and so does the graph~$\GG_t$. The minimal Euclidean distance between two such lines is called the \emph{width} of~$\GG_t$.
\end{defi}

By the properties of the curve shortening flow, see Theorem~\ref{theo:CSF}, the width of~$\GG_t$ is nonincreasing. Moreover, we have

\begin{lemm}\label{lem:strip}
Suppose that the asymptotic direction of the geodesic represented by~$\GG_0$ is rational. Then, the width of~$\GG_t$ tends to zero as $t$ goes to infinity.
\end{lemm}

\begin{proof}
By contradiction, assume that there exists $w_0>0$ such that for every $t \geq 0$, the graph~$\GG_t$ is of width greater than~$w_0$.

Since the asymptotic direction~$\theta_0$ of~$\GG_0$ is rational, then the projection~$\gamma_0$ of~$\GG_0$ to the torus~$\T^2$ closes up. The cover of~$\T^2$ corresponding to the subgroup generated by the homotopy class of~$\gamma_0$ is a flat cylinder~$C=~S^1\times\R$ containing~$\gamma_0$. The projection~$\gamma_t$ of~$\GG_t$ to~$C$ lies in the minimal annulus $S^1\times I$ of~$C$ containing~$\gamma_0$. Since the curvature of~$\gamma_t$ uniformly converges to zero as $t$ goes to infinity, see Lemma~\ref{lem:curv}, the closed curve~$\gamma_t$ of bounded length becomes arbitrarily $C^1$-close to a circle of~$C$, for $t$ large enough. Thus, the width of~$\gamma_t$, and so of~$\GG_t$, tends to zero as $t$ goes to infinity.


%\forget
\begin{comment}
Suppose that the asymptotic direction~$\theta_0$ of~$\GG_0$ is irrational. There exists~$L=L(\theta_{0})>0$ such that every segment of Euclidean length~$L$ and direction~$\theta_0$ in the flat two-torus~$\T^2=\R^2/\Z^2$ decomposes the square $(0,1) \times (0,1)$, obtained by removing the horizontal and vertical circles $S^1 \times \{0\}$ and~$\{0\} \times S^1$ from~$\T^2$, into domains of width less than~$w_0$. Since the curvature of~$\GG_t$ uniformly converges to zero as $t$ goes to infinity, see Lemma~\ref{lem:curv}, some subarc of~$\GG_t$ becomes arbitrarily $C^1$-close to a segment of Euclidean length~$L$ and direction~$\theta_0$, for $t$ large enough. By our choice of $L$, the projection of this segment to $\T^{2}$ decomposes the square $(0,1)\times(0,1)$ into domains of width less than $w_{0}$ and transversally intersects the projection of $\GG_{t}$ whose width is greater than $w_{0}$ by assumption. Since some subarc of $\GG_{t}$ is arbitrarily close to this segment, it follows that the projection of $\GG_{t}$ to the torus has a transverse self-intersection. In order to derive a contradiction, observe that the geodesic $\GG_{0}$ has no self-intersection and that its translates under $\Z^{2}$ do not intersect each other, otherwise the projection of $\GG_{0}$ to $\T^{2}$ would have a self-intersection. By the properties of the curve shortening flow, the same holds for the graph $\GG_{t}$. In particular, the projection of $\GG_{t}$ to $\T^{2}$ has no self-intersection. Hence a contradiction.
\end{comment}
%\forgotten

\end{proof}

Combining the zero convergence results of the curvature of~$\GG_t$ and of its width when the asymptotic direction of~$\GG_0$ is rational, see Lemma~\ref{lem:curv} and Lemma~\ref{lem:strip}, we immediately deduce Theorem~\ref{theo:affine} in the rational case. To prove the theorem in the remaining irrational case, it is enough to show that the conclusion of Lemma~\ref{lem:strip} still holds when the asymptotic direction of~$\GG_0$ is irrational. This is done in the next section where a stronger result is proved; see Theorem~\ref{theo:unifwidth}.



%\forget
\begin{comment}
\section{The space of oriented geodesics}

Suppose that the Finsler metric~$F$ on the two-torus~$M$ has no conjugate points. Then the geodesic flow action of the lift~$\bar{F}$ of~$F$ on the universal cover~$\bar{M} \simeq \R^2$ of~$M$ is proper\dots.



Recall that the quotient manifold theorem, see~\cite{lee}[Theorem~7.10], asserts that if $G$ is a Lie group acting smoothly, freely and properly on a smooth manifold~$N$, then the quotient space $N/G$ is a topological manifold with a unique smooth structure such that the quotient map $N \to N/G$ is a smooth submersion. This result applies to the action~$\rho_{\bar{F}}$ of the geodesic flow of~$\bar{F}$ on~$U \bar{M}$.

\begin{defi}
Denote by
\[
\Gamma_{\bar{F}}=U \bar{M}/\rho_{\bar{F}}
\]
the quotient manifold and by
\[
q_{\bar{F}}:U \bar{M} \to \Gamma_{\bar{F}}
\]
the quotient submersion. The quotient manifold~$\Gamma_{\bar{F}}$ represents the space of unparametrized oriented geodesics of the manifold~$\bar{M}$ with the Finsler metric~$\bar{F}$.

Denote also by~$\Gamma_0$ the space of unparametrized oriented lines in the Euclidean plane~$\R^2$. When $\bar{F}$ is the Euclidean metric on~$\bar{M} =
\R^2$, we have $\Gamma_{\bar{F}} = \Gamma_0$.
\end{defi}

Consider the map
\[
\varphi: \Gamma_{\bar{F}} \to \Gamma_0
\]
taking a geodesic of~$\bar{M} \simeq \R^2$ to its limit in~$\R^2$ under the curve shortening flow~\eqref{eq:CSF}, see Theorem~\ref{theo:affine}.

\begin{prop}
The map $\varphi: \Gamma_{\bar{F}} \to \Gamma_0$ is an homeomorphism.
\end{prop}

\begin{proof}

Our initial curve $\gamma$ is wedged between two straight lines. First, let us prove that the supremum of $\gamma$ which is the height with respect to the orthogonal is attained. By contradiction, suppose that our curve is asymptotic to the lines, then if we translate $\gamma$ by a rational direction, we shall also translate the two straight lines with a slight downward shift, hence this translated curve will be asymptotic to a curve slightly below with a self-intersection which is impossible since there cannot exist any intersections with translations by $\Z^{2}$. \\

Let us show by contradiction that $\varphi:\Gamma_{\bar{F}} \to \Gamma_0$ is injective. By the standard identification, see Proposition~\ref{prop:identification}, the map~$\varphi$ takes every vertical line to itself and sends a non-vertical geodesic to a non-vertical line. By contradiction, take two different non-vertical geodesics~$\gamma_1$ and~$\gamma_2$ in~$\Gamma_{\bar{F}}$ with $\varphi(\gamma_1)=\varphi(\gamma_2)$. The geodesics~$\gamma_i$ can be represented as the graph of a smooth function $u_i:\R \to \R$.\\


Assume that $\gamma_{1}$ and $\gamma_{2}$ are contained in two strips with different directions, since the asymptotic directions is preserved, then we know that they will converge to two lines with also different asymptotic directions. In addition, if $\gamma_{1}$ and $\gamma_{2}$ intersect, then after applying the geodesic flow, they will also intersect and hence we get the injectivity.\\

Assume now that $\gamma_{1}$ and $\gamma_{2}$ have the same asymptotic direction and set $\delta=\inf_{x\in\R}(\gamma_{2}(x)-\gamma_{1}(x))$. If $\delta>0$, then
\end{proof}
%\forgotten
\end{comment}



\section{Convergence of the width at irrational directions} \label{sec:uniform}


We show that the width of a geodesic in the universal cover of a Finsler two-torus without conjugate points converges to zero under the curve shortening flow for geodesics whose asymptotic direction lies in a small enough neighborhood of a fixed irrational direction. We deduce that the straight lines obtained at the limit vary continuously at irrational directions.


\begin{theo}\label{theo:unifwidth}
Let $M=(\T^2,F)$ be a Finsler two-torus without conjugate points. Denote by~$\bar{M}=(\R^2,\bar{F})$ its universal Finsler cover. Let $\theta$ be an irrational direction. For every $\varepsilon>0$, there exists $t_0>0$ such that for every geodesic~$\gamma$ of~$\bar{M}$ with asymptotic direction close enough to~$\theta$ and every~$t \geq t_0$, the width of~$\gamma_t$ satisfies $\width(\gamma_t) < \varepsilon$.
\end{theo}

\begin{rema}
The point of Theorem~\ref{theo:unifwidth} is that $t_0$ does not depend on~$\gamma$ as long as its asymptotic direction remains close enough to the irrational direction~$\theta$. Of course, Theorem~\ref{theo:unifwidth} also applies when the asymptotic direction of the geodesic~$\gamma$ is irrational, which gives an analogue of Lemma~\ref{lem:strip} in the irrational case and allows us to derive Theorem~\ref{theo:affine} in this case too. We will need this stronger version of Theorem~\ref{theo:affine} to establish the continuity of the Euclidean curve shortening flow at the limit for Finsler geodesics with irrational asymptotic directions.
\end{rema}

We decompose the proof of the proposition into several lemmas.

\begin{defi}
The \emph{vertical distance} between two subsets $A, B \subset \R^{2}$ is defined as
\[
\sup_{V} \inf\{d(x,y) \mid x\in A\cap V, y\in B\cap V\}
\]
where $V$ runs over all vertical lines of~$\R^2$ intersecting~$A$ and~$B$, and $d$ represents the Euclidean distance in $\R^{2}$.
\end{defi}

The following lemma is a quantitative version of the fact that an irrational geodesic in the square flat torus is dense.

\begin{lemm}\label{lem:dioph}
Let $\theta$ be an irrational direction. For every $\varepsilon \in (0,1)$, there exists $L=L(\theta,\varepsilon)>0$ such that every segment~$c$ of~$\R^2$ of direction close enough to~$\theta$ which projects onto an interval of the $x$-axis of length at least~$L$ passes below a point~$x_+$ of~$\Z^2$ and above a point~$x_-$ of~$\Z^2$, both at vertical distance less than~$\varepsilon$ from the segment~$c$.
\end{lemm}


\begin{proof}
Let $\varepsilon \in (0,1)$. Fix a point~$\bar{x} \in \Z^2$. Let $\bar{x}_+$ and~$\bar{x}_-$ be the points of~$\R^2$ at vertical distance~$\varepsilon$ from~$\bar{x}$, with $\bar{x}_+$ above~$\bar{x}$ and $\bar{x}_-$ below~$x$. Let~$\bar{m}_+$ and~$\bar{m}_-$ be the midpoints of~$[\bar{x},\bar{x}_+]$ and~$[\bar{x},\bar{x}_-]$. Denote by~$\pi(\bar{m}_+)$ and~$\pi(\bar{m}_-)$ the projections of~$\bar{m}_+$ and~$\bar{m}_-$ to~$\T^2$, where $\pi:\R^2 \to \T^2$ is the covering projection. Suppose that the segment~$c$ is directed along the irrational direction~$\theta$. If the segment~$c$ is long enough (depending on~$\theta$ and~$\varepsilon$), its projection to~$\T^2$ passes at (vertical) distance less than~$\frac{\varepsilon}{2}$ from any point, and in particular from~$\pi(\bar{m}_+)$ and~$\pi(\bar{m}_-)$. Thus, the segment~$c$ passes at vertical distance less than~$\frac{\varepsilon}{2}$ from some $\Z^2$-translates~$m_+$ and~$m_-$ of~$\bar{m}_+$ and~$\bar{m}_-$ (\ie, some lifts of~$\pi(\bar{m}_+)$ and~$\pi(\bar{m}_-)$). Define~$x_+$ and~$x_-$ as the corresponding $\Z^2$-translates of~$\bar{x}_+$ and~$\bar{x}_-$. By construction, the segment~$c$ passes below~$x_+$ and above~$x_-$ at vertical distance less than~$\varepsilon$ from them. The same holds for a segment~$c$ of direction close enough to $\theta$.
\end{proof}

Denote by~$(\vec{\imath},\vec{\jmath})$ the canonical basis of~$\R^2$. The following lemma is a quantitative version of the fact that the curve shortening flow straighten curves.

%\forget
\begin{comment}
The following lemma shows that the tangent vector of a non-vertical geodesic in a Finsler two-torus without conjugate points stays away from a vertical direction.


\begin{lemm}\label{lem:horiz}
Assume that the horizontal and vertical curves of~$M$ are geodesic. For every \mbox{$\nu >0$}, there exists $C_\nu \in (0,1)$ such that every geodesic~$\gamma$ of~$\bar{M}$ whose asymptotic direction forms an angle lying between~$-\nu$ and~$\nu$ with the horizontal direction satisfies $|\langle \gamma',\vec{\jmath} \rangle| \leq C_\nu \, \lVert \gamma' \rVert$.
\end{lemm}


\begin{proof}
To complete\dots
\end{proof}
%\forgotten
\end{comment}


\begin{lemm}\label{lem:L'}
Assume that the horizontal and vertical curves of~$M$ are geodesic. For every $\varepsilon \in (0,1)$ and every~$L>0$, there exists $t_0=t_0(\varepsilon,L)>0$ such that for every $t \geq t_0$ and every geodesic~$\gamma$ of~$\bar{M}$ whose asymptotic direction forms an angle lying between~$-\frac{\pi}{4}$ and~$\frac{\pi}{4}$ with the horizontal direction, the following holds: every arc of~$\gamma_t$ over an interval~$I$ of length~$L$ of the $x$-axis is at vertical distance at most~$\varepsilon$ from a segment of~$\R^2$ over the same interval~$I$.
\end{lemm}

\begin{proof}
By Theorem~\ref{theo:hedlund}, the geodesic~$\gamma$ lies in a strip~$S$ of width~$w$ with the same asymptotic direction, where $w$ does not depend on~$\gamma$ (only on the Finsler metric on~$M$). Since the vertical curves of~$\bar{M}$ are geodesic, the tangent vector~$\gamma'$ is uniformly bounded away from a vertical direction, that is, $|\langle \gamma',\vec{\jmath} \rangle| \leq C \, \lVert \gamma' \rVert$ for some constant $C \in (0,1)$ not depending on the geodesic~$\gamma$. Otherwise, the geodesic~$\gamma$ would be close to a vertical geodesic and would leave the strip~$S$ whose vertical width is bounded.

By Lemma~\ref{lem:curv}, there exists~$c_0>0$ not depending on~$\gamma$ such that the curvature~$\kappa(\gamma_t)$ of~$\gamma_t$ satisfies
\[
|\kappa(\gamma_t)| \leq \frac{c_0}{\sqrt{t}}.
\]
Thus, for $t$ large enough, the curve~$\gamma_t$ is almost straight. Hence the desired result.
\end{proof}

We can now proceed to the proof of Theorem~\ref{theo:unifwidth}.

\begin{proof}[Proof of Theorem~\ref{theo:unifwidth}]
By Proposition~\ref{prop:identification}, we can assume that the horizontal and vertical curves of~$\T^2$ are geodesics. Switching the roles of the $x$- and~$y$-axis, and the orientation of~$\gamma$ if necessary, we can further assume that the angle between the asymptotic direction of~$\gamma$ and the horizontal vector~$\vec{\imath}$ lies between~$-\frac{\pi}{4}$ and~$\frac{\pi}{4}$. Recall that the curve~$\gamma_t$ lies in a (closed) Euclidean strip~$S_t=S(\gamma_t)$ of~$\R^2$ bounded by two straight lines~$\Delta_+=\Delta_+(\gamma_t)$ and~$\Delta_-=\Delta_-(\gamma_t)$ with the same asymptotic direction as~$\gamma$, with $\Delta_+$ above~$\Delta_-$ in~$\R^2$; see Theorem~\ref{theo:hedlund}.

Let us show that if the asymptotic direction of~$\gamma$ is close enough to~$\theta$ (more precisely, if the asymptotic direction of~$\gamma$ is in the neighborhood of~$\theta$ given by Lemma~\ref{lem:dioph}), then the width of~$S_t$ is less than~$14\varepsilon$ for $t \geq t_0$. Take a point~$p_+=p_{+}(\gamma_t)$ of~$\gamma_t$ at vertical distance at most~$\varepsilon$ from~$\Delta_+$ (below~$\Delta_+$) and denote by~$\bar{p}_+$ its projection to the $x$-axis. Let $I$ be the interval of the $x$-axis centered at~$\bar{p}_+$ of length~$L$, where $L$ is given in Lemma~\ref{lem:dioph}. Consider the arc~$\alpha_+$ of~$\gamma_t$ over~$I$. By Lemma~\ref{lem:L'}, this arc~$\alpha_+$ is at vertical distance at most~$\varepsilon$ from a segment~$c$ of~$\R^2$. Thus, the segment~$c$ lies below the line~$\Delta_+ +\varepsilon \vec{\jmath}$ above~$\Delta_+$ at vertical distance~$\varepsilon$ from~$\Delta_+$ and passes through a point at vertical distance at most~$\varepsilon$ from~$p_+$. Therefore, the segment~$c$ lies at vertical distance at most~$6\varepsilon$ from~$\Delta_++\varepsilon \vec{\jmath}$. We deduce that the arc~$\alpha_+$ of~$\gamma_t$ lies below~$\Delta_+$ at vertical distance at most~$6\varepsilon$ from~$\Delta_+$. See~Figure~\ref{fig:relative}. 


\begin{figure}[htbp!]
%\vspace{1cm}
\newcommand*\svgwidth{6cm}
\centering
\includegraphics[scale=1.8]{Strip.pdf}
\caption{Relative positions of the curves $\alpha_+$, $c$ and~$\Delta_+$}\label{fig:relative}
\end{figure}


Apply Lemma~\ref{lem:dioph} to the segment lying in~$\Delta_+ - 6 \varepsilon \vec{\jmath}$ above the interval~$I$ of length~$L$, assuming the asymptotic direction of~$\Delta_+$ is close enough to the irrational direction~$\theta$. This yields a point~$x_-$ of~$\Z^2$ below~$\Delta_+-6 \varepsilon \vec{\jmath}$ at vertical distance at most~$\varepsilon$ from it and so at vertical distance at most~$7\varepsilon$ from~$\Delta_+$. It follows that the geodesic arc~$\alpha_+$ of~$\gamma_t$ passes above the point~$x_-$ of~$\Z^2$ at vertical distance at most~$7\varepsilon$ from~$\Delta_+$. Similarly, the curve~$\gamma_t$ passes below a point~$x_+$ of~$\Z^2$ at vertical distance at most~$7\varepsilon$ from~$\Delta_-$.

Let $T$ be the translation of~$\R^2$ by the integral vector~$x_--x_+ \in \Z^2$. Note that $T$ is an isometry both for the Euclidean metric and the Finsler metric. By construction, the point~$x_-$ lies below~$\gamma_t$ and above $T(\gamma_t)$, and is at vertical distance at most~$7\varepsilon$ from~$\Delta_+$ and~$T(\Delta_-)$. The geodesics~$\gamma$ and~$T(\gamma)$ have the same asymptotic direction and do not intersect. The same holds for their images~$\gamma_t$ and~$T(\gamma_t)$ under the curve shortening flow; see Section~\ref{sec:convergenceCSF}. In addition, $T(\gamma_{t})$ is proved to be below $\gamma_{t}$ and $T(\Delta_{-})$ is below $\Delta_{-}$, hence $T(\Delta_{-})$ bounds $\gamma_{t}$ from below. It follows that the curve~$\gamma_t$ lies in the strip of width at most~$14 \varepsilon$ bounded by~$\Delta_+$ on top and by~$T(\Delta_-)$ at the bottom. Thus, $\width(\gamma_t) \leq 14 \varepsilon$.
\end{proof}


%\forget
\begin{comment}
\begin{proof}
By Proposition~\ref{prop:identification}, we can assume that the horizontal and vertical curves of~$\T^2$ are geodesic. Switching the roles of the $x$- and~$y$-axis, and the orientation of~$\gamma$ if necessary, we can further assume that the angle between the asymptotic direction of~$\gamma$ and the horizontal vector~$\vec{\imath}$ lies between~$-\frac{\pi}{4}$ and~$\frac{\pi}{4}$. By Lemma~\ref{lem:horiz}, this implies that the tangent vector~$\gamma'$ remains uniformly away from a vertical direction, that is, $|\langle \gamma',\vec{\jmath} \rangle| \leq C \, ||\gamma'||$ for some constant $C=C_{\pi/4} \in (0,1)$ not depending on the geodesic~$\gamma$.

Let $\varepsilon >0$. By~\cite[Proposition~4.4]{EH89}, there exists~$C_0>0$ not depending on~$\gamma$ such that the curvature~$K(\gamma_t)$ of~$\gamma_t$ satisfies
\[
|K(\gamma_t)| \leq \frac{C_0}{\sqrt{t}}.
\]
This immediately implies the following result.

\begin{claim} \label{claim:segment}
There exists $t_0>0$ (not depending on~$\gamma$) such that for every $t \geq t_0$, every arc of~$\gamma_t$ over an interval~$I$ of length~$L'$ of the $x$-axis is at vertical distance at most~$\varepsilon$ from a segment of~$\R^2$ over the same interval~$I$.
\end{claim}


Recall that the minimizing Finsler geodesic~$\gamma_t$ lies in a minimal (closed) Euclidean strip~$S_t$ of~$\R^2$ bounded by two straight lines~$\Delta_+=\Delta^t_+$ and~$\Delta_-=\Delta^t_-$ with the same asymptotic direction as~$\gamma$, with $\Delta_+$ above~$\Delta_-$ in~$\R^2$; see Theorem~\ref{theo:hedlund}.

Let us show that for every geodesic~$\gamma$ of asymptotic direction close enough to~$\theta$, the width of~$S_t$ is less than~$14\varepsilon$ for $t \geq t_0$. Take a point~$x_+$ of~$\gamma_t$ at vertical distance at most~$\varepsilon$ from~$\Delta_+$ (below~$\Delta_+$) and denote by~$\pi(x_+)$ its projection to the $x$-axis. Let $I$ be the interval of the $x$-axis centered at~$\pi(x_+)$ of length~$L''=\max\{L,L'\}$, where $L$ and~$L'$ are given in Lemma~\ref{lem:dioph} and Claim~\ref{claim:segment}. Consider the arc~$\alpha_+$ of~$\gamma_t$ over~$I$. By Claim~\ref{claim:segment}, this arc~$\alpha_+$ is at vertical distance at most~$\varepsilon$ from a segment~$c$ of~$\R^2$. Thus, the segment~$c$ lies below the line~$\Delta_+ +\varepsilon \vec{\jmath}$ above~$\Delta_+$ at vertical distance~$\varepsilon$ from~$\Delta_+$ and passes through a point at vertical distance at most~$\varepsilon$ from~$x_+$. Therefore, the segment~$c$ lies at vertical distance at most~$6\varepsilon$ from~$\Delta_++\varepsilon \vec{\jmath}$. We deduce that the arc~$\alpha_+$ of~$\gamma_t$ lies below~$\Delta_+$ at vertical distance at most~$6\varepsilon$ from~$\Delta_+$.

\begin{figure}[htbp!]
\vspace{1cm}
\newcommand*\svgwidth{6cm}
\failedinput{dessin.pdf_tex}
\caption{Relative positions of the curves $\alpha_+$, $c$ and~$\Delta_+$}\label{fig:relative}
\end{figure}

Apply Lemma~\ref{lem:dioph} to the segment lying in~$\Delta_+ - 6 \varepsilon \vec{\jmath}$ above the interval~$I$ of length~$L'' \geq L'$, assuming the asymptotic direction of~$\Delta_+$ is close enough to the irrational direction~$\theta$. This yields a point~$m_+$ of~$\Z^2$ below~$\Delta_+-6 \varepsilon \vec{\jmath}$ at vertical distance at most~$\varepsilon$ from it and so at vertical distance at most~$7\varepsilon$ from~$\Delta_+$. It follows that the geodesic arc~$\alpha_+$ of~$\gamma_t$ passes above the point~$m_+$ of~$\Z^2$ at vertical distance at most~$7\varepsilon$ from~$\Delta_+$. Similarly, the geodesic~$\gamma_t$ passes below a point~$m_-$ of~$\Z^2$ at vertical distance at most~$7\varepsilon$ from~$\Delta_-$.

Let $T$ be the translation of~$\R^2$ by the integral vector~$m_+-m_- \in \Z^2$. By construction, the point~$m_+$ lies below~$\gamma_t$ and above $T(\gamma_t)$, and is at vertical distance at most~$7\varepsilon$ from~$\Delta_+$ and~$T(\Delta_-)$. It follows that the geodesic~$\gamma_t$ lies in the strip of width at most~$14 \varepsilon$ bounded by~$\Delta_+$ on top and by~$T(\Delta_-)$ at the bottom. Thus, $\width(\gamma_t) \leq 14 \varepsilon$.
\end{proof}
\end{comment}
%\forgotten

By Theorem~\ref{theo:affine} and Proposition~\ref{prop:direction}, a geodesic in the universal cover of a Finsler two-torus converges to an (oriented) straight line of~$\R^2$ with the same asymptotic direction. An oriented straight line~$\Delta$ of~$\R^2$ is determined by its asymptotic direction~$\theta(\Delta)$ and its signed distance~$p(\Delta)$ to the origin in~$\R^2$; see Section~\ref{sec:crofton} for a more detailed description. The following result which is a consequence of the previous proposition, is about the regularity of the signed distance at the limit when the geodesic varies.


\begin{prop}\label{prop:pcontinuous}
Let $M=(\T^2,F)$ be a Finsler two-torus without conjugate points. Denote by~$\bar{M}=(\R^2,\bar{F})$ its universal Finsler cover. Then the function
\[
\begin{array}{rccc}
p_\infty : & U\bar{M} & \to & \R \\
& v & \mapsto & p\left(\gamma_v^\infty\right)
\end{array}
\]
is continuous at vectors~$v$ generating a geodesic~$\gamma_v$ with irrational asymptotic direction.
\end{prop}

\begin{proof}
Fix a vector~$v_0$ generating a geodesic~$\gamma_{v_0}$ with irrational asymptotic direction. Let $\varepsilon >0$. By Theorem~\ref{theo:unifwidth}, we can choose~$t_0>0$ such that for every geodesic~$\gamma_v$ of~$\bar{M}$ with $v \in U\bar{M}$ close enough to~$v_0$, the curve~$\gamma_v^{t_0}$ lies in a (closed) Euclidean strip of width less than~$\varepsilon$. Since its limit~$\gamma_v^\infty$ lies in the same strip, the Hausdorff distance between~$\gamma_v^{t_0}$ and~$\gamma_v^\infty$ in~$\R^2$ is less than~$\varepsilon$.

For $v \in U\bar{M}$ close enough to~$v_0$, the curve~$\gamma_v^{t_0}$ is at Hausdorff distance less than~$\varepsilon$ from~$\gamma_{v_0}^{t_0}$ on any given compact subset of~$\R^2$. It follows that for $v \in U\bar{M}$ close enough to~$v_0$, the Hausdorff distance between the limit straight lines~$\gamma^\infty_v$ and~$\gamma^\infty_{v_0}$ on any given compact subset of~$\R^2$ is less than~$3\varepsilon$. Thus,
\[
\left|p\left(\gamma_v^\infty\right)-p\left(\gamma_{v_0}^\infty\right)\right| \leq \left|p\left(\gamma_v^\infty\right)-p\left(\gamma_{v}^{t_0}\right)\right| + \left|p\left(\gamma_{v}^{t_0}\right)-p\left(\gamma_{v_0}^{t_0}\right)\right| + \left|p\left(\gamma_{v_0}^{t_0}\right)-p\left(\gamma_{v_0}^\infty\right)\right| < 3\varepsilon.
\]
Hence, the function~$p_\infty:U\bar{M} \to \R$ is continuous at~$v_0$.
\end{proof}

It is unclear whether the function $p_\infty:U\bar{M} \to \R$ is continuous at vectors generating geodesics pointing in rational directions. The following example illustrates possible issues. Note however that these issues may not occur for geodesics in the universal cover of a Finsler torus without conjugate points.

\begin{exam}\label{ex:noncontinuous}
We can construct a smooth family of graphs asymptotic to the lines $y=mx$ converging to the horizontal line $y=1$ in the smooth topology on compact sets, as $m$ goes to zero. Such a family converges to a non-continuous family of lines under the curve shortening flow, namely the family formed of the lines $y=mx$ for $m>0$ and $y=1$ for $m=0$. See Figure~\ref{fig:non-continuity}.

\begin{figure}[htbp!]
\newcommand*\svgwidth{7cm}
\begin{center}
\includegraphics[scale=0.5]{dessin-1.pdf}
\end{center}
\caption{Non-continuity of the limit under the curve shortening flow.}\label{fig:non-continuity}
\end{figure}
\end{exam}

\section{Analytic expression of the limit under the curve shortening flow}

In this section, we derive the following analytic expression of the limit under the curve shortening flow of a function whose graph represents a (minimizing) geodesic in the universal cover of a Finsler two-torus without conjugate points; see Theorem~\ref{theo:affine}.

\begin{prop}\label{prop:limit}
Let $M=(\T^2,F)$ be a Finsler two-torus without conjugate points. Denote by~$\bar{M}=(\R^2,\bar{F})$ its universal Finsler cover. Let $\bu:\R \to \R$ be a smooth function whose graph~$\GG$ in~$\R^2$ is a (minimizing) Finsler geodesic of~$\bar{M}$ with the same asymptotic direction as the line $y=ax$. Then, the limit~$\bu_\infty$ of~$\bu$ under the curve shortening flow~\eqref{eq:GCSF} satisfies
\[
\bu_\infty(0) = \lim_{R\, \to\,+\infty} \frac{1}{2R} \int_{-R}^R \bu(x) - ax \, dx
\]
and, more generally,
\begin{equation}\label{eq:limit+}
\bu_\infty(x_0) = \lim_{R\,\to\,+\infty} \frac{1}{2R} \int_{-R}^R \bu(x+x_0) - ax \, dx
\end{equation}
for every $x_0 \in \R$. In particular, the limits exist.

Actually, the convergence in~\eqref{eq:limit+} is uniform, that is,
\[
\lim_{R\,\to\,+\infty} \sup_{x_0\,\in\,\R} \frac{1}{2R} \left| \int_{-R}^R \bu(x+x_0) - ax - \bu_\infty(x_0) \, dx \right| =0.
\]
\end{prop}

\begin{proof}
Assume that $a=0$. By Theorem~\ref{theo:affine}, the solution of the curve shortening flow~\eqref{eq:GCSF} with initial condition~$\bu$ converges to a constant function $\bu_{\infty}=\bu_{\infty}(0)$. Now, recall that the solution of the curve shortening flow equation~\eqref{eq:GCSF} converges uniformly to the solution of the heat equation with the same initial condition; see~\cite[Theorem~1.4]{NT07}. Thus, the solution of the heat equation with initial condition~$\bu$ uniformly converges to the same constant function $\bu_{\infty}=\bu_{\infty}(0)$. Since the initial condition~$\bu$ is bounded, see Lemma~\ref{lem:strip}, it follows from the expression of the limit of the solution of the heat equation, see~\cite{RE66, RE67}, that
\[
\bu_\infty(0) = \lim_{R\,\to\,+\infty} \frac{1}{2R} \int_{-R}^R \bu(x) \, dx.
\]
See also~\cite[\S~3]{NT07}. By a change of variable, we obtain~\eqref{eq:limit+}. Actually, since the solution of the heat equation with the initial condition~$\bu$ uniformly converges, the convergence in~\eqref{eq:limit+} is uniform.

In the general case (\ie, when~$a$ is not necessarily zero), we apply the same argument to~$\bv(x)=\bu(x) - ax$.
\end{proof}


\begin{rema}
By Theorem~\ref{theo:affine}, the limit function~$\bu_\infty$ is an affine function of the form~$\bu_\infty(x)=ax+b$. It follows from Proposition~\ref{prop:limit} that
\begin{align}
a = \bu_\infty(1) - \bu_\infty(0) &= \lim_{R\,\to\,+\infty} \frac{1}{2R} \int_{-R}^R \bu(x+1) - \bu(x) \, dx\label{eq:lima}
\\
\intertext{and}
b = \bu_\infty(0) &= \lim_{R\,\to\,+\infty} \frac{1}{2R} \int_{-R}^R \bu(x) - ax \, dx.\label{eq:limb}
\end{align}
Thus, the limit affine function~$\bu_\infty:\R \to \R$ is given by the limits of some linear integrals of~$\bu$.
\end{rema}


\begin{rema}
As for Theorem~\ref{theo:affine}, the conclusion of Proposition~\ref{prop:limit} still holds in the rational case when $M$ has conjugate points (as long as the geodesic represented by~$\clG$ is minimizing).
\end{rema}


\section{Curve shortening flow and unit bundle diffeomorphism} \label{sec:diffeo}

We show that the deformation of the geodesic foliation of a Finsler two-torus without conjugate points under the curve shortening flow induces a family of diffeomorphisms on the unit tangent bundle of the torus. This gives rise to a deformation of the Heber foliation on the cotangent bundle.

Let $M=(\T^2,F)$ be a Finsler two-torus without conjugate points, where $\T^2=\R^2/\Z^2$. Denote by~$\bar{M}=(\R^2,\bar{F})$ the universal Finsler cover of~$M$. The isometric action of~$\Z^2$ by deck transformations on~$\bar{M}$ induces a natural action on~$U\bar{M}$ and~$T\bar{M}$.

\begin{defi}\label{def:Psi}
Consider the map
\[
\overline{\Psi}_t:U \bar{M} \to T \R^2 \setminus \{0\}
\]
defined as
\[
\overline{\Psi}_t(v) = (\gamma_v^t)'(0)
\]
for every $v \in U \bar{M}$ and~$t \in [0,\infty)$, where $(\gamma_v^t)$ is the Euclidean curve shortening flow of the Finsler geodesic~$\gamma_v$ induced by~$v$; see Definition~\ref{def:CSF}. Note that $\overline{\Psi}_0$ is the identity map on~$U \bar{M}$. Since $\Z^2$ acts by isometries both on~$\bar{M}$ and~$\R^2$, the map~$\overline{\Psi}_t$ is $\Z^2$-equivariant.

%\forget
\begin{comment}
The unit tangent bundle~$U\T^2$ of the Finsler metric on~$\T^2$ naturally identifies with the unit tangent bundle~$U_0\T^2$ of the canonical metric on the square flat torus~$\T^2=\R^2/\Z^2$ by radial projection on each tangent plane. That is, $U\T^2 \simeq U_0 \T^2$.
\end{comment}
%\forgotten

Consider also the map
\[
\Psi_t:U\bar{M} \to U_0 \R^2
\]
to the unit tangent bundle~$U_0\R^2$ of the Euclidean plane~$\R^2$ defined as
\[
\Psi_t(v) = \pi\left[\overline{\Psi}_t(v)\right] = \pi\left[(\gamma_v^t)'(0)\right]
\]
where $\pi:T\bar{M} \setminus \{0\} \to U_0 \R^2$ is the radial projection onto the unit circle of each tangent plane of the Euclidean plane~$\R^2 = \bar{M}$. The map~$\Psi_t$ is also $\Z^2$-equivariant and its quotient map is a $\pi_1$-isomorphism since it is the radial projection of the identity map for~$t=0$.
\end{defi}

The following result is derived from an analysis of the (parabolic) partial differential equation satisfied by the curve shortening flow. This is the analogue of~\cite[Proposition~4.4]{Sa19} for the family of diffeomorphisms of the unit tangent bundle of the round projective plane obtained by applying the curve shortening flow to the geodesics of a Zoll Finsler metric.


\begin{theo}\label{theo:Psi}
Let $M=(\T^2,F)$ be a Finsler torus without conjugate points. For every $t \in [0,+\infty)$, the map $\Psi_t:U\bar{M} \to U_0\R^2$ is a diffeomorphism.

Its quotient map, still denoted by $\Psi_t:UM \to U_0\T^2$, is also a diffeomorphism.
\end{theo}



\begin{proof}
First, let us show that the map~$\overline{\Psi}_{t}:U\bar{M}\to T\R^{2}$ is an immersion.

For $t=0$, this is true since, by construction, for every $v \in U\bar{M}$,
\[
\overline{\Psi}_0(v) = \gamma_v'(0) = v.
\]

Assume~$t>0$. Let $v_0 \in U\bar{M}$ and $\zeta \in T_{v_0}U\bar{M}$. Let $v=v(\lambda)$ be a smooth vector variation in~$U\bar{M}$ with~$v(0)=v_0$ and~$v'(0)=\zeta$. We want to show that for every $t>0$, the differential
\[
d\overline{\Psi}_t(v_0): T_{v_0} U \bar{M} \to T_{\overline{\Psi}_t(v_0)} T \R^2
\]
of~$\overline{\Psi}_t$ at~$v_0$ is injective. That is, if the derivative~$d\overline{\Psi}_t(v_0)(\zeta)$ of~$\overline{\Psi}_t(v(\lambda))$ vanishes at~$\lambda=0$, then the vector~$\zeta=v'(0)$ is zero.

By Proposition~\ref{prop:identification}, we can assume that the horizontal and vertical lines of~$\bar{M}$ are geodesics. Switching the horizontal and vertical lines if necessary, we can assume that the vector~$v_0$ is nonvertical and therefore each geodesic~$\gamma_\lambda = \gamma_{v(\lambda)}$ of~$\bar{M}$ is represented in~$\bar{M}=\R^2$ as the graph of some smooth function~$\bu_0(.,\lambda)$ for $\lambda$ close enough to zero. The map $(x,y) \mapsto (x,y-\bu_0(x,0))$ is a diffeomorphism of~$\R^2$ taking vertical lines to vertical lines and the geodesic~$\gamma_0$ to the $x$-axis. In this new coordinate system, every curve~$\gamma_{\lambda}^t$ with $\lambda$ close enough to zero obtained by applying the curve shortening flow to~$\gamma_\lambda$, see~\eqref{eq:CSF}, is represented as the graph
\[
\{ (x,\bu(x,t,\lambda)) \, \mid \, x \in \R \}
\]
of a smooth function~$\bu(.,t,\lambda)$ with $\bu(.,0,0)=0$. More specifically,
\[
\gamma_\lambda^t(s) = (x(s,t,\lambda),\bu(x(s,t,\lambda),t,\lambda))
\]
where $s$ is the arclength-parameter of~$\gamma_\lambda^t$ for the Finsler metric~$\bar{F}$ (with an orientation compatible with~$x$). Note that the partial derivative~$x_s$ does not vanish (and is positive). Observe also that $\bu_x(.,0,0)=\bu_{xx}(.,0,0)=0$.

In the new coordinate system, the equation of the curve shortening flow is no longer given by~\eqref{eq:GCSF}. It satisfies a new equation which has been derived in~\cite[Eq.~(3.2)]{An90} and~\cite[Appendix]{gage90}. More precisely, the function
\[
\bu : \R \times [0,\infty) \times (-\varepsilon,\varepsilon) \to \R
\]
satisfies the following (parabolic) partial differential equation of the curve shortening flow
\begin{equation}\label{eq:utF}
\bu_t = \FF\left(x,\bu, \bu_x, \bu_{xx}\right)
\end{equation}
where $\FF$ is a smooth function defined on~$\R^4$ with
\[
\FF_q(x,u,p,q) >0
\]
which can be expressed in terms of the coefficients of the Euclidean metric in the new coordinate system (where $\gamma_0$ coincides with the $x$-axis). Here, the subscripts refer to the partial differentiations. Note that we do not need an explicit form of $\FF$.

%\forget
\begin{comment}
Under the curve shortening flow~\eqref{eq:CSF}, the curve~$\gamma_{\lambda}^t$ is represented as the graph
\[
\{ (x,\bu(x,t,\lambda)) \, \mid \, x \in \R \}
\]
where $(x,t) \mapsto \bu(x,t,\lambda)$ is the solution of the graphical curve shortening flow~\eqref{eq:GCSF} with initial condition~$\bu_0(.,\lambda)$. More specifically,
\[
\gamma_\lambda^t(s) = (x(s,t,\lambda),\bu(x(s,t,\lambda),t,\lambda))
\]
where $s$ is the arclength-parameter of~$\gamma_\lambda^t$ for the Finsler metric and $\bu(x,0,\lambda) = \bu_0(x,\lambda)$. Note that the partial derivative~$x_s$ does not vanish.
\end{comment}
%\forgotten

The tangent vector $\overline{\Psi}_t(v(\lambda)) = (\gamma_{\lambda}^t)'(0) \in U_0 \R^2$ is given by
\begin{equation}\label{eq:uuuu}
\overline{\Psi}_t(v(\lambda)) =
\left(x,\bu,x_s,x_s \, \bu_x \right)
\end{equation}
where $x_*=x_*(0,t,\lambda)$ and $\bu_*=\bu_*(x(0,t,\lambda),t,\lambda)$. Taking the derivative of~\eqref{eq:uuuu} with respect to~$\lambda$, we obtain for~$\lambda=0$
\begin{equation}\label{eq:dPsi}
d\overline{\Psi}_t(v_0)(\zeta) = \left(x_\lambda, x_\lambda \, \bu_x + \bu_\lambda, x_{s \lambda}, x_{s \lambda} \, \bu_x + x_s \, x_\lambda \, \bu_{xx} + x_s \, \bu_{x \lambda} \right)
\end{equation}
where $x_* = x_*(0,t,0)$ and $\bu_* = \bu_*(x(0,t,0),t,0)$. 


Suppose that $d\overline{\Psi}_\tau(v_0)(\zeta)=0$. Recall that $x_s$ does not vanish. In this case, the functions $x_\lambda$, $\bu_\lambda$, $x_{s \lambda}$ and~$\bu_{x \lambda}$ vanish at $s=0$, $t=\tau$, $\lambda=0$.
%~$(x(0),\tau,0)$.
In particular, the function $\bv=\bu_\lambda$ has a multiple zero at~$(x(0,\tau,0),\tau,0)$, \ie, both $\bv$ and~$\bv_x$ vanish at this point.

Using the fact that $\bv$ has a multiple zero and that the Finsler metric has no conjugate points, our goal is to show that $\zeta=0$. 

The following result shows that $\bv$ vanishes when~$\lambda=0$.

\begin{lemm}\label{lem:v}
For every $x \in \R$ and $t \geq 0$, we have
\[
\bv(x,t,0)=0.
\]
\end{lemm}

\begin{proof}
Let us derive the evolution equation of~$\bv$. Differentiating the relation~\eqref{eq:utF} with respect to~$\lambda$ leads to the following parabolic partial differential equation
\begin{equation}\label{eq:vt}
\bv_t = a(x,t,\lambda) \, \bv_{xx} + b(x,t,\lambda) \, \bv_x + c(x,t,\lambda) \, \bv
\end{equation}
where $a=\FF_q(x,\bu,\bu_x,\bu_{xx})$, $b=\FF_p(x,\bu,\bu_x,\bu_{xx})$ and $c=\FF_u(x,\bu,\bu_x,\bu_{xx})$.

%\forget 
\begin{comment}
Differentiating the quasilinear parabolic equation~\eqref{eq:GCSF}
\[
\bu_t = \frac{\bu_{xx}}{1+\bu_x^2}
\]
satisfied by~$\bu = \bu(x,t,\lambda)$ with respect to~$\lambda$, we obtain for $\lambda=0$ the following parabolic partial differential equation
\begin{equation}\label{eq:vt}
\bv_t = a(x,t) \, \bv_{xx} + b(x,t) \, \bv_x
\end{equation}
where $a=\frac{1}{1+\bu_x^2} >0$ and $b=-\frac{2 \, \bu_x \, \bu_{xx}}{(1+\bu_x^2)^2}$. Here, $\bu_x = \bu_x(x,t,0)$ and $\bu_{xx} = \bu_{xx}(x,t,0)$.
\end{comment}
%\forgotten


By Lemma~\ref{lem:curv} (still for $\lambda=0$), the function~$\bu$ has linear growth with respect to~$x \in \R$ and all the partial derivatives of~$\bu$ (and so $a$, $a^{-1}$, $b$ and all their partial derivatives) are uniformly bounded with respect to~$x \in \R$ and $t \geq 0$.

It follows from~\cite[Theorem~B]{An88} that the function~$\bv(.,t,0)$ has at least two zeros for $t < \tau$ (the number of zeros does not increase under the flow and~$\bv(.,\tau,0)$ has a double zero). In particular, $\bv(.,0,0)$ vanishes at least twice.



Consider the geodesic variation
\[
\gamma_\lambda(s) = (x(s,0,\lambda), \bu(x(s,0,\lambda),0,\lambda)).
\]
Taking the derivative with respect to~$\lambda$ and using that $\bu_x(.,0,0)=0$, we obtain for $\lambda=0$ the Jacobi field
\[
J=(x,\bu,x_\lambda, \bv)
\]
where $x_* = x_*(s,0,0)$, $\bu = \bu(x(s,0,0),0,0)$ and $\bv = \bu_\lambda = \bu_\lambda(x(s,0,0),0,0)$. Since $s \mapsto x(s,0,0)$ is a diffeomorphism and $\bv(.,0,0)$ vanishes twice, the Jacobi field~$J$ is parallel to the horizontal axis~$\gamma_0$, at two points~$s_1$ and~$s_2$. That is, $J(s_i) = \alpha_i \, \gamma_0'(s_i)$ with $\alpha_i \in \R$, for $i=1,2$.

Consider the decomposition
\[
J = \left(\frac{s_2-s}{s_2-s_1} \alpha_1 + \frac{s_1-s}{s_1-s_2} \alpha_2 \right) \gamma' + J_\perp.
\]
By construction, $J_\perp$ is a Jacobi field along~$\gamma_0$ which vanishes twice, namely at~$s_1$ and~$s_2$. Since the Finsler metric has no conjugate points, the Jacobi field~$J_\perp$ is trivial. Thus, $J$ is parallel to~$\gamma_0'$ and $\bv(.,0,0)$ is constant equal to zero. Now, since $\bv$ satisfies the parabolic partial differential equation~\eqref{eq:vt}, we deduce from the initial condition~$\bv(.,0,0)=0$ that $\bv(.,.,0)$ is zero.
\end{proof}

We can now derive the desired result.

\begin{lemma}\label{lem:xi}
We have $\zeta=0$.
\end{lemma}

\begin{proof}
By Lemma~\ref{lem:v}, the function~$\bv(.,.,0)$ and its derivative~$\bv_x(.,.,0)$ are zero. In particular, $\bu_\lambda(.,0,0)$ and~$\bu_{x \lambda}(.,0,0)$ are zero. Besides, we also have $\bu_x(.,0,0)=\bu_{xx}(.,0,0)=0$. It follows from the relation~$d\overline{\Psi}_0(v_0)(\zeta)=\zeta$ and the expression~\eqref{eq:dPsi} that
\[
\zeta = d\overline{\Psi}_0(v_0)(\zeta) = (x_\lambda, 0, x_{s \lambda}, 0)
\]
where $x_*=x_*(0,0,0)$. By our choice of coordinate, the geodesic~$\gamma_0$ agrees with the $x$-axis and its tangent vector~$v_0$ is horizontal. Now, since $\zeta \in T_{v_0} U\R^2$ is tangent at~$v_0$ to the unit tangent bundle of~$\R^2$ with the Finsler metric, the vector~$(x_{s \lambda},0)$ of~$\R^2$ formed of the last two coordinates of~$\zeta$ in~$\R^4$ is not colinear to~$v_0$ unless it is trivial. Since both vectors are horizontal, it follows that $x_s(0,0,0)=0$.

To show that $x_\lambda$ also vanishes at~$(0,0,0)$, we consider a smooth variation of \emph{horizontal} vectors~$v(\lambda)$ with $v(0)=v_0$ and~$v'(0)=\zeta$ (recall that both~$v_0$ and~$\zeta$ are horizontal). Since $v(\lambda)$ is tangent to the geodesic~$\gamma_0$, the induced geodesic variation~$\gamma_\lambda$ is given by a change of variable
\[
\gamma_\lambda(s) = \gamma_0(s+\bar{s}(\lambda))
\]
where $\bar{s}:\R \to \R$ is a smooth function with~$\bar{s}(0)=0$ so that
\begin{equation}\label{eq:vL}
\gamma'_\lambda(0) = \gamma'_0(\bar{s}(\lambda)) = v(\lambda).
\end{equation}
Following the curve shortening flow, this implies that
\[
\gamma_\lambda^t(s) = \gamma_0^t(s+\bar{s}(\lambda)) = \gamma_0(s+\bar{s}(\lambda)).
\]
In particular,
\[
\overline{\Psi}_t(v(\lambda)) = (\gamma_\lambda^t)'(0) = \gamma_0'(\bar{s}(\lambda)).
\]
Taking the derivative of this expression with respect to~$\lambda$, we obtain for~$\lambda=0$
\[
d\overline{\Psi}_t(v_0)(\zeta) = \bar{s}'(0) \, \frac{\partial}{\partial s} (\gamma_0')_{|s=0}.
\]
Observe that the vector~$\frac{\partial}{\partial s} (\gamma_0')_{|s=0}$ of~$T_{\overline{\Psi}_t(v_0)} T \R^2$ is nonzero since the vector~$\gamma_0'(0)$ formed of its first two coordinates is nonzero.

Now, by assumption, $d\overline{\Psi}_\tau(v_0)(\zeta)=0$. This implies that $\bar{s}'(0)=0$. Differentiating the equation~\eqref{eq:vL} with respect to~$\lambda$ and plugging in~$\lambda=0$, we derive that $\zeta=v'(0)$ is zero.
\end{proof}

It follows from Lemma~\ref{lem:xi} that the map $\overline{\Psi}_t:U\bar{M} \to T \R^2 \setminus\{0\}$ is an immersion.

Let us show that this immersion is transverse to the rays~\mbox{$\R_+^* u = \{ su \mid s >0\}$}, where the vector~$u$ runs over~$U\R^2$. By contradiction, assume that there exist $v \in U\bar{M}$ and a nonzero vector $\zeta \in T_v U\bar{M}$ such that the unit tangent vector~$u$ pointing to~$\overline{\Psi}_\tau(v)$ and the vector~$d\overline{\Psi}_\tau(v)(\zeta)$ are colinear in the tangent space $T_{\overline{\Psi}_\tau(v)} T\R^2$ to~$T\R^2 \setminus \{0\}$ at~$\overline{\Psi}_\tau(v)$. In the previous coordinate system, see~\eqref{eq:uuuu} and~\eqref{eq:dPsi}, the unit tangent vector~$u$ is proportional to
\[
u \sim (0,0,1,\bu_x)
\]
where $\bu_x=\bu_x(x(0,\tau,0),\tau,0)$, while
\[
d\overline{\Psi}_\tau(v)(\zeta) = \left(x_\lambda, x_\lambda \, \bu_x + \bu_\lambda, x_{s \lambda}, x_{s \lambda} \, \bu_x + x_s \, x_\lambda \, \bu_{xx} + x_s \, \bu_{x \lambda} \right)
\]
where $x_*=x_*(0,\tau,0)$ and $\bu_*=\bu_*(x(0,\tau,0),\tau,0)$. Since these two vectors are colinear, the functions~$x_\lambda$, $\bu_\lambda$ and the determinant
\[
\begin{vmatrix}
1 & x_{s \lambda} \\
\bu_x & x_{s \lambda} \, \bu_x + x_s \, x_\lambda \, \bu_{xx} + x_s \, \bu_{x \lambda}
\end{vmatrix} = x_s \, (x_\lambda \, \bu_{xx} + \bu_{x \lambda})
\]
vanish at~$(s,t,\lambda)=(0,\tau,0)$. So does the function~$\bu_{x \lambda}$ (recall that $x_s$ does not vanish). Thus, both~$\bv$ and~$\bv_x$ vanish at~$(0,\tau,0)$, where $\bv=\bu_\lambda$. That is, the function~$\bv$ has a multiple zero at~$(0,\tau,0)$. By Lemmas~\ref{lem:v} and~\ref{lem:xi}, this implies that $\zeta=0$, which is absurd. Therefore, the map $\overline{\Psi}_t:U\bar{M} \to T\R^2 \setminus \{0\}$ is transverse to the rays of~$T\R^2 \setminus \{0\}$. \\

This implies that the map $\Psi_t: U\bar{M} \to U_0\R^2$ defined from~$\overline{\Psi}_t$ by taking the radial projection $\pi:T\R^2 \setminus \{0\} \to U_0\R^2$ is a local diffeomorphism. Since the map~$\Psi_t$ is $\Z^2$-equivariant and the action of~$\Z^2$ on~$U\bar{M}$ is cocompact, it follows that the map~$\Psi_t$ is a proper local diffeomorphism. Therefore, it is a covering map. Now, the quotient map of~$\Psi_t$ is a $\pi_1$-isomorphism as observed at the end of Definition~\ref{def:Psi}. Hence, the covering~$\Psi_t:U\bar{M} \to U_0\R^2$ is a diffeomorphism and so is its quotient map.
\end{proof}

\begin{enonce}
{Observation } \label{obs:foliation-deformation}
Using the natural identifications $U_0\T^2 \simeq UM$ (by radial projection on each tangent plane) and $U^*M \simeq UM$ (by the Legendre transform), Theorem~\ref{theo:Psi} yields a family of diffeomorphisms $U^*M \to U^*M$ which, by homogeneity, extend to diffeomorphisms $T^*M \to T^*M$. We still denote by $\Psi_t:T^*M \to T^*M$ this family of diffeomorphisms. Since the asymptotic directions of the geodesics of~$M$ are preserved under the curve shortening flow, see Proposition~\ref{prop:direction}, each diffeomorphism $\Psi_t:T^*M \to T^*M$ sends the Heber foliation of~$T^*M$ to a continuous foliation of~$T^*M$ by Lipschitz Lagrangian graphs. Furthermore, if the Heber foliation is smooth, so is the image foliation.
\end{enonce}



%\forget
\begin{comment}
Now, observe that
\[
ds = F_{(x,\bu)}(1,\bu_x)
\]
That is, $x_s = F_{(x,\bu)}(1,\bu_x)$. \\

Recall that $d\Psi_0(v_0)(\zeta) = \zeta$. Since $\bu_\lambda(.,0,0)$ is equal to zero, we derive from~\eqref{eq:dPsi}
\begin{align*}
\zeta & = d\Psi_0(v_0)(\zeta) \\
& = x_\lambda \, \left(1,\bu_x,-\frac{\bu_x \bu_{xx}}{(1+\bu_x^2)^\frac{3}{2}},\frac{\bu_{xx}}{(1+\bu_x^2)^\frac{3}{2}} \right).
\end{align*}
with $s=0$, $t=0$, $\lambda=0$. Now, since $v_0$ is parallel to~$(1,\bu_x)$ and $\zeta \in T_{v_0} U \T^2$ is tangent to~$U \T^2$ at~$v_0$\dots


The curve~$\gamma_v^\tau$ defines a proper embedding of~$\R$ into~$\R^2$. This embedding extends to an embedding $h:\R \times (-\varepsilon,\varepsilon) \to \R$ onto the $\varepsilon$-tubular neighborhood of~$\gamma_v^\tau$ given by normal coordinates. Note that $h(.,0)=\gamma_v^\tau$.

In this normal coordinate system around~$\gamma_v^\tau$, every curve~$\gamma_\lambda^t$ with~$(\lambda,t)$ close enough to~$(0,\tau)$ CHK can be represented as the graph
\[
\{ (x,\bu(x,t,\lambda)) \in \R \times (-\varepsilon,\varepsilon) \mid x \in \R \}
\]
of some function~$\bu(.,t,\lambda)$ with $\bu(x,\tau,0)=0$ for every $x \in \R$. More specifically, in this normal coordinate system
\[
\gamma_\lambda^t(s) = (x(s,t,\lambda),\bu(x(s,t,\lambda),t,\lambda))
\]
where $s$ is the Finsler arclength-parameter of~$\gamma_\lambda^t$. Note that $x_s$ does not vanish.

The tangent vector of~$\gamma_\lambda^t$ at $s=0$ is given by
\begin{equation}\label{eq:tangent}
(\gamma_\lambda^t)'(0) = (x,\bu,x_s,\bu_x \, x_s)
\end{equation}
where $x_*=x_*(0,t,\lambda)$ and $\bu_* = \bu_*(x(0,t,\lambda),t,\lambda)$. Up to radial projection onto the unit Euclidean circle in the tangent plane $\Psi_t(v(\lambda))$ identifies with~$(\gamma_\lambda^t)'(0)$. Taking the derivative of~\eqref{eq:tangent} with respect to~$\lambda$, we obtain for $\lambda=0$
\[
d\Psi_t(v_0)(\zeta) = (x_\lambda, \bu_x \, x_\lambda + \bu_\lambda, x_{s \lambda}, \bu_{xx} \, x_s \, x_\lambda + \bu_{x \lambda} \, x_s + \bu_x \, x_{s \lambda})
\]
where $x_* = x_*(0,t,0)$ and $\bu_*=\bu_*(0,t,0)$. Since $x_s(0,\tau,0)=1$, we simplify this expression as
\[
d\Psi_\tau(v_0)(\zeta) = (x_\lambda, \bu_x \, x_\lambda + \bu_\lambda, x_{s \lambda}, \bu_{xx} \, x_\lambda + \bu_{x \lambda} + \bu_x \, x_{s \lambda}).
\]

Suppose $d\Psi_\tau(v_0)(\zeta) = 0$. In this case, the functions $x_\lambda$, $\bu_\lambda$, $x_{s \lambda}$ and~$\bu_{x \lambda}$ vanish at $s=0$, $t=\tau$, $\lambda=0$. In particular, the function~$\bv = \bu_\lambda$ has a multiple zero at~$(x(0,\tau,0),\tau,0)$, that is, both $\bv$ and~$\bv_\lambda$ vanish at this point.

BY~???, the function $\bu:\R \times (\tau -\delta,\tau + \delta) \times (-\eta,\eta) \to (-\varepsilon,\varepsilon)$ satisfies the following parabolic partial differential equation of the curve shortening flow
\begin{equation}\label{eq:F}
\bu_t = \FF(x,\bu,\bu_x,\bu_{xx})
\end{equation}
where $\FF$ is a smooth function defined on~$\R^4$ with
\[
\FF_q(x,u,p,q) > 0
\]
which can be expressed in terms of the coefficients of the Euclidean metric in the normal coordinate system. Differentiating~\eqref{eq:F} with respect to~$\lambda$, we obtain the following parabolic partial differential equation for $\bv=\bu_\lambda$
\[
\bv_t = a(x,t,\lambda) \, \bv_{xx} + b(x,t,\lambda) \, \bv_x + c(x,t,\lambda) \, \bv
\]
where $a=\FF_q(x,\bu,\bu_x,\bu_{xx})$, $b=\FF_p(x,\bu,\bu_x,\bu_{xx})$ and $c=\FF_u(x,\bu,\bu_x,\bu_{xx})$.
\end{comment}
%\forgotten


%\forget
\begin{comment}
We will use the following result from~\cite{An88} about the number of zeros of a parabolic partial differential equation.

\begin{theorem}
Let $u:\R\times [0,T]\to \R$ be a bounded solution of the equation
\[
u_{t}=a(x,t)u_{xx}+b(x,t)u_{x}+c(x,t)u,
\]
where $a$, $a^{-1}$, $a_{t}$, $a_{x}$, $a_{xx}$, $b$, $b_{t}$, $b_{x}$ and~$c$ are bounded functions. Then, for any $t\in[0,T]$, the number of zeros of~$u(.,t)$
\[
z(t)=\#\{x\in\R \, | \, u(x,t)=0\}
\]
is finite. In addition, if, for some $t_{0}\in[0,T]$, $u(.,t_{0})$ has a multiple zero (\ie, if $u$ and~$u_{x}$ vanish simultaneously), then $z(t)$ drops as $t$ increase beyond~$t_{0}$.
\end{theorem}
\end{comment}
%\forgotten



\section{Geodesic flow deformation} \label{sec:deformation}

Using the curve shortening flow and the family of diffeomorphisms of the previous section, we construct a deformation of the geodesic flow of a Finsler two-torus without conjugate points to the geodesic flow of the square flat two-torus.

Let $M=(\T^2,F)$ be a Finsler two-torus. Denote by $\bar{M}=(\R^2,\bar{F})$ the universal Finsler cover of $M$. Identify $U\bar{M} \simeq U_0\R^2$ by radial projection on each tangent plane. With this identification, the geodesic flow of the (quadratically convex) Finsler metric~$\bar{F}$ on~$\bar{M}$ induces a smooth free proper $\Z^2$-equivariant $\R$-action on~$U_0\R^2$ by conjugation. This action is denoted by
\[
\rho:\R \times U_0\R^2 \to U_0\R^2
\]
and is defined as
\[
\rho(s,v) = \gamma_v'(s)
\]
for every $s \in \R$ and~$v \in U_0\R^2$, where $\gamma_v$ is the arclength parametrized $\bar{F}$-geodesic induced by~$v$. In this expression, the vector~$v \in U_0\R^2$ is identified with a vector of~$U\bar{M}$ and the vector~$\gamma_v'(s) \in U\bar{M}$ is identified with a vector of~$U_0\R^2$ by radial projection. By construction, the orbits of the action~$\rho$ of~$\R$ on~$U_0\R^2$ project down to geodesics of~$\bar{M}$. 

The following result yields a deformation of the action
\[
\rho:\R \times U_0\R^2 \to U_0\R^2
\]
induced by the geodesic flow of~$\bar{M}$ into the corresponding action induced by the geodesic flow of the Euclidean plane~$\R^2$.

\begin{theo}\label{theo:convergence}
Let $M=(\T^2,F)$ be a Finsler two-torus without conjugate points and $\bar{M}=(\R^2,\bar{F})$ be its universal Finsler cover. Then there exists a smooth free proper $\Z^2$-equivariant $\R$-action
\[
\rho_t:\R \times U_0\R^2 \to U_0\R^2
\]
induced by the curve shortening flow, which starts at $\rho_{0}=\rho$, varies smoothly with respect to~$t \in [0,\infty)$, $s \in \R$ and~$v \in U_0\R^2$, and converges to the action $\rho_\infty:\R \times U_0\R^2 \to U_0\R^2$ of the geodesic flow of the Euclidean plane~$\R^2$. Here, the convergence is in the compact-open $C^k$-topology for any given~\mbox{$k \geq 0$}.

Furthermore, for every~$t \in [0,\infty]$, every $\rho_t$-orbit projects to an embedding of~$\R$ into~$\R^2$ under the canonical projection $U_0\R^2 \to \R^2$.
\end{theo}

\begin{proof}
\begin{comment}
The unit tangent bundle $U\T^{2}$ of $\T^{2}$ with the metric $F$ identifies naturally with the unit tangent bundle $U_{0}\T^{2}$ of $\T^{2}$ equipped with the canonical metric $F_{0}$ on $\T^{2}$ by radial projection on each tangent plane. Using this identification, the smooth proper free action of $\R$ on $U\T^{2}$ given by the geodesic flow of $F$ induces a smooth proper free action on $U_{0}\T^{2}$ by conjugation. This $\R$-action is denoted by
\[
\rho_{0}:\R\rightarrow \mbox{Diff}(U_{0}\T^{2})
\]
and defined as
\[
\rho_{0}(t)(v)=\gamma_{v}'(t)
\]
for every $t\in\R$ and $v\in U_{0}\T^{2}$, where the vector $v$ is identified with a vector of $UT^{2}$ by radial projection. Similarly, the vector $\gamma_{v}'(t)$ of $U\T^{2}$ is identified with a vector of $U_{0}\T^{2}$ by radial projection.
%\forgotten
\end{comment}
First, let us modify the $\rho$-action without changing its orbits through a Euclidean arclength reparametrization. For every Finsler geodesic \mbox{$\gamma:\R \to \bar{M}$}, let $\hat{\gamma}$ be the orientation-preserving arclength reparametrization of~$\gamma$ with the same initial point, \ie, $\hat{\gamma}(0)=\gamma(0)$, with respect to the Euclidean metric on~$\R^2$. Denote by $\sigma:\R \to \R$ the corresponding change of parameter with
\[
\hat{\gamma}(s) = \gamma(\sigma(s)).
\]
For every $t \in [0,1]$, define $\gamma_t:\R \to \R^2$ as
\[
\gamma_t(s) = \gamma(t \, \sigma(s) + (1-t) \, s)
\]
for every $s \in \R$. Clearly, the curve~$\gamma_t$ is a proper, regular, orientation-preserving reparametrization of~$\gamma$ with the same initial point as~$\gamma$. By construction, the isotopy~$(\gamma_t)$ connects~$\gamma_0 = \gamma$ to~$\gamma_1 = \hat{\gamma}$. Note that the isotopy~$(\gamma_t)$ smoothly depends on~$\gamma$ (where the curve space on the plane is endowed with the metrizable compact-open $C^{k}$-topology). 

This reparametrization allows us to define an $\R$-action deformation of~$\rho$
\[
\rho_t:\R \times U_0\R^2 \to U_0\R^2
\]
with~$\rho_0=\rho$ such that
\[
\rho_t(s,v) = (\gamma_v)_t'(s)
\]
for every $s \in \R$ and~$v \in U_0\R^2$. Note that the $\rho_t$-orbits remain the same for every $t \in [0,1]$ and that
\[
\rho_1(s,v) = (\widehat{\gamma_v})'(s).
\]

Next, we extend the deformation~$(\rho_t)$ of the geodesic flow to~$t \geq 1$ using the isotopy of diffeomorphisms $\Psi_t:U\bar{M} \to U_0\R^2$ given by Theorem~\ref{theo:Psi}. For every $v \in U_0\R^2$, consider the unique curve~$\gamma_u^t$ tangent to~$v$ at~$s=0$ and pointing in the same direction as~$v$. That is, $u=\Psi_t^{-1}(v)$ with the identification $U\bar{M} \simeq U_0\R^2$. Reparametrize this curve proportionally to its Euclidean arclength into~$\widehat{\gamma_u^t}$ preserving both its initial point and its orientation. Define the $\R$-action
\[
\rho_t:\R \times U_0\R^2 \to U_0\R^2
\]
such that $\rho_{t+1}(s,v)$ is the vector of~$U_0\R^2$ tangent to~$\widehat{\gamma_u^t}$ at the point of parameter~$s$ for every $t \geq 0$. That is,
\[
\rho_{t+1}(s,v) = \left(\widehat{\gamma_{\Psi_t^{-1}(v)}^t} \right)'(s)
\]
for every~$t \geq 0$, $s \in \R$ and~$v \in U_0\R^2$. Since $\Psi_t: U_0\R^2 \simeq U\bar{M} \to U_0\R^2$ is a diffeomorphism, see Theorem~\ref{theo:Psi}, the map~$\rho_t(s,.)$ is also a diffeomorphism of~$U_0\R^2$. Clearly, the $\R$-action~$\rho_t$ on~$U_0\R^2$ is smooth, free and proper. Since $\Z^2$ acts by isometries both on~$\bar{M}$ and~$\R^2$, the $\R$-action $\rho_t:\R \times U_0\R^2 \to U_0\R^2$ is $\Z^2$-equivariant. It also satisfies the symmetry property
\[
\rho_t(s,-v) = - \rho_t(-s,v)
\]
for every $t \in [0,\infty)$, $s \in \R$ and~$v \in U_0\R^2$. It follows from Lemma~\ref{lem:curv} that for every $\varepsilon>0$ and every $t\geq 0$ large enough, the $C^k$-norm of the curvature of the curves~$\widehat{\gamma_u^t}$ in the Euclidean plane~$\R^{2}$ is at most~$\varepsilon$ for every~$u\in U_{0}\R^{2}$. Hence, these curves are uniformly close to segments in any given compact set of~$\R^2$ as $t$ goes to infinity. By construction, this implies that the action $\rho_{t}$ is $C^{k}$-close to the action $\rho_{\infty}$ induced by the geodesic flow of the Euclidean plane $\R^{2}$ on any compact set for $t$ large enough. Moreover, every $\rho_t$-orbit is transverse to the fibers of $U_{0}\R^{2}\to\R^{2}$ and projects to an embedding of~$\R$ into~$\R^2$ by the canonical projection $U_0\R^2 \to \R^2$.
\end{proof}



\begin{rema}
The Heber foliation of~$T^*\bar{M}$ is deformed into the canonical Heber foliation of~$T^*\R^2$ (\ie, the foliation induced by straight lines via the Legendre transform) under the curve shortening flow of the Euclidean plane $\R^{2}$.
\end{rema}



%\forget
\begin{comment}
The quadratically convex condition (as opposed to a mere convex condition) allows us to define a \emph{geodesic flow} for~$F$ acting on the unit tangent bundle~$UM$ of~$M$, see~\cite{Be78}. The geodesic flow of a Finsler metric~$F$ on~$M$ defines a smooth proper free action of~$\R$ on~$UM$
\[
\rho_F: \R \to \Diff(UM)
\]
given by
\[
\rho_F(\theta)(v) = \gamma_v'(\theta)
\]
where $\gamma_v$ is the (arclength parametrized) $F$-geodesic induced by~$v$.
\end{comment}
%\forgotten


%\forget
\begin{comment}
\section{Minimizing geodesics on the two-torus}

Hedlund (distance to a geodesic line) + Bangert (rotational number) + Arnaud et al. (foliation of the unit tangent bundle).


Suppose $F$ is a Finsler metric on the torus $\T^{2}=\R^{2}/\Z^{2}$ or, equivalently, $F$ is a $\Z^{2}$-periodic Finsler metric on $\R^{2}$. A minimal geodesic segment is a curve $c:[a,b]\rightarrow\R^{2}$ such that
\[
d(c(x),c(y))=|x-y|
\]
for all $x,y\in[a,b]$, where $d$ denotes the distance on $\R^{2}$ induced by the Finsler metric $F$. A geodesic is a curve $c:\R\rightarrow\R^{2}$ such that, for all $x\in\R$, there exists $\varepsilon>0$ such that $c_{|[x-\varepsilon,x+\varepsilon]}$ is a minimal geodesic segment. In this paper, we are interested in geodesics such that $c_{|[a,b]}$ is a minimal geodesic segment for all compact intervals $[a,b]\subset\R$. In other words, the segment $c_{|[a,b]}$ is the shortest curve joining $c(a)$ and $c(b)$. If $F$ is euclidean then all geodesics which are straight lines are minimal. To obtain a Finsler metric on $\R^{2}/\Z^2$, we require that all translations $T_k$, $k\in\Z^2$, be isometries of $F$:
\[
\mbox{for all } (p,q)\in\R^{2}\times\R^{2} \mbox{ and all } k\in\Z^2, \mbox{ we have } F(p+k,q)=F(p,q).
\]
The $F$-length of a piecewise $C^1$-curve $\gamma:[a,b]\rightarrow\R^2$ is defined by
\[
L_{F}(\gamma):=\int_{a}^{b}F(\gamma(t),\dot\gamma(t))dt.
\]
The $F$-length of a piecewise $C^1$-curve in $\R^{2}/\Z^{2}$ is the $F$-length of any of its lifts to $\R^{2}$. Hence, one may define a distance function $\bar{d}$ on $\R^{2}/\Z^{2}$ given a symmetric Finsler metric $F$ on $\R^{2}/\Z^{2}$,
\[
\bar{d}(p,q)=\inf\{L_{F}(\gamma)| \gamma:[a,b]\rightarrow\R^{2}/\Z^{2} \mbox{ piecewise } C^{1}, \gamma(a)=p, \gamma(b)=q\}.
\]


In this section, we present some characteristics of minimal geodesics on the two-torus.


Let us start by defining the \emph{rotation number} of a minimal trajectory. We consider the space $\R^{\Z}=\{x|\mbox{ } x:\Z\rightarrow\R\}$ of bi-infinite sequences of real numbers with the product topology. An element $x\in\R^{\Z}$ will also be denoted by $(x_{i})_{i\in\Z}$. The set of minimal trajectories will be denoted by $\MM$. There exists an action $S$ of the group $\Z^{2}$ on $\R^{\Z}$ by order-preserving homeomorphisms: if $(a,b)\in\Z^{2}$ and $x\in\R^\Z$ then
\[
S_{(a,b)}x=y\quad\mbox{where}\quad y_{i}=x_{i-a}+b.
\]
This action corresponds to translating $\mbox{graph}(x)\subset\R^{2}$ by $(a,b)\in\Z^{2}$. An element $x\in\R^{\Z}$ is said to be periodic with period $(p,q)\in(\Z\times\{0\})\times\Z$ if and only if $S_{(p,q)}x=x$.\\

The following definition of a rotation number of a trajectory $x\in\R^\Z$ was given in~\cite{Ba88}

\begin{defi}
Let $x\in\R^{\Z}$ then $\tilde{\alpha}(x)$ is the rotation number of $x$ where $\tilde{\alpha}:\MM\rightarrow\R$ is a continuous map satisfying the following properties:
\begin{enumerate}
\item $\tilde{\alpha}(x)=\lim_{|i|\rightarrow +\infty}\frac{x_{i}}{i}$.
\item If $x\in\MM$ is periodic with $(q,p)$ then $\tilde{\alpha}(x)=\frac{p}{q}$.
\item $\tilde{\alpha}$ is invariant under $T$, i.e. $\tilde{\alpha}(S_{(a,b)}x)=\tilde{\alpha}(x)$ for all $(a,b)\in\Z^{2}$.
\end{enumerate}
\end{defi}

The following theorem is known as the fundamental result of Hedlund.

\begin{theo}
For every $r\in\R$, there exists a minimizing geodesic of rotation number $r$. In addition, this minimizing geodesic is at a bounded distance of an Euclidean line.
\end{theo}

\begin{prop}\label{prop:foliation}
Let $T\neq 0$ and $\alpha$ be a $T$-periodic closed geodesic on the Finsler $2$-torus without conjugate points. Then all closed geodesics which are homotopic to $\alpha$ are $T$-periodic and they define a smooth foliation of $\T^{2}$.
\end{prop}

This proposition helps to prove the following theorem which was stated in~\cite{AABZ15} and proved in the case of Riemannian metrics by Heber~\cite{Heb94},

\begin{theo}
Let $H$ be a Tonelli Hamiltonian on $T^{*}\T^{n}$. Then $H$ has no conjugate points if and only if there is a continuous foliation of $T^{*}\T^{n}$ by Lipschitz, Lagrangian, flow-invariant graphs.
\end{theo}



\section{Standard identification of a Finsler two-torus without conjugate points}

In this section, we prove that one may identify $\T^2$ with $S^{1}\times S^{1}$ so that the horizontal and vertical lines are geodesics. Thus, the nonvertical geodesic lines are graphs over the horizontal lines.

\begin{prop}
Let $\alpha,\beta:S^{1}\rightarrow \T^{2}$ be two simple closed geodesics parametrized by arclength with $\alpha(0)=\beta(0)=x_{0}$ for some $x_{0}\in\T^{2}$ such that they intersect transversely one time. \\
Then there exists a diffeomorphism $\phi:S^{1}\times S^{1}\rightarrow \T^{2}$ such that $\phi(.,0)=\alpha$ and $\phi(0,.)=\beta$.
\end{prop}

\begin{proof}

Let $\Tilde{M}$ be the universal covering of $\T^{2}$. Using Proposition~\ref{prop:foliation}, we know that there exists a foliation of $\T^{2}$, one has
\begin{enumerate}
\item for every $t\in S^{1}$, there exists a unique closed geodesic $\alpha_{t}$ of $\T^{2}$ homotopic to $\alpha$ with $\alpha_{t}(0)=\beta(t)$.
\item for every $s\in S^{1}$, there exists a unique closed geodesic $\beta_{t}$ of $\T^{2}$ homotopic to $\beta$ with $\beta_{s}(0)=\alpha(s)$.
\end{enumerate}
By Proposition~\ref{prop:foliation}, the closed geodesics $\alpha_{t}$ and $\beta_{s}$ define a vertical and a horizontal smooth foliations of $\T^{2}$, thus $\alpha_{t}$ and $\beta_{s}$ depend in a smooth way of $t$ and $s$. Hence, one may construct the map
\[
\begin{array}{ccccc}
\phi : & S^{1}\times S^{1} & \to & \T^{2} \\
& (s,t) & \mapsto & \alpha_{t}\cap\beta_{s} \\
\end{array}
\]
The map $\phi$ is well defined in the sense that the geodesics $\alpha_{t}$ and $\beta_{s}$ intersect in a single point for every $(s,t)\in S^{1}\times S^{1}$. Indeed, if there exists $t\in S^{1}$ and $s\in S^{1}$ such that $\alpha_{t}$ and $\beta_{s}$ intersect at two points, then by going to the universal covering of $\T^{2}$ and since these geodesics are homotopic, we get that these two points are conjugates, hence we get a contradiction.\\
For $s\in S^{1}$ fixed, we know that the vector field $\frac{\partial\phi}{\partial t}$ is parallel to $\beta'_{s}$ and cannot be zero, otherwise if there exists $(s,t)\in S^{1}\times S^{1}$ such that $\frac{\partial\phi}{\partial t}(s,t)=0$, hence going to the universal covering $\Tilde{M}$, we would get a Jacobi vector field such that it vanishes at two different points which contradicts the fact that there is no conjugate points. In the same way, one may prove that for $t\in S^{1}$ fixed, the vector $\frac{\partial \phi}{\partial s}$ is parallel $\alpha'_{t}$ and cannot be zero. Since both vectors are non-colinear, one may conclude that $\phi$ is a local diffeomorphism.\\
On $S^{1}\times S^{1}$ the horizontal and vertical lignes intersect at a single point, their image on $\T^{2}$ via the map $\phi$ will also intersect at a single point otherwise there will be conjugate points. Hence the degree of $\phi$ which is the number of intersections of curves on $\T^{2}$ is $1$, thus $\phi$ is a global diffeomorphism.
\end{proof}

Thus, we proved that the vertical and horizontal curves are geodesics, and if we have another curve, it cannot be tangent at a vertical line, otherwise they will be equal, hence the non-vertical lines are transverse over the horizontal lines.
\end{comment}
%\forgotten



\section{Crofton's formula for Finsler metrics without conjugate points} \label{sec:crofton}

We review a general Crofton formula for Finsler metrics without conjugate points on surfaces.


Let $M=(\T^2,F)$ be a Finsler two-torus without conjugate points. Denote by~$\bar{M}=(\R^2,\bar{F})$ the universal Finsler cover of~$M$, where the Finsler metric~$\bar{F}$ is the lift of~$F$. We will identify~$T\bar{M}$ with~$T^{*}\bar{M}$, and~$U\bar{M}$ with~$U^{*}\bar{M}$ via the Legendre transform; see~\eqref{eq:Legendre}. Using these identifications, the action~$\rho_{\bar{F}}$ of~$\R$ on~$T\bar{M}$ (resp.~$U\bar{M}$) given by the geodesic flow of~$\bar{F}$ induces an action on~$T^*\bar{M}$ (resp.~$U^{*}\bar{M}$) by conjugation by the Legendre transform, namely the cogeodesic flow of~$\bar{F}$. Both $\R$-actions will be denoted by~$\rho_{\bar{F}}$. Note that the $\R$-orbits of~$\rho_{\bar{F}}$ on~$U^{*}\bar{M}$ are transverse to the contact structure given by the kernel of the tautological one-form~$\alpha$ defined in~\eqref{eq:alpha}.


Recall that the quotient manifold theorem, see~\cite[Theorem~7.10]{lee}, asserts that if $G$ is a Lie group acting smoothly, freely and properly\footnote{A Lie group $G$ acts properly on a manifold $N$ if for every compact set $K\subset N$, the subset $\{g\in G| \mbox{ }g.K\cap K\neq\emptyset\}$ is compact.} on a smooth manifold~$N$, then the quotient space $N/G$ is a topological manifold with a unique smooth structure such that the quotient map $N \to N/G$ is a smooth submersion (actually, a locally trivial fibration by Ehresmann's fibration theorem~\cite{ehresmann}). This result applies to the $\R$-action~$\rho_{\bar{F}}$ of the cogeodesic flow of~$\bar{F}$ on~$U^* \bar{M}$. Indeed, since $\bar{M}$ is simply connected and the Finsler metric~$\bar{F}$ on~$\bar{M}$ has no conjugate points, the cogeodesic flow action on~$U^* \bar{M}$ is free and proper. Denote by
\[
\Gamma_{\bar{F}}=U^* \bar{M}/\rho_{\bar{F}}
\]
the quotient manifold and by
\[
q_{\bar{F}}:U^* \bar{M} \to \Gamma_{\bar{F}}
\]
the quotient fibration. The quotient manifold~$\Gamma_{\bar{F}}$ represents the space of unparametri\-zed oriented geodesics of the Finsler metric~$\bar{F}$ without conjugate points. It identifies with~$S^1 \times \R$, that is,
\[
\Gamma_{\bar{F}} \simeq S^1 \times \R.
\]
Indeed, recall that every geodesic of~$\bar{M}$ has an asymptotic direction and that the geodesics of~$\bar{M}$ with a given asymptotic direction foliate~$\bar{M}$; see Section~\ref{sec:foliation}. Conversely, every direction is the asymptotic direction of a geodesic of~$\bar{M}$; see Section~\ref{sec:foliation}. Thus, every oriented unparametrized geodesic of~$\bar{M}$ is determined by its asymptotic direction and its signed Euclidean distance to the origin of~$\R^2$ where the sign is positive if the orientation of the geodesic matches the orientation of the Euclidean circle centered at the origin tangent to it, and negative otherwise. This yields the desired one-to-one correspondence between~$\Gamma_{\bar{F}}$ and~$S^1 \times \R$.

The isometric action of~$\Z^2$ on~$\bar{M}$ by deck transformations induces a natural action on~$\Gamma_{\bar{F}}$. Note that the action of~$\Z^2$ on~$\Gamma_{\bar{F}}$ is neither free, nor continuous.

%\forget 
\begin{comment}
By the quotient manifold theorem, the $\R$-action on $U^{*}M$ given by the geodesic flow $\rho_{F}$ gives rise to a quotient manifold
\[
\Gamma_{F}=U^{*}M/\rho_{F}
\]
diffeomorphic to $S^{1}\times\R$???, representing the space of unparametrized oriented geodesics of the Finsler metric~$F$ without conjugate points on~$M$, and a quotient submersion
\begin{equation}\label{QuoSub}
q_{F}:U^{*}M\rightarrow\Gamma_{F}.
\end{equation}
\end{comment}
%\forgotten

By construction, the fibration $q_{\bar{F}}: U^* \bar{M} \to \Gamma_{\bar{F}}$ is $\Z^2$-equivariant. It takes a unit cotangent vector~$\xi \in U^*\bar{M}$ to the unparametrized oriented $\bar{F}$-geodesic of~$\bar{M}$ it generates. Thus, for every $\gamma \in \Gamma_{\bar{F}}$, the projection~$p(q_{\bar{F}}^{-1}(\gamma))$, where $p:T^{*}\bar{M} \rightarrow \bar{M}$ is the canonical projection, represents the unparametrized geodesic of~$\bar{F}$ on~$\bar{M}$ given by~$\gamma$. We will sometimes identify $\gamma$ and $p(q_{\bar{F}}^{-1}(\gamma))$. 

Consider the double fibration
\[
\xymatrix{ & U^* \bar{M} \ar[dl]_p \ar[dr]^{q_{\bar{F}}} \ar[r]^i & T^*\bar{M} \\
\bar{M} & & \Gamma_{\bar{F}}}
\]
where $i:U^{*}\bar{M}\hookrightarrow T^{*}\bar{M}$ is the canonical injection. Note that the product map $p\times q_{\bar{F}}:U^{*}\bar{M} \rightarrow \bar{M}\times\Gamma_{\bar{F}}$ is an embedding. Since the canonical symplectic form~$\omega$ is invariant under the cogeodesic flow, see~\eqref{eq:symplec}, then there exists a unique $\Z^2$-invariant symplectic area form~$\Omega_{\bar{F}}$ on~$\Gamma_{\bar{F}}$ such that
\begin{equation}\label{symlequ}
q_{\bar{F}}^{*} \, \Omega_{\bar{F}}=i^{*} \omega.
\end{equation}

We will need the following result about the Crofton formula on Finsler surfaces established in~\cite[Theorem~5.2]{AB}.

\begin{theo}\label{theo:crofton}
The length of every smooth curve $c$ on $\bar{M}$ with the Finsler metric~$\bar{F}$ satisfies the following equation
\begin{equation}\label{lengthequ}
\length_{\bar{F}}(c)=\frac{1}{4}\int_{\gamma\,\in\,\Gamma_{\bar{F}}}\#(\gamma\cap c)\, |\Omega_{\bar{F}}|
\end{equation}
where $|\Omega_{\bar{F}}|$ is the smooth positive $\Z^2$-invariant area density on $\Gamma_{\bar{F}}$ induced by~$\Omega_{\bar{F}}$.
\end{theo}

\begin{rema}
The Crofton formula~\eqref{lengthequ} shows that the Finsler metric~$\bar{F}$ is uniquely determined by the fibration~$q_{\bar{F}}$ (and the symplectic area form~$\Omega_{\bar{F}}$ on~$\Gamma_{\bar{F}}$ derived from~$q_{\bar{F}}$).
\end{rema}


\begin{rema}\label{rem:lambda}
We will denote by~$\lambda_{\bar{F}}$ the smooth positive $\Z^2$-invariant measure on~$\Gamma_{\bar{F}}$ corresponding to the area density~$|\Omega_{\bar{F}}|$.
\end{rema}

We derive the following corollary.

\begin{coro}\label{coro:closing}
Let $c$ be a closed geodesic of~$M$ and $\langle c \rangle$ be the subgroup of~$\pi_1(\T^2)=\Z^2$ generated by the homotopy class of~$c$. Then
\[
\length_{F}(c) = \frac{1}{4} \, \lambda_{\bar{F}}(\Gamma_{\bar{F}} / \langle c \rangle).
\]
In this relation, we still denote by~$\lambda_{\bar{F}}$ the push-forward under the quotient map $\Gamma_{\bar{F}} \to \Gamma_{\bar{F}} / \langle c\rangle$ of the restriction of $\lambda_{\bar{F}}$ to a Borel fundamental domain.
\end{coro}

\begin{proof}
Let $\bar{c}$ be a lift of~$c$ in~$\bar{M}$ (of the same length). Every generic geodesic $\gamma \in \Gamma_{\bar{F}}$ intersects the geodesic arc~$\bar{c}$ at most once. Moreover, for every generic geodesic $\gamma \in \Gamma_{\bar{F}}$, there is a unique $\langle c \rangle$-translate of~$\gamma$ intersecting the geodesic arc~$\bar{c}$. Thus, by Theorem~\ref{theo:crofton}, we obtain
\[
\length_{F}(c) = \frac{1}{4} \int_{\gamma\,\in\,\Gamma_{\bar{F}}} \#(\gamma\cap \bar{c})\, d\lambda_{\bar{F}} = \frac{1}{4} \, \lambda_{\bar{F}}(\Gamma_{\bar{F}} / \langle c \rangle).
\]
\end{proof}


This leads us to the following definition.

\begin{defi}\label{def:closing}
A smooth positive $\Z^2$-invariant measure $\lambda$ on $\Gamma=\Gamma_{\bar{F}}$ satisfies the \emph{$F$-closing condition} if
\[
\lambda(\Gamma / \alpha) = \lambda_{\bar{F}}\left(\Gamma_{\bar{F}} / \alpha\right)
\]
for every $\alpha \in \Z^2$.
\end{defi}

\begin{rema}\label{rem:closing}
The closing condition is stable under convex combinations: if we have two measures $\lambda_{1}$ and $\lambda_{2}$ on~$\Gamma$ satisfying the $F$-closing condition, then any convex combination of $\lambda_{1}$ and $\lambda_{2}$ will also satisfy the $F$-closing condition.
\end{rema}



\section{Constructing Finsler metrics with prescribed geodesics} \label{sec:prescribe}

In this section, we go over the geometric construction of Finsler metrics with prescribed geodesics on a surface given by \'Alvarez Paiva and Berck in~\cite{AB} and adapt this construction to our situation.

Let $\R^2$ be the Euclidean plane with the natural action of~$\Z^2$ by translations. Consider the double fibration of the bundle~$S^* \R^2$ of cooriented contact elements on~$\R^2$
\begin{equation}\label{eq:doublefibration}
\xymatrix{ & S^* \R^2 \simeq U_0^* \R^2 \ar[dl]_p \ar[dr]^{q} \\
\R^2 & & \Gamma}
\end{equation}
where $p:S^*\R^2 \to \R^2$ is the canonical projection and $q:S^*\R^2 \to \Gamma$ is a $\Z^2$-equivariant fibration onto an oriented surface~$\Gamma$ endowed with a $\Z^2$-action (in our case, $\Gamma=S^1 \times \R$). Identify $S^* \R^2$ with the unit cotangent bundle~$U_0^*\R^2$ of the Euclidean metric on~$\R^2$ using the canonical identification. We will assume the following:
\begin{enumerate}
\item\label{legendre} The fibration \mbox{$q:S^*\R^2 \to \Gamma$} is Legendrian (\ie, its fibers are Legendrian curves~$\gamma$ of~$S^*\R^2 \simeq U_0^*\R^2$ with respect to the tautological contact structure~$\alpha$, see~\eqref{eq:alpha}, that is, $\gamma^* \alpha=0$). 
\item \label{embedding} The product map $p \times q : S^*\R^2 \to \R^2 \times \Gamma$ is an embedding.
\end{enumerate}

By~\cite[Theorem~3.3]{AB}, for every area form~$\Omega$ on~$\Gamma$, there exists a unique Finsler metric~$\bar{F}$ on~$\R^2$ satisfying the Crofton formula
\[
\length_{\bar{F}}(c)=\frac{1}{4} \int_{\gamma\,\in\,\Gamma} \#(\gamma\cap c) \, |\Omega|
\]
for any piecewise smooth curve~$c$ on~$\R^2$. In this formula, we identify $\gamma \in \Gamma$ with the curve~$p(q^{-1}(\gamma))$ of~$\R^2$.

The Finsler metric~$\bar{F}$ is given by the Gelfand transform
\begin{equation}\label{eq:Gelfand}
\bar{F}:=p_*\left(q^*|\Omega|\right):T\R^2 \to \R
\end{equation}
of the area density~$|\Omega|$; see~\cite[Theorem~2.2]{AB}, which is defined as follows. For every $v \in T_x \R^2$, take the area density on the fiber $p^{-1}(x) = S_x^* \R^2$ by contracting~$q^* |\Omega|$ at each point~$\xi \in S_x^*\R^2$ with every vector~$\hat{v} \in T_\xi S^*\R^2$ such that $dp(\hat{v})=v$. By definition, $\bar{F}(x;v)$ is the integral of this area density over~$S_x^*\R^2$.

More precisely, for every~$x \in \R^2$, there exists a unique non-vanishing one-form~$\beta_{x}$ on~$S^{*}_{x}\R^2$ such that
\[
\bar{F}(x,v)=\int_{\xi\,\in\,S^{*}_{x}\R^2} |\xi(v)| \, \beta_{x}
\]
for every $v\in T_{x}\R^2$, where the one-form~$\beta_{x}$ defined by
\[
\left(q^* \Omega\right)_{(x,\xi)} = p^* \xi \wedge \beta_{(x,\xi)}
\]
depends smoothly on~$x$; see~\cite[Lemma~2.3]{AB}.

Alternatively, the Finsler metric~$\bar{F}$ can be defined from a smooth positive measure~$\lambda$ on~$\Gamma$ by the same formula
\[
\length_{\bar{F}}(c)=\frac{1}{4} \int_{\gamma\,\in\,\Gamma} \#(\gamma\cap c) \, d\lambda
\]
using the correspondence between smooth positive measures~$\lambda$ on~$\Gamma$ and area densities~$|\Omega|$ of area forms~$\Omega$ on~$\Gamma$.

From now on, we will assume that the smooth positive measure~$\lambda$ on~$\Gamma$ is $\Z^2$-invariant. In this case, the Finsler metric~$\bar{F}$ on~$\R^2$ passes to the quotient and induces a Finsler metric~$F$ on~$\T^2$. We will sometimes denote this metric by~$F_\lambda$ to emphasize that the construction is induced by the measure~$\lambda$ on~$\Gamma$ (and the double fibration~\eqref{eq:doublefibration}). Similarly, we denote by
\begin{equation}\label{eq:Mlambda}
M_\lambda=\left(\T^2,F_\lambda\right)
\end{equation}
the two-torus~$\T^2$ with the Finsler metric~$F_\lambda$. 


In order to apply the previous construction to the fibration induced by a smooth free proper $\Z^2$-equivariant $\R$-action
\[
\rho:\R \times U_0\R^2 \to U_0\R^2
\]
(whose orbits are transverse to the contact structure~$\ker \alpha$ induced by the tautological one-form~$\alpha$), we first modify the action to ensure that the corresponding fibration is Legendrian. To this end, consider the map
\[
\Tau:U_0 \R^2 \to U_0 \R^2
\]
sending every vector $v \in U_{0,x} \R^2$ to the unique vector $w \in U_{0,x} \R^2$ such that \mbox{$\LL(w)(v)=0$} with $(v,w)$ positively oriented. Here, $\LL$ is the Legendre transform of the Euclidean metric and the letter~$\Tau$ stands for turn. The map $\Tau:U_0\R^2 \to U_0\R^2$ and its restriction $\Tau_x:U_{0,x} \R^2 \to U_{0,x} \R^2$ are diffeomorphisms. Thus, the $\R$-action~$\rho$ on~$U_0\R^2$ induces a smooth free proper $\Z^2$-equivariant $\R$-action~$\bar{\rho}$ on~$U_0 \R^2$ by conjugation by~$\Tau$ defined as
\[
\bar{\rho}(s,u) = \Tau^{-1}\big(\rho(s,\Tau(u))\big)
\]
for every $s \in \R$ and~$u \in U_0\R^2$. This action induces a smooth free proper $\Z^2$-equivariant $\R$-action on~$U_0^* \R^2 \simeq U_0 \R^2$ by conjugation by the Legendre transform, which is still denoted by~$\bar{\rho}$.

Consider the $\Z^2$-equivariant fibration
\[
q_{\bar{\rho}}:U_0^* \R^2 \to \Gamma_{\bar{\rho}}
\]
induced by the smooth free proper $\Z^2$-equivariant $\R$-action~$\bar{\rho}$ on~$U_0^* \R^2$, where $\Gamma_{\bar{\rho}} = U_0^* \R^2 / \bar{\rho}$. Since the actions~$\rho$ and~$\bar{\rho}$ are conjugate, there is a natural diffeomorphism~$\Gamma_{\bar{\rho}} \simeq \Gamma_\rho$.

We shall need the following results about the fibration $q_{\bar{\rho}}$ corresponding to the assumptions~\eqref{legendre} and~\eqref{embedding}.

\begin{lemm}
The fibers of~$q_{\bar{\rho}}:U_0^{*}\R^{2} \rightarrow \Gamma_{\bar{\rho}}$ are Legendrian with respect to the contact structure induced by~$\alpha$ on~$U_0^{*}\R^{2}$.
\end{lemm}

\begin{proof}
For every $\gamma\in\Gamma_{\bar{\rho}}$, let $\xi\in q_{\bar{\rho}}^{-1}(\gamma)$ and $X\in T_{\xi} q_{\bar{\rho}}^{-1}(\gamma)$. Here, we identify $\gamma \in \Gamma_{\bar{\rho}}$ with the curve~$p(q_{\bar{\rho}}^{-1}(\gamma))$ of~$\R^2$. The vector $dp_{\xi}(X)$ is tangent to $\gamma$ at $p(\xi)$. By the definition of $q_{\bar{\rho}}$, the vectors tangent to~$\gamma$ at~$p(\xi)$ lie in the kernel of~$\xi$. Hence, $\xi(dp_{\xi}(X))=0$. That is, $\alpha_{\xi}(X)=0$ as desired.
\end{proof}

\begin{lemm}
The product map $\phi_{\bar{\rho}}=p \times q_{\bar{\rho}}:U^{*}_{0}\R^{2}\rightarrow\R^{2}\times\Gamma_{\bar{\rho}}$ is an embedding.
\end{lemm}

\begin{proof}
Let us start by proving that $\phi_{\bar{\rho}}$ is an immersion. Consider $\xi\in U_0^{*}\R^{2}$ and $X\in T_{\xi} q_{\bar{\rho}}^{-1}(\gamma)$ such that $d\phi_{\bar{\rho}}(\xi)(X)=0$. Thus, the vector $X$ lies in the kernel of the differential of the fibration~$q_{\bar{\rho}}$ at~$\xi$. It follows that $X$ is tangent to the $\bar{\rho}$-orbit of $U_0^{*}\R^{2}$ at~$\xi$. Since the restriction of~$p$ to each $\bar{\rho}$-orbit of~$U_0^*\R^2$ is an embedding into~$\R^2$, the relation~$dp_\xi(X)=0$ implies that~$X=0$. Thus, $d\phi_{\bar{\rho}}$ is injective and $\phi_{\bar{\rho}}$ is an immersion.

Let us prove now that $\phi_{\bar{\rho}}$ is injective. Let $\xi_{1}, \xi_{2} \in U_0^{*}\R^{2}$ such that $\phi_{\bar{\rho}}(\xi_{1})=\phi_{\bar{\rho}}(\xi_{2})$. That is, $p(\xi_{1})=p(\xi_{2})$ and $q_{\bar{\rho}}(\xi_{1})=q_{\bar{\rho}}(\xi_{2})$. Thus, the covectors~$\xi_{1}$ and~$\xi_{2}$ are based at the same point~$x$ of~$\R^{2}$. Now, the projections of the $\bar{\rho}$-orbits to $\R^{2}$ are embeddings of $\R$ into $\R^{2}$. Therefore, for every $\gamma\in\Gamma_{\bar{\rho}}$ such that $x\in p(q_{\bar{\rho}}^{-1}(\gamma))$, there exists a unique $\xi\in U^{*}_{0,x}\R^{2}$ such that $q_{\bar{\rho}}(\xi)=\gamma$. Since $q_{\bar{\rho}}(\xi_1) = q_{\bar{\rho}}(\xi_2)$, this implies that $\xi_{1}=\xi_{2}$.

Since the map~$\phi_{\bar{\rho}}$ is clearly proper, it is an embedding.
\end{proof}


Now, we can apply the construction given by~\eqref{eq:Mlambda} to the double fibration
\[
\xymatrix{ & U_0^* \R^2 \ar[dl]_p \ar[dr]^{q_{\bar{\rho}}} \\
\R^2 & & \Gamma_{\bar{\rho}} \simeq \Gamma_\rho}
\]
induced by the smooth free proper $\Z^2$-equivariant $\R$-action
\[
\rho:\R \times U_0\R^2 \to U_0\R^2.
\]
Thus, to any smooth positive $\Z^2$-invariant measure~$\lambda$ on~$\Gamma_\rho$ corresponds a Finsler metric~$F_\lambda$ on~$\T^2$. Moreover, the geodesics of the lift~$\bar{F}_\lambda$ of~$F_\lambda$ on~$\R^2$ coincide with the projection of the curves corresponding to the fibers of~$q_{\bar{\rho}}$ or~$q_{\rho}$. More precisely, the geodesics of~$\bar{F}_\lambda$ coincide with the curves $p(q_{\rho}^{-1}(\gamma))$, where $\gamma$ runs over~$\Gamma_{\rho}$. We will sometimes identify $\gamma$ with $p(q_{\rho}^{-1}(\gamma))$. Under this identification, the space of geodesics of~$\bar{F}_\lambda$ agrees with~$\Gamma_\rho$.

\begin{enonce}{Observation} \label{obs:ncp}
In practice, we will apply this construction to a double fibration where the curves~$p(q_{\rho}^{-1}(\gamma))$ intersect each other at most once (for instance, when $\Gamma$ is the space of geodesics of the universal cover~$\bar{M}$ of a Finsler two-torus~$M$ without conjugate points or the space of their images under the Euclidean curve shortening flow). In this case, the resulting Finsler metric~$F_\lambda$ has no conjugate points since the geodesics on its universal cover intersect each other at most once.
\end{enonce}



\section{Conjugate geodesic flows and the curve shortening flow} \label{sec:conjugate}


We show that, at the limit, the curve shortening flow preserves the natural measure on the space of geodesics of a Finsler metric on the two-torus with geodesic flow conjugate to the geodesic flow of a flat Finsler metric. 


Let $M_\diamond=(\T^2,F_\diamond)$ be a flat Finsler two-torus and $M=(\T^2,F)$ be a Finsler two-torus. Suppose that the geodesic flow of~$M$ is conjugate to the geodesic flow of~$M_\diamond$. This means that there exists a smooth diffeomorphism
\[
h:UM_\diamond \to UM
\]
which intertwines the geodesic flows of~$M_\diamond$ and~$M$, that is,
\[
\varphi_t \circ h = h \circ \varphi^\diamond_t
\]
where $\varphi^\diamond_t:UM_\diamond \to UM_\diamond$ and $\varphi_t:UM \to UM$ are the geodesic flows of~$M_\diamond$ and~$M$ on their unit tangent bundles.

We will need the following result proved in~\cite[\S~IV]{cro90}.

\begin{lemm}\label{lem:isom}
The map on geodesics induced by the conjugacy $h:UM_\diamond \to UM$ induces an isomorphism $h_*:\pi_1(M_\diamond) \to \pi_1(M)$.
\end{lemm}

\begin{proof}
First, note that $UM_{\diamond}$ is homeomorphic to $S^{1}\times S^{1}\times S^{1}$ and that $\pi_{1}(UM_{\diamond})$ is isomorphic to~$\Z^{3}$ with generators $a_{1}$, $a_{2}$, $a_{3}$. One may assume that $a_{1}$ and $a_{2}$ come from tangent vector fields to closed geodesics on~$M_{\diamond}$, while $a_{3}$ comes from the tangent fiber. In particular, there is a natural identification between the lattice~$\Z^{2}$ spanned by~$a_{1}$ and~$a_{2}$, and $\pi_{1}(M_{\diamond})$. This identification is given by lifting a closed geodesic to its tangent vector field in~$UM_{\diamond}$. Consider the homomorphism $(p\circ h)_{*}:\rspan\{a_{1},a_{2}\} \simeq \pi_1(M_\diamond) \to \pi_{1}(M)$, where $p:UM \to M$ is the canonical projection. This homomorphism is surjective since each element of~$\pi_{1}(M)$ can be represented by a closed geodesic~$\gamma$ of~$M$ and the inverse image of~$\gamma$ by~$h$ is a geodesic~$\gamma_\diamond$ of~$M_\diamond$, hence it lies in the span of~$a_{1}$ and~$a_{2}$. Now, every surjective homomorphism $\Z^2 \to \Z^2$ is an isomorphism. Hence the result.
\end{proof}

\begin{enonce}{Observation} \label{obs:id}
The isomorphism $h_*:\pi_1(M_\diamond) \to \pi_1(M)$, where $\pi_1(M_\diamond)=\pi_1(M)=\Z^2$, extends to an automorphism $A \in \GL_2(\Z)$ of~$\T^2$. Its inverse $A^{-1}:\T^2 \to \T^2$ is an isometry between $(A^{-1})^*F_\diamond$ and~$F_\diamond$ which induces a diffeomorphism $h_{A^{-1}}:U_{(A^{-1})^*F_\diamond} \T^2 \to U_{F_\diamond} \T^2$ between the unit tangent bundles of~$(A^{-1})^*F_\diamond$ and~$F_\diamond$. Replacing~$F_\diamond$ with~$(A^{-1})^*F_\diamond$, and $h$ with $h \circ h_{A^{-1}}$, we can assume that the map on geodesics induced by the conjugacy $h:UM_\diamond \to UM$ agrees with the identity map between~$\pi_1(M_\diamond)$ and~$\pi_1(M)$.
\end{enonce}

The following result can be extracted from~\cite[\S~IV]{cro90}.

\begin{lemm}\label{lem:ncp}
A Finsler two-torus~$M$ whose geodesic flow is conjugate to that of a flat Finsler torus $M_{\diamond}$ has no conjugate points.
\end{lemm}

\begin{proof}
We claim that every closed geodesic $\gamma$ in $M$ is the shortest in its homotopy class. To see this, let~$\tau$ be a closed geodesic homotopic to~$\gamma$. By Lemma~\ref{lem:isom}, the corresponding geodesics~$\gamma_{\diamond}$ and~$\tau_{\diamond}$ in~$M_{\diamond}$ are homotopic and hence have the same length (since $M_{\diamond}$ is a flat Finsler torus). Thus, the geodesics~$\gamma$ and~$\tau$ of~$M$ have the same length. Since this applies as well to all iterates of~$\gamma$, we see that the lift~$\bar{\gamma}$ of~$\gamma$ to the universal cover~$\bar{M}$ of~$M$ is minimizing and hence has no conjugate points. Now, observe that since~$M_{\diamond}$ is a flat Finsler torus, the closed geodesics of~$M_{\diamond}$ induce a dense subset in~$UM_{\diamond}$. The same holds on~$M$ through the conjugacy $h:UM_\diamond \to UM$. We deduce that $M$ has no conjugate points.
\end{proof}

\begin{rema}
It is unknown whether a Finsler torus without conjugate points has a geodesic flow conjugate to the geodesic flow of a flat Finsler torus.
\end{rema}

We will also need to introduce the equivariant diffeomorphism between spaces of geodesics induced by the geodesic flow conjugacy.

\begin{defi}\label{def:conjugacy}
The conjugacy $h:UM_\diamond \to UM$ between the geodesic flows of~$M_\diamond$ and~$M$ lifts to a $\Z^2$-equivariant conjugacy $\bar{h}:U\bar{M}_\diamond \to U\bar{M}$ between the lifted geodesic flows on the universal covers~$\bar{M}_\diamond$ and~$\bar{M}$. The conjugacy $\bar{h}:U\bar{M}_\diamond \to U\bar{M}$ induces a $\Z^2$-equivariant smooth diffeomorphism
\[
\bh:\Gamma_\diamond \to \Gamma
\]
between the spaces of unparametrized oriented geodesics of~$\bar{M}_\diamond$ and~$\bar{M}$.
\end{defi}

Now, assume that the conjugacy $h:UM_\diamond \to UM$ induces the identity map between~$\pi_1(M_\diamond)$ and~$\pi_1(M)$; see Observation~\ref{obs:id}. In this case, the map $\bh:\Gamma_\diamond \to \Gamma$ takes an oriented straight line to a geodesic line of the same asymptotic direction.

Recall that $\Gamma_\diamond$ identifies with the space of oriented straight lines of~$\R^2$ parametrized by~$S^1 \times \R$, that is, $\Gamma_\diamond \simeq S^1 \times \R$. Here, an oriented straight line in~$\R^2$ is determined by its asymptotic direction and its signed Euclidean distance to the origin of~$\R^2$; see Section~\ref{sec:crofton}.


Consider the deformation $\rho_t:\R \times U_0^*\R^2 \to U_0^*\R^2$ of the cogeodesic flow~$\rho_{\bar{F}}$ on~$U^*\bar{M} \simeq U_0\R^2$; see Section~\ref{sec:deformation}. Denote by
\[
\Gamma_t = U^*\R^2/\rho_t
\]
the quotient manifold and by
\[
q_t:U_0^*\R^2 \to \Gamma_t
\]
the quotient fibration; see Section~\ref{sec:crofton}. Since the $\Z^2$-translations are Euclidean isometries, the curve shortening flow induces a family of~$\Z^2$-equivariant diffeomorphisms
\begin{equation}\label{eq:ft}
f_t:\Gamma \to \Gamma_t
\end{equation}
which, at the limit, gives rise to a $\Z^2$-equivariant map
\[
f_\infty:\Gamma \to \Gamma_\diamond;
\]
see Theorems~\ref{theo:Psi} and~\ref{theo:convergence}. By construction, the map $f_\infty:\Gamma \to \Gamma_\diamond$ takes a geodesic in the universal cover~$\bar{M}$ of~$M$ to a straight line in the plane obtained at the limit by applying the curve shortening flow; see Theorem~\ref{theo:affine}. This map is measurable as a limit of continuous functions, it is continuous at irrational directions, see Propositions~\ref{prop:pcontinuous} and~\ref{prop:pinfty} below, but it is unclear whether it is continuous on~$\Gamma$; see Example~\ref{ex:noncontinuous}. In particular, it is unclear whether the function $p_\infty:S^1 \to \R$ in the following proposition is continuous at every direction~$\theta \in S^1$.

The map $f_\infty:\Gamma \to \Gamma_\diamond$ is bijective and admits a simple expression after reparametri\-zation by the map $\bh:\Gamma_\diamond \to \Gamma$ induced by the conjugacy.

\begin{prop}\label{prop:pinfty}
There exists a function $p_\infty:S^1 \to \R$ continuous at every irrational direction such that
\[
\left(f_\infty \circ \bh\right)(\theta,p) = \left(\theta,p+p_\infty(\theta)\right)
\]
for every irrational direction~$\theta \in S^1$ and every $p \in \R$.
\end{prop}

\begin{proof}
Let $\Delta=(\theta,p) \in \Gamma_{\diamond}$ be a straight line in~$\R^2$. Its image~$f_\infty \circ \bh(\Delta)$ is the limit of the curve~$\bh(\Delta)$ under the curve shortening flow. Since the map $\bh:\Gamma_\diamond \to \Gamma$ and the curve shortening flow preserve the asymptotic direction of a curve, the image~$f_\infty \circ \bh(\Delta)$ can be represented by~$(\theta,p') \in \Gamma_\diamond$. For $\Delta_0=(\theta,0)$, we have $f_\infty \circ \bh(\Delta_0)=(\theta,p_\infty(\theta))$, where the function \mbox{$p_\infty:S^1 \to \R$} is continuous at every irrational direction; see Proposition~\ref{prop:pcontinuous}.


The signed Euclidean distance~$d_{\pm}(\gamma \cdot \Delta, \Delta)$ between the parallel oriented lines~$\gamma \cdot \Delta$ and~$\gamma$ depends only on~$\gamma$ and~$\theta$, but not on~$p$. Thus,
\[
d_{\pm}\big(\gamma \cdot (f_\infty \circ \bh(\Delta)),f_\infty \circ \bh(\Delta)\big) = d_{\pm}(\gamma \cdot \Delta,\Delta).
\]
By $\Z^2$-equivariance of $f_\infty \circ \bh$, we obtain
\[
f_\infty \circ \bh(\gamma \cdot \Delta_0) = \gamma \cdot (f_\infty \circ \bh(\Delta_0)) = \big(\theta,p_\infty(\theta)+d_{\pm}(\gamma \cdot \Delta_0,\Delta_0)\big).
\]

Observe that
\[
\gamma \cdot \Delta_0 = \big(\theta,d_{\pm}(\gamma \cdot \Delta_0,\Delta_0)\big)
\]
under the identification~$\Gamma_\diamond \simeq S^1 \times \R$.

Now, if $\theta$ is irrational, the subset
\[
\left\{ \gamma \cdot \Delta_0 = \left(\theta,d_{\pm}(\gamma \cdot \Delta_0,\Delta_0)\right) \,\middle|\, \gamma \in \Z^2 \right\}
\]
is dense in~$\{ (\theta,p) \in \Gamma_\diamond \mid p \in \R \}$. By continuity of~$f_\infty \circ \bh$ at irrational directions, it follows that
\[
f_\infty \circ \bh(\theta,p) = \left(\theta,p+p_\infty(\theta)\right)
\]
for every $(\theta,p) \in \Gamma_\diamond \simeq S^1 \times \R$ of irrational direction.
\end{proof}

We deduce the following measure invariance property.

\begin{prop}\label{prop:measure}
The $\Z^2$-equivariant map $f_\infty \circ \bh: \Gamma_\diamond \to \Gamma_\diamond$ preserves the measure~$\lambda_\diamond$, that is,
\[
(f_\infty \circ \bh)_* \lambda_\diamond = \lambda_\diamond.
\]
\end{prop}

\begin{proof}
It is enough to show that the measures $\lambda_\diamond$ and $(f_\infty \circ \bh)_* \lambda_\diamond$ agree on the rectangles $[\theta_1,\theta_2] \times [p_1,p_2]$ of~$\Gamma_\diamond \simeq S^1 \times \R$. Since the flat Finsler metric~$F_\diamond$ is invariant by the translations of~$\R^2$, the measure~$\lambda_\diamond$ associated to~$F_\diamond$ is invariant by the action on~$\Gamma_\diamond$ induced by these translations. It follows that the smooth measure~$\lambda_\diamond$ on~$\Gamma_\diamond \simeq S^1 \times \R$ can be written as
\[
\lambda_\diamond = K(\theta) \, |d\theta \wedge dp|
\]
where $K:S^1 \to \R$ is a smooth function which does not depend on~$p$. By Proposition~\ref{prop:pinfty}, the map $f_\infty \circ \bh: \Gamma_\diamond \to \Gamma_\diamond$ agrees with the map $(\theta,p) \mapsto (\theta,p+p_\infty(\theta))$ almost everywhere on~$\Gamma_\diamond$, where $p_\infty:S^1 \to \R$ is continuous at every irrational direction. Thus,
\begin{align*}
\left[(f_\infty \circ \bh)_* \lambda_\diamond\right]\left([\theta_1,\theta_2] \times [p_1,p_2]\right) & = \int_{\theta_1}^{\theta_2} \int_{p_1-p_\infty(\theta)}^{p_2-p_\infty(\theta)} K(\theta) \, dp \, d\theta \\
& = \int_{\theta_1}^{\theta_2} \int_{p_1}^{p_2} K(\theta) \, dp \, d\theta \\
& = \lambda_\diamond\left([\theta_1,\theta_2] \times [p_1,p_2]\right).
\end{align*}
Hence the desired result.
\end{proof}


%\forget
\begin{comment}
Since it is unclear whether the function~$p_\infty:S^1 \to \R$ is continuous, we will proceed with an approximation argument. By convolution with a mollifier, there exists a continuous family of smooth functions $p_t:S^1 \to \R$ with $t \in [0,\infty)$ starting at~$p_0=0$ and converging to~$p_\infty$ in~$L^1(S^1)$. Consider the family of maps
\[
\widehat{f}_t:\Gamma \to \Gamma_t
\]
defined by
\[
\widehat{f}_t = f_t \circ (\bh \circ P_t \circ \bh^{-1})
\]
for $t \in [0,\infty]$, where
\[
\begin{array}{rccc}
P_t: & \Gamma_0 & \longrightarrow & \Gamma_0 \\
& (\theta,p) & \longmapsto & (\theta,p-p_t(\theta))
\end{array}
\]
This family of maps starts at~$\widehat{f}_0=f_0=\id_{\Gamma_0}$ and ends at~$\widehat{f}_\infty$, where the composite
\[
\widehat{f}_\infty \circ \bh:\Gamma_0 \to \Gamma_0
\]
is the identity map. Clearly, the maps~$\widehat{f}_t: \Gamma \to \Gamma_t$ are measurable and $\Z_2$-equivariant for $t \in [0,\infty]$. Furthermore, they are smooth for $t \in [0,\infty)$.

Now, the push-forward measures $(\widehat{f}_t \circ \bh)_* \lambda_\star$ on~$\Gamma_t$ are smooth and~$\Z^2$-invariant, even for $t=\infty$ where $(\widehat{f}_\infty \circ \bh)_* \lambda_\star = \lambda_\star$ on~$\Gamma_\infty = \Gamma_\star$. These measures induce a family of Finsler metrics~$(F_t)$ on~$\T^2$ starting at~$F'$ and ending at~$F_\star$; see Section~\ref{sec:prescribe}.

We need to check that $F_t$ converges to~$F_\star$ with respect to the uniform topology. That is, the distance they induce on~$\T^2$ are uniformly close.
\end{comment}
%\forgotten



%\forget
\begin{comment}
\section{Smooth geodesic foliations and conjugate geodesic flows} \label{sec:smoothfoliation}


We shall need the following result, first stated (but not used) in~\cite{CK95} without proof. Note that the converse is obvious.

\begin{theo}\label{theo:conjugacy}
If the geodesic foliation of $UM$ is smooth, then the geodesic flow of $(M,F_{t})$ is $C^1$-conjugate to the geodesic flow of a flat Finsler metric.
\end{theo}

\begin{proof}
We know that the curve shortening flow induces a one-parameter family of~$\Z^2$-equivariant diffeomorphisms $f_t:\Gamma \to \Gamma_t$ converging to a $\Z^2$-equivariant map $f_\infty:\Gamma \to \Gamma_{\diamond}$
\[
\xymatrix{
U_{F}^{*}\R^{2}\ar[d] \ar[r] & U_{F_{t}}^{*}\R^{2\ar[d]}\ar[d] \ar[r] & U_{F_{\diamond}}^{*}\R^{2}\ar[d]\\
\Gamma \ar[r] & \Gamma_{t} \ar[r] & \Gamma_{\diamond}}
\]
where $f_\infty$
takes a geodesic in the universal cover of~$M$ to a straight line in the plane obtained at the limit by applying the curve shortening flow. Recall that $U_{F_{\diamond}}^{*}\R^{2}$ and $U_{F_{t}}^{*}\R^{2}$ can be identified with $\R^{2}\times S^{1}$, hence for the flat metric, each element of $U_{F_{\diamond}}^{*}\R^{2}$ can be identified with $(p,s,\theta)$, where $p$ represents the distance from the origin, $s$ represents the length of the segment and $\theta$ is the angle with respect to the horizontal axis. In addition, each element of $U_{F_{t}}\R^{2}$ can be identified with $(p,B_{\theta}(x),\theta)$, where $B_{\theta}$ is the



In addition, the length spectrum of $(M,F)$ and $(M,F_{t})$ is preserved. Indeed, Let $\gamma$ be a closed geodesic on $M$ and $\bar{\gamma}=p(\gamma)$ where $p:U\R^{2}\rightarrow \R^{2}$ is the canonical projection. One has that every curve $c\in\Gamma$ intersects $\bar{\gamma}$ at most 1 time. By translation, if we take all the orbits of $c$ by $<\gamma>$, then, only one curve in these orbits will intersect $\bar{\gamma}$. Thus, using the closing condition, one has
\[
L_{F}(\gamma)=\int_{c\in\Gamma}\#(\bar{\gamma}\cap c) d\mu=\int_{\Gamma/<\gamma>}d\mu=\mu(\Gamma/<\gamma>)= L_{F_{t}}(\gamma).
\]
\end{proof}
\end{comment}
%\forgotten



\section{Smooth Heber foliations and conjugate geodesic flows} \label{sec:smoothfoliation}

We show that two Finsler two-tori without conjugate points having the same marked length spectrum and smooth Heber foliations have the same dynamics, namely their geodesic flows are conjugate. Note that this result has already been stated (without proof) in the introduction of~\cite{CK95}, at least for Riemannian metrics. 

We introduce the definition of the stable norm in the special case of a Finsler two-torus and refer to~\cite{gro99} for a general discussion.

\begin{defi}
Let $M=(\T^2,F)$ be a Finsler two-torus. Denote by~$\bar{M}=(\R^2,\bar{F})$ its universal Finsler cover. The \emph{stable norm} of~$M$ is a norm on~$H_1(\T^2;\R)$ defined as follows. For every $\gamma \in H_1(\T^2;\Z)$,
\[
|\gamma|_{\rst} = \lim_{k\,\to\,+\infty} \frac{d_{\bar{M}}\left(o,\gamma^k \cdot o\right)}{k}
\]
where $o \in \bar{M}$ is a fixed origin. Note that the limit exists and does not depend on the origin~$o$. The function~$|.|_{\rst}$ defined by this function extends to a vector-space norm on~$H_1(\T^2;\R)$, the so-called stable norm of~$M$, still denoted by~$|.|_{\rst}$. Alternatively, the stable norm of an integral homology class can be defined as the minimal length of a cycle representing this integral homology class; see~\cite[Lemma~4.32]{gro99}.
\end{defi}


Let $M=(\T^2,F)$ be a Finsler two-torus without conjugate points. We will endow~$H_1(\T^2;\R)$ with the stable norm~$|.|_{\rst}$ of~$M$. Denote by~$\B \subseteq H_1(\T^2;\R)$ the unit ball of the stable norm of~$M$. The space~$S^1$ of asymptotic directions of~$\bar{M}$ identifies with the boundary~$\partial \B$ of~$\B$. Moreover, we have the following regularity result obtained in~\cite[Lemma~5]{MS11} using Aubry-Mather theory and weak KAM theory.

\begin{lemm}
The stable norm~$|.|_{\rst}$ of a Finsler two-torus~$M$ without conjugate points is~$C^1$, except possibly at~$0$. In particular, the boundary~$\partial \B$ of~$\B$ admits a unique tangent line at every point.
\end{lemm}

\begin{proof}
Since $M$ has no conjugate points, the cotangent space~$T^*M$ admits a Heber foliation; see Theorem~\ref{theo:foliation}. By~\cite[Lemma~5]{MS11}, the existence of a Heber foliation implies that Mather's average action~$\beta_F:H_1(\T^2;\R) \to \R$ is~$C^1$. We do not need the general definition of~$\beta_F$, simply that in our case $\beta_F = \frac{1}{2} |. |_{\rst}^2$; see~\cite[Proposition~3.3]{BM08}. This implies that the stable norm~$|.|_{\rst}$ of~$M$ is~$C^1$, except possibly at~$0$.
\end{proof}

For the rest of this section, we will assume that the Heber foliation is smooth; see Definition~\ref{def:heber}. Thus, the unit cotangent bundle~$U^*\bar{M}$ of~$\bar{M}$ identifies with~$\R^2 \times S^1$ through the smooth diffeomorphism
\begin{align}\label{eq:identification1}
\begin{split}
U^*\bar{M} & \rightarrow \R^2 \times S^1 \\
(x,\xi) & \mapsto (x,\theta)
\end{split}
\end{align}
where $\theta$ is the asymptotic direction of the geodesic induced by~$\xi$; see~\eqref{eq:UM}. The inverse map of this diffeomorphism is given by~$\xi=-dB_\theta(x)$; see~\eqref{eq:v}.


%\forget
\begin{comment}
The unit cotangent bundle~$U^*\bar{M}$ of~$\bar{M}$ identifies with~$\R^2 \times S^1$ as follows
\begin{align}\label{eq:identification1}
\begin{split}
U^*\bar{M} & \rightarrow \R^2 \times S^1 \\
(x,\xi) & \mapsto (x,\theta)
\end{split}
\end{align}
where $\theta$ is the asymptotic direction (of the projection to~$\R^2$) of the cogeodesic induced by~$\xi$; see Definition~\ref{def:asymptotic}. By assumption, this map is a diffeomorphism. Using this identification, the Busemann function based at the origin~$o \in \R^2$ and pointing in the direction of~$\theta$ is defined as
\[
B_\theta(x) = \lim_{t \to +\infty} d(x,c_\theta(t)) -t
\]
where $c_\theta$ is the geodesic ray arising from~$o$ with asymptotic direction~$\theta$. Note that $c_\theta$ is the projection of the cogeodesic induced by~$\xi$ to~$\R^2$. Via the Legendre transform, the differential~$dB_\theta$ at~$x$ corresponds to the unit tangent vector at~$x$ generating a geodesic ray with asymptotic direction~$\theta$. Thus, the inverse map of~\eqref{eq:identification1} is given by~$\xi = dB_\theta(x)$. By assumption, both the map~\eqref{eq:identification1} and its inverse are smooth. We refer to Gromov99 and Heber for further details on Busemann functions and their regularity.
\end{comment}
%\forgotten



Consider the $\Z^2$-invariant symplectic area form~$\Omega$ on the space~$\Gamma$ of unparametrized oriented geodesics on~$\bar{M}$; see~\eqref{symlequ}. Define $\Gamma \to S^1$ as the quotient of~$U^* \bar{M} \to \R^2 \times S^1 \to S^1$ under the cogeodesic flow, where the first map $U^* \bar{M} \to \R^2 \times S^1$ is the diffeomorphism~\eqref{eq:identification1}. Let $\eta$ be the pullback of the canonical one-form of~$S^1$ under the smooth map $\Gamma \to S^1$ taking a geodesic of~$\bar{M}$ to its asymptotic direction~$\theta$. Observe that the one-form~$\eta$ does not vanish on~$\Gamma$ and is $\Z^2$-invariant (by $\Z^2$-invariance of $\Gamma \to S^1$). Thus, there exists a $\Z^2$-invariant differential one-form~$\zeta$ on~$\Gamma$ such that
\begin{equation}\label{eq:etabeta}
\Omega = \eta \wedge \zeta.
\end{equation}
The one-form~$\zeta$ is not uniquely defined. Indeed, every other $\Z^2$-invariant one-form~$\zeta'$ on~$\Gamma$ with $\Omega = \eta \wedge \zeta'$ satisfies
\[
\zeta' = \zeta + h \, \eta
\]
where $h:\Gamma \to \R$ is a smooth $\Z^2$-invariant function on~$\Gamma$. Still, we can think of~$\zeta$ as ``the quotient of the two-form~$\Omega$ by the one-form~$\eta$''. We will write~$\Omega/\eta$ for the class of $\Z^2$-invariant differential one-forms~$\zeta$ on~$\Gamma$ with $\Omega = \eta \wedge \zeta$. For a curve~$c$ of~$\Gamma$ formed of geodesics of~$\bar{M}$ with the same asymptotic direction, we will write
\begin{equation}\label{eq:O/}
\int_c \Omega/\eta := \int_c \zeta
\end{equation}
where $\zeta$ is a representative of~$\Omega/\eta$. Observe that the one-form~$\eta$ vanishes along~$c$. Thus, even though~$\zeta$ is not uniquely defined, the value of the integral~\eqref{eq:O/} does not depend on the one-form~$\zeta$ representing~$\Omega/\eta$, which justifies the notation.


The unit cotangent bundle~$U^*\bar{M}$ also identifies with~\mbox{$\Gamma \times \R = S^1 \times \R \times \R$.} Loosely speaking, a unit covector $(x,\xi) \in U^*\bar{M}$ is represented by the geodesic in~$\Gamma$ (of asymptotic direction~$\theta$) it induces and a parameter~$s$ measuring the signed distance to the origin of the geodesic (defined as the unique point of the geodesic where the Busemann function~$B_\theta$ vanishes); see Figure~\ref{fig:tps}.

\begin{figure}[htbp!]
%\newcommand*\svgwidth{7cm}
\centering
\includegraphics[scale=1.7]{dessin-2.pdf}
%\end{center}
\caption{$(\theta,p,s)$ coordinates}\label{fig:tps}
\end{figure}

More precisely, this identification is given by
\begin{align}\label{eq:identification2}
\begin{split}
U^*\bar{M} & \rightarrow S^1 \times \R \times \R \\
(x,\xi) & \mapsto (\theta,p,s)
\end{split}
\end{align}
In this expression, $\theta$ is the asymptotic direction of the geodesic induced by~$\xi$, \mbox{$s=B_\theta(x)$} and
\[
p = \int_c \Omega/\eta
\]
where $c$ is the path of~$\Gamma$ formed of the geodesics of~$\bar{M}$ (with the same asymptotic direction~$\theta$) lying between the geodesic of asymptotic direction~$\theta$ passing through the origin~$o \in \R^2$ and the geodesic of asymptotic direction~$\theta$ passing through~$x$. Since the Heber foliation is assumed to be smooth, the map~\eqref{eq:identification2} is also a smooth diffeomorphism.

%\forget
\begin{comment}
\[
p = \int_{\bar{c}} \iota_{X_\theta} \omega
\]
is the integral of the one-form~$\iota_{X_\theta} \omega$ along~$\bar{c}$, where $X_\theta$ is the vector field on~$U^*\bar{M}$ corresponding to~$(x,\theta,0,1) \in T(\R^2 \times S^1)$ through the identification~\eqref{eq:identification1}, the differential two-form~$\omega$ is the restriction to~$U^*\bar{M}$ of the canonical symplectic form on~$T^*\bar{M}$, see~\eqref{eq:symplec}, and $\bar{c} = \LL(c')$ is the natural lift to~$U^*\bar{M}$ (via the Legendre map~$\LL$) of any path~$c$ of~$\bar{M}$ joining the origin~$o \in \R^2$ to~$x$. Here, $\iota_{X_\theta}$ is the operator that contracts a differential form with the vector field~$X_\theta$. It follows from the assumption on the smoothness of the geodesic flow that the map~\eqref{eq:identification2} is also a diffeomorphism.
\end{comment}
%\forgotten


Under this identification, the symplectic area form~$\Omega$ on~$\Gamma$ can be written as
\[
\Omega = \eta \wedge dp
\]
(but we won't need this) and the action of the cogeodesic flow of~$\bar{M}$ on~$U^*\bar{M}$ is given by
\begin{equation}\label{eq:s+t}
\varphi_t(\theta,p,s) = (\theta,p,s+t).
\end{equation}
Finally, the action of~$\Z^2$ on~$U^* \bar{M}$ takes the following form. There exist two additive functions $\vec{p}_\theta:\Z^2 \to \R$ and $\vec{s}_\theta:\Z^2 \to \R$ such that for every $\gamma \in \Z^2$, we have
\begin{equation}\label{eq:action}
\gamma \cdot (\theta,p,s) = \left(\theta,p+\vec{p}_\theta(\gamma),s+\vec{s}_\theta(\gamma)\right).
\end{equation}
The functions $\vec{p}_\theta:\Z^2 \to \R$ and $\vec{s}_\theta:\Z^2 \to \R$ can be defined as
\begin{align}
\vec{p}_\theta(\gamma) & = \int_{(\theta,x)}^{(\theta,\gamma \cdot x)} \Omega/\eta \label{eq:vp}
\\
\vec{s}_\theta(\gamma) & = B_\theta(\gamma \cdot x) - B_\theta(x) \label{eq:vs}
\end{align}
where the integral is along the path of~$\Gamma$ formed of the geodesics of~$\bar{M}$ (with the same asymptotic direction~$\theta$) lying between the geodesic of asymptotic direction~$\theta$ passing through~$x$ and the geodesic of asymptotic direction~$\theta$ passing through~$\gamma \cdot x$, where $x$ is any point of~$\bar{M}$.

It remains to show the following.

\begin{lemm}\leavevmode
\begin{enumerate}
\item\label{lemm12.3.1} The functions~$\vec{p}_\theta$ and~$\vec{s}_\theta$ given by the formulas~\eqref{eq:vp} and~\eqref{eq:vs} do not depend on $x\in\bar{M}$.
\item\label{lemm12.3.2} The functions~$\vec{p}_\theta$ and~$\vec{s}_\theta$ are additive.
\item\label{lemm12.3.3} The relation~\eqref{eq:action} for the action of~$\gamma \in \Z^2$ is satisfied.
\end{enumerate}
\end{lemm}

\begin{proof}
It follows from the $\Z^2$-invariance of the one-form~$\zeta$ representing~$\Omega/\eta$ that the formula~\eqref{eq:vp} defining $\vec{p}_{\theta}$ does not depend on~$x \in \bar{M}$. Indeed, for every $x,x'\in \bar{M}$, we have
\begin{align*}
\vec{p}_\theta(\gamma)&=\int_{(\theta,x)}^{(\theta,x')}\Omega/\eta +\int_{(\theta,x')}^{(\theta,\gamma\cdot x')}\Omega/\eta+\int_{(\theta,\gamma\cdot x')}^{(\theta,\gamma\cdot x)}\Omega/\eta\\
&=\int_{(\theta,x)}^{(\theta,x')}\Omega/\eta +\int_{(\theta,x')}^{(\theta,\gamma\cdot x')}\Omega/\eta+\int_{(\theta, x')}^{(\theta, x)}\Omega/\eta\\
&=\int_{(\theta,x')}^{(\theta,\gamma\cdot x')}\Omega/\eta.
\end{align*}

Similarly, it directly follows from~\cite[Corollary~2.6]{CS86}, using the fact that the (minimal) closed geodesics in every free homotopy class of~$M$ cover the torus, that for every rational direction~\mbox{$\theta \in S^1$} and every~$\gamma \in \Z^2$, the difference $B_{\theta}(\gamma\cdot x)-B_{\theta}(x)$ does not depend on $x \in \bar{M}$. By a continuity argument, see Proposition~\ref{prop:heber}, the same holds true for every~$\theta\in S^{1}$ and $\gamma\in\Z^{2}$.

Since the functions~$\vec{p}_\theta$ and~$\vec{s}_\theta$ do not depend on $x\in\bar{M}$, this ensures that both functions are additive and that the relation~\eqref{eq:action} for the action of~$\gamma \in \Z^2$ is satisfied.
\end{proof}


%\forget
\begin{comment}
It follows from the $\Z^2$-invariance of the one-form~$\zeta$ representing~$\Omega/\eta$ that the formula~\eqref{eq:vp} defining $\vec{p}_{\theta}$ does not depend on~$x \in \bar{M}$. Indeed, for every $x,x'\in \bar{M}$, we have
\begin{align*}
\vec{p}_\theta(\gamma)&=\int_{(\theta,x)}^{(\theta,x')}\Omega/\eta +\int_{(\theta,x')}^{(\theta,\gamma\cdot x')}\Omega/\eta+\int_{(\theta,\gamma\cdot x')}^{(\theta,\gamma\cdot x)}\Omega/\eta\\
&=\int_{(\theta,x)}^{(\theta,x')}\Omega/\eta +\int_{(\theta,x')}^{(\theta,\gamma\cdot x')}\Omega/\eta+\int_{(\theta, x')}^{(\theta, x)}\Omega/\eta\\
&=\int_{(\theta,x')}^{\theta,\gamma\cdot x')}\Omega/\eta.
\end{align*}
Similarly, it follows from the $\Z^{2}$-invariance of the differential~$dB_{\theta}$, \cite[Corollary~2.6]{CS86}, the definition of $B_{\theta}$ for both rational and irrational directions, and a continuity argument that for every $\theta\in S^{1}$ and $\gamma\in\Z^{2}$, the difference $B_{\theta}(\gamma\cdot x)-B_{\theta}(x)$ does not depend on $x$. Since $\vec{p}_\theta$ and~$\vec{s}_\theta$ do not depend on $x\in\bar{M}$, this also ensures that these functions are additive and that the relation~\eqref{eq:action} for the action of~$\gamma$ is satisfied.
\end{comment}
%\forgotten



The additive functions $\vec{p}_\theta:H_1(\T^2;\Z) \to \R$ and $\vec{s}_\theta:H_1(\T^2;\Z) \to \R$ extend to linear forms $\vec{p}_\theta:H_1(\T^2;\R) \to \R$ and $\vec{s}_\theta:H_1(\T^2;\R) \to \R$.

\begin{lemm}\label{lem:linear}
Let $\theta \in \partial \B \simeq S^1$.
\begin{enumerate}
\item\label{linear1} The linear form $\vec{p}_\theta:H_1(\T^2;\R) \to \R$ satisfies
\begin{equation*}\label{eq:kervp}
\ker \vec{p}_\theta = \R \, \theta.
\end{equation*}
\item\label{linear2} The linear form $\vec{s}_\theta:H_1(\T^2;\R) \to \R$ is of norm~$1$ (with respect to the stable norm of~$M$) and satisfies $\vec{s}_\theta(-\theta)=1$. In particular, the kernel of~$\vec{s}_\theta$ is parallel to the tangent line to~$\partial \B$ at~$-\theta$. 
\end{enumerate}
\end{lemm}

\begin{proof}\ \\*[0.20em]
\noindent\eqref{linear1} Let $\theta \in \partial \B$ be a rational direction and~$\gamma \in H_1(\T^2;\Z)$ be an integral vector pointing in the direction of~$\theta$. The geodesic~$(\theta,x)$ of~$\bar{M}$ of asymptotic direction~$\theta$ passing through~$x$ coincides with its image~$(\theta,\gamma \cdot x)$ under the translation of~$\gamma$. Thus, $\vec{p}_\theta(\gamma)=0$. Since $\theta = \frac{\gamma}{|\gamma|_{\rst}}$, it follows that $\vec{p}_\theta(\theta)=0$. The same relation also holds true if $\theta$ is an irrational direction by a continuity argument. Hence, $\ker \vec{p}_\theta = \R \, \theta$.

\noindent\eqref{linear2} As previously, let $\theta \in \partial \B$ be a rational direction and~$\gamma \in H_1(\T^2;\Z)$ be an integral vector pointing in the direction of~$\theta$. By Proposition~\ref{prop:Hfoliation}.\eqref{geod1}, the displacement function $d_\gamma : \bar{M} \to \R$ of~$\gamma \in H_1(\T^2;\Z)$, defined as $d_\gamma(x)=d_{\bar{M}}(x,\gamma \cdot x)$, is constant equal to~$|\gamma|_{\rst}$. Since the Busemann function~$B_\theta:\bar{M} \to \R$ is $1$-Lipschitz, we obtain
\[
|\vec{s}_\theta(\gamma)| = |B_\theta(\gamma \cdot x) - B_\theta(x)| \leq d_{\bar{M}}(x,\gamma \cdot x) = |\gamma|_{\rst}.
\]
Thus, the linear form~$\vec{s}_\theta:H_1(M;\R) \to \R$ is of norm at most~$1$.

Furthermore, since $\gamma$ points in the direction of~$\theta$, we have
\[
\vec{s}_\theta(\gamma) = B_\theta(\gamma \cdot x) - B_\theta(x) = - d_{\bar{M}}(x,\gamma \cdot x) = - |\gamma|_{\rst}.
\]
Since $\theta = \frac{\gamma}{|\gamma|_{\rst}}$, we obtain $\vec{s}_\theta(\theta) = -1$. The same relation also holds true if $\theta$ is an irrational direction by a continuity argument.

Therefore, $\lVert \vec{s}_\theta \rVert = 1$ and $\vec{s}_\theta(-\theta) = 1$ for every~$\theta \in \partial \B$.
\end{proof}


This construction allows us to prove the following result.

\begin{theo}\label{theo:conjugacy}
Let $M_1=(\T^2,F_{1})$ and $M_2=(\T^2,F_{2})$ be two Finsler two-torus without conjugate points whose Heber foliation are smooth. Suppose that the two metrics have the same marked length spectrum. Then the geodesic flows of~$M_1$ and~$M_2$ are conjugate.
\end{theo}

\begin{proof}
Denote by $\vec{p}_{\theta,i}:H_1(\T^2;\R) \to \R$ and $\vec{s}_{\theta,i}:H_1(\T^2;\R) \to \R$ the linear forms corresponding to the metric~$F_i$ defined in~\eqref{eq:vp} and~\eqref{eq:vs}. By Lemma~\ref{lem:linear}$\MK$\eqref{linear1}, the linear forms~$\vec{p}_{\theta,1}$ and~$\vec{p}_{\theta,2}$ have the same kernel and are therefore proportional. That is,
\[
\vec{p}_{\theta,2} = \kappa(\theta) \, \vec{p}_{\theta,1}
\]
for some smooth function~$\kappa:S^1 \to \R$. Since the metrics $F_1$ and $F_2$ have the same length spectrum, they also have the same stable norm. It follows from Lemma~\ref{lem:linear}$\MK$\eqref{linear2} that the linear forms~$\vec{s}_{\theta,1}$ and~$\vec{s}_{\theta,2}$ coincide for every direction~$\theta \in S^1$.

Using the identification~\eqref{eq:identification2}, consider the map
\begin{align*}
U^*\bar{M}_1 & \rightarrow U^*\bar{M}_2 \\
(\theta,p,s) & \mapsto \left(\theta,\kappa(\theta) \, p,s \right).
\end{align*}
By the action of the cogeodesic flow in this coordinate system, see~\eqref{eq:s+t}, this map $U^*\bar{M}_1 \to U^*\bar{M}_2$ conjugates the cogeodesic flows on~$\bar{M}_1$ and~$\bar{M}_2$. Furthermore, the map $U^*\bar{M}_1 \to U^*\bar{M}_2$ passes to the quotient under the $\Z^2$-action on~$\bar{M}_1$ and~$\bar{M}_2$; see~\eqref{eq:s+t}. Thus, the quotient map $U^*M_1 \to U^*M_2$ conjugates the cogeodesic flows of~$M_1$ and~$M_2$, and therefore the geodesic flows of~$M_1$ and~$M_2$ through their Legendre transforms, as required.
\end{proof}



%\forget
\begin{comment}
Under this identification, the symplectic area form~$\Omega$ on~$\Gamma$ can be written
\[
\Omega = \eta \wedge dp
\]
(but we won't need this) and the action of the cogeodesic flow of~$\bar{M}$ on~$U^*\bar{M}$ is given by
\[
\varphi_t(\theta,p,s) = (\theta,p,s+t).
\]
The action of~$\Z^2$ on~$U^* \bar{M}$ takes the following form.

\begin{proposition}
There exist two additive functions $\vec{p}_\theta:\Z^2 \to \R$ and $\vec{s}_\theta:\Z^2 \to \R$ such that for every $\gamma \in \Z^2$, we have
\begin{equation}\label{eq:action}
\gamma \cdot (\theta,p,s) = (\theta,p+\vec{p}_\theta(\gamma),s+\vec{s}_\theta(\gamma)).
\end{equation}

Furthermore, the functions $\vec{p}_\theta:\Z^2 \to \R$ and $\vec{s}_\theta:\Z^2 \to \R$ can be defined as
\begin{align*}
\vec{p}_\theta(\gamma) & = \int_{(\theta,x)}^{(\theta,\gamma \cdot x)} \Omega/\eta \\
\vec{s}_\theta(\gamma) & = B_\theta(\gamma \cdot x) - B_\theta(x).
\end{align*}
where the integral is along the path of~$\Gamma$ formed of the geodesics of~$\bar{M}$ (with the same asymptotic direction~$\theta$) lying between the geodesic of asymptotic direction~$\theta$ passing through~$x$ and the geodesic of asymptotic direction~$\theta$ passing through~$\gamma \cdot x$, where $x$ is any point of~$\bar{M}$. In particular, these expressions do not depend on~$x \in \bar{M}$.
\end{proposition}

\begin{proof}
It follows from the $\Z^2$-invariance of the one-forms~$\zeta$ representing~$\Omega/\eta$ and the differential~$dB_\theta$ of the Busemann functions that the formulas defining~$\vec{p}_\theta$ and~$\vec{s}_\theta$ do not depend on~$x \in \bar{M}$. This also ensures that the relation~\eqref{eq:action} for the action of~$\gamma$ is satisfied.
\end{proof}
\end{comment}
%\forgotten



\section{Deforming Finsler two-tori}


Given a Finsler metric on the two-torus whose geodesic flow is conjugate to that of a flat Finsler metric, we show how to canonically connect the two metrics through Finsler metrics with the same dynamics. For this, we make use of the deformation of the geodesic foliation given by the Euclidean curve shortening flow, see Section~\ref{sec:convergenceCSF}, the invariance of the natural measure on the space of geodesics at the limit under this flow, see Section~\ref{sec:conjugate}, the construction of Finsler metrics with prescribed geodesics given by the Crofton formula, see Section~\ref{sec:prescribe}, and the relationship between conjugate geodesic flows and smooth Heber foliations; see Section~\ref{sec:smoothfoliation}. 

Let us start with the following observation.

\begin{prop}\label{Heber:Smooth}
Let $M_1=(\T^2,F_{1})$ and $M_2=(\T^2,F_{2})$ be two Finsler two-torus without conjugate points whose geodesic flows are conjugate. If the Heber foliation of $M_{1}$ is smooth then the Heber foliation of $M_{2}$ is also smooth.
\end{prop}

%\forget
\begin{comment}
\begin{proof}
Since the cogeodesic flow of $M_{1}$ and $M_{2}$ are conjugate, there exists a smooth $\Z^2$-equivariant diffeomorphism $\bar{h}:U^{*}\bar{M}_{1} \to U^{*}\bar{M}_{2}$ which intertwines the cogeodesic flows of~$\bar{M}_{1}$ and~$\bar{M}_{2}$. We can assume that the homomorphism $h_*:\pi_1(M_1) \to \pi_1(M_2)$ induced by the quotient map $h:U^*M_1 \to U^*M_2$ is the identity map; see Section~\ref{sec:conjugate}. Now, observe that the smooth diffeomorphism $\bar{h}:U^{*}\bar{M}_{1} \to U^{*}\bar{M}_{2}$ takes the Heber foliation of~$U^*\bar{M}_1$ to the Heber foliation of~$U^*\bar{M}_2$. More precisely, see~\eqref{eq:v}, the conjugacy $\bar{h}:U^{*}\bar{M}_{1} \to U^{*}\bar{M}_{2}$ takes the $\theta$-leaf~$\FF_\theta^1$ of~$U^*\bar{M}_1$ to the $\theta$-leaf~$\FF_\theta^2$ of~$U^*\bar{M}_2$, where $
\FF_\theta^i=\{ dB_\theta^i(x) \in U^*\bar{M}_i \mid x \in \bar{M}_i \}. $ Thus, if the Heber foliation of~$U^*\bar{M}_1$ is smooth, so is the Heber foliation of~$U^*\bar{M}_2$.
\end{proof}
\end{comment}
%\forgotten



%\forget
\begin{comment}
\begin{proof}
Let $M_{1}=(\T^{2},F_{1})$ and $M_{2}=(\T^{2},F_{2})$ be two Finsler metrics without conjugate points. Since the cogeodesic flow of $F_{1}$ and $F_{2}$ are conjugate, there exists a smooth diffeomorphism
\begin{align*}
\bar{h}:U^{*}\bar{M}_{1}&\to U^{*}\bar{M}_{2}\\
(x_{1},\xi_{1})&\to (x_{2},\xi_{2}).
\end{align*}
which intertwines the geodesic flows of~$\bar{M}_{1}$ and~$\bar{M}_{2}$. The smoothness of the Heber foliation of the metric $F_{1}$ implies that the Busemann map $B_{1}:\bar{M}_{1}\times S^{1}\to U^{*}\bar{M_{1}}$ associated to $F_{1}$ is a smooth diffeomorphism. In addition, using~\eqref{eq:identification1}, we get the following diagram
\[
\begin{tikzcd}
\bar{M}_{1}\times S^{1} \arrow{r}{B_{1}}\arrow{d}{f} &U^{*}\bar{M_{1}} \arrow{d}{\bar{h}} \\
\bar{M}_{2}\times S^{1} \arrow[swap]{r}{B_{2}} & U^{*}\bar{M}_{2}
\end{tikzcd}
\]
where $f$ is the map defined as follows
\begin{align*}
f:\bar{M}_{1}\times S^{1}&\to \bar{M}_{2}\times S^{1}\\
(x_{1},\theta_{1})&\to (x_{2},\theta_{2}),
\end{align*}
with $\theta_{1}$ and $\theta_{2}$ are the asymptotic directions of the geodesics induced by~$\xi_{1}$ and $\xi_{2}$. The asymptotic direction is preserved under the conjugacy of the cogeodesic flow, therefore we have $\theta_{1}=\theta_{2}$. By the smoothness of the maps $\bar{h}$ and $B_{1}$, we know that $x_{2}$ depends smoothly on $x_{1}$ and $\theta_{1}$. Thus $f$ is smooth and the same follows for $B_{2}$. In the same way, one may prove that the map $B_{2}^{-1}$ is smooth.
\end{proof}
\end{comment}
%\forgotten

\begin{proof}
Since the geodesic flow of~$M_1$ and $M_2$ are conjugate, there exists a smooth $\Z^2$-equivariant diffeomorphism $\bar{h}:U \bar{M}_1 \to U\bar{M}_2$ which intertwines the geodesic flows of~$\bar{M}_1$ and~$\bar{M}_2$. We can assume that the homomorphism $h_*:\pi_1(M_1) \to \pi_1(M_2)$ induced by the quotient map is the identity; see Observation~\ref{obs:id}. This implies that a unit vector generating a geodesic of~$\bar{M}_1$ with asymptotic direction~$\theta$ is sent under $\bar{h}:U\bar{M}_1 \to U\bar{M}_2$ to a unit vector generating a geodesic of~$\bar{M}_2$ with the same asymptotic direction~$\theta$. Thus, we have the following diagram

\[
\begin{tikzcd}[row sep=large]
(x_1,v_1) \in U\bar{M}_1 \arrow{d}{h_1} \arrow{r}{\bar{h}} & (x_2,v_2) \in U\bar{M}_2 \arrow{d}{h_2} \\
(x_1,\theta) \in \R^2 \times S^{1} & (x_2,\theta) \in \R^2 \times S^{1}
\end{tikzcd}
\]
where the two vertical maps~$h_1$ and~$h_2$ are the inverse Heber homeomorphism of~$M_1$ and~$M_2$; see Definition~\ref{def:heber}. By assumption, the first vertical map~$h_1$ is a diffeomorphism.

To prove that the Heber foliation of~$M_2$ is smooth, we need to show that the second vertical map~$h_2:U\bar{M}_2 \to \R^2 \times S^1$ is also a diffeomorphism. First, observe that the composition $U\bar{M}_2 \to U\bar{M}_1 \to \R^2 \times S^1 \to S^1$ taking~$(x_2,v_2)$ to~$\theta$ is smooth. This ensures that the vertical map~$h_2$ is smooth. Now, we only need to show that its inverse~$h_2^{-1}$ is also smooth.~Fix $(x_2,\theta) \in \R^2 \times S^1$. By Proposition~\ref{prop:identification}, we can assume that the horizontal and vertical curves of~$\bar{M}_1$ are geodesics. Without loss of generality, we can also assume that the direction~$\theta$ is not horizontal, otherwise we switch the horizontal and vertical axis of~$\bar{M}_1$ and~$\bar{M}_2$. Thus, the horizontal axis of~$\bar{M}_1$ transversely intersects every curve of the geodesic foliation with asymptotic direction~$\theta$ exactly once. As $x$ runs over the horizontal axis of~$\bar{M}_1$, the image under the conjugacy of the geodesic~$\gamma_v$ generated by the vector~$h_1^{-1}(x,\theta)= (x,v) \in U\bar{M}_1$ runs over the geodesic foliation of~$\bar{M}_2$ with asymptotic direction~$\theta$ (recall that the conjugacy preserves the asymptotic direction). More precisely, there is a unique vector $(x,v) \in U\bar{M}_1$ and a unique real~$t \in \R$ such that the basepoint of the vector~$h(\gamma_v'(t)) \in U\bar{M}_2$ agrees with~$x_2$. The vector~$v$ and the real~$t$ vary smoothly with~$(x_2,\theta)$ and the same goes for $(x_1,v_1)=\gamma_v'(t) \in U\bar{M}_1$ and $(x_2,v_2)=\bar{h}(x_1,v_1) \in U\bar{M}_2$. As previously noticed, we have $h_2(x_2,v_2) = (x_2,\theta)$ or equivalently $(x_2,v_2) = h_2^{-1}(x_2,\theta)$. It follows that the inverse of the second vertical map~$h_2:U\bar{M}_2 \to \R^2 \times S^1$ in the diagram is also smooth. Hence, the map~$h_2:U\bar{M}_2 \to \R^2 \times S^1$ is a diffeomorphism.
\end{proof}



We can now prove our main theorem.

\begin{theo}\label{theo:fin}
Let $M=(\T^2,F)$ be a Finsler two-torus whose geodesic flow is conjugate to the geodesic flow of a flat Finsler two-torus~$M_\diamond = (\T^2,F_\diamond)$. Then there exists a canonical deformation~$(F_t)_{t\,\geq\,0}$ of Finsler metrics on~$\T^2$ with $F_0=F$ such that
\begin{enumerate}
\item\label{theo13.2.1} the geodesic flow of~$F_t$ is conjugate to the geodesic flow of~$F_\diamond$;
\item the metric~$F_t$ converges to~$F_\diamond$ for the uniform convergence topology, up to isometry, as $t$ goes to infinity. \label{fin2}
\end{enumerate}
\end{theo}

%\forget
\begin{comment}
\begin{definition}\label{def:closing}
A measure $\lambda$ on $\Gamma$ satisfies the \emph{$F$-closing condition} if
\[
\lambda(\Gamma / \langle \gamma \rangle)=\mbox{length}_{F}(\gamma)
\]
for every closed geodesic~$\gamma$ on~$M$, where $\langle \gamma \rangle$ is the subgroup of~$\pi_1(\T^2)=\Z^2$ generated by the homotopy class of~$\gamma$.
\end{definition}

\begin{remark}
The closing condition is stable under convex combinations: if we have two measures $\lambda_{1}$ and $\lambda_{2}$ satisfying the $F$-closing condition, then any convex combination of $\lambda_{1}$ and $\lambda_{2}$ will also satisfy the $F$-closing condition.
\end{remark}
\end{comment}
%\forgotten


\begin{proof}
First, observe that since the geodesic flow of~$M$ is conjugate to that of the flat Finsler torus~$M_\diamond$, the metric~$F$ has no conjugate points by Lemma~\ref{lem:ncp} and its Heber foliation is smooth by Proposition~\ref{Heber:Smooth}. Denote by \mbox{$h:UM_\diamond \to UM$} the conjugacy between the geodesic flows of~$M_\diamond$ and~$M$. By Observation~\ref{obs:id}, we can assume that the conjugacy $h:UM_\diamond \to UM$ induces the identity map between~$\pi_1(M_\diamond)$ and~$\pi_1(M)$. In particular, this implies that the metrics~$F_\diamond$ and~$F$ have the same marked length spectrum. Let $\bh:\Gamma_\diamond \to \Gamma$ be the $\Z^2$-equivariant smooth diffeomorphism between the spaces of unparametrized oriented geodesics of~$\bar{M}_\diamond$ and~$\bar{M}$ induced by~$h:UM_\diamond \to UM$; see Definition~\ref{def:conjugacy}. Denote by $\lambda_\diamond$ and~$\lambda$ the smooth positive $\Z^2$-invariant measures on~$\Gamma_\diamond$ and~$\Gamma$ induced by~$F_\diamond$ and~$F$; see Remark~\ref{rem:lambda}. The measures
\[
\lambda_t = (1-t) \, \lambda + t \, \bh_*\lambda_\diamond
\]
with $0 \leq t \leq 1$ connecting~$\lambda$ to the push-forward~$\bh_* \lambda_\diamond$ of~$\lambda_\diamond$ by the $\Z^2$-equivariant smooth diffeomorphism $\bh:\Gamma_\diamond \to \Gamma$ are smooth positive $\Z^2$-invariant measures on~$\Gamma$. The family~$(\lambda_t)$ of measures on~$\Gamma$ induces a family~$(F_t)$ of Finsler metrics on~$\T^2$ joining~$F_0=F$ to~$F_1=F_{\bh_* \lambda_\diamond}$ with the same geodesic space~$\Gamma$; see the end of Section~\ref{sec:prescribe}. Since the metrics~$F_\diamond$ and~$F$ have the same marked length spectrum, the measures~$\bh_* \lambda_\diamond$ and~$\lambda$ on~$\Gamma$ satisfy the same $F$-closing condition; see Definition~\ref{def:closing} and Corollary~\ref{coro:closing}. By stability of the $F$-closing condition under convex combinations, see Remark~\ref{rem:closing}, the measures~$\lambda_t$ also satisfy the $F$-closing condition. It follows that the Finsler metrics~$F_t$ induced by the measures~$\lambda_t$ have the same marked length spectrum with the same geodesic space~$\Gamma$ as~$F$; see the end of Section~\ref{sec:prescribe} and Corollary~\ref{coro:closing}. By construction, these metrics have no conjugate points, see Observation~\ref{obs:ncp}, and their Heber foliations are smooth (since the metrics~$F_t$ and~$F$ have the same geodesics and the Heber foliation of~$F$ is smooth). Applying Theorem~\ref{theo:conjugacy}, we obtain that the Finsler metrics~$F_t$ have conjugate geodesic flows. 

Now, we need to deform the metric~$F_{\bh_* \lambda_\diamond}$ into~$F_\diamond$ through a family of Finsler metrics on~$\T^2$ with conjugate geodesic flows. Instead of deforming the measure~$\lambda$ induced by~$F$ on the geodesic space~$\Gamma$ as previously, we will deform the geodesics. Consider the family of~$\Z^2$-equivariant maps $f_t: \Gamma \to \Gamma_t$ induced by the curve shortening flow; see~\eqref{eq:ft}. The $\Z^2$-equivariant maps $ f_t \circ \bh: \Gamma_\diamond \to \Gamma_t $ are smooth for~$t \in [0,\infty)$ and measurable for $t=\infty$; see Proposition~\ref{prop:pinfty}. Thus, the push-forward measures $\lambda_t=(f_t \circ \bh)_* \lambda_\diamond$ on~$\Gamma_t$ are smooth, positive and~$\Z^2$-invariant for every $t \in [0,\infty]$, even for $t=\infty$ where $(f_\infty \circ \bh)_* \lambda_\diamond = \lambda_\diamond$ on~$\Gamma_\infty = \Gamma_\diamond$ by Proposition~\ref{prop:measure}. The family~$(\lambda_t)$ of measures on~$\Gamma_t$ induces a family~$(F_t)$ of Finsler metrics on~$\T^2$ starting at~$F_{\bh_* \lambda_0}$ and ending at~$F_\diamond$; see the end of Section~\ref{sec:prescribe}. Indeed, by the Crofton formula, see Theorem~\ref{theo:crofton}, the Finsler metric corresponding to the measure~$\lambda_t=(f_t \circ \bh)_* \lambda_\diamond$ on~$\Gamma_t$ when $t=\infty$ is the flat Finsler metric~$F_\diamond$. Now, by $\Z^2$-invariance of the maps $f_t:\Gamma \to \Gamma_t$ for $t \in [0,\infty]$, we have
\[
\lambda_t(\Gamma_t/\langle \alpha \rangle) = \bh_* \lambda_\diamond(\Gamma/\langle \alpha \rangle) = \lambda(\Gamma/\langle \alpha \rangle)
\]
for every $\alpha \in \Z^2$ (recall that the measures~$\bh_* \lambda_\diamond$ and~$\lambda$ on~$\Gamma$ satisfy the same $F$-closing condition). It follows that the Finsler metrics~$F_t$ with $t \in [0,\infty]$ have the same marked length spectrum by Corollary~\ref{coro:closing}. Furthermore, since their geodesic space is given by the deformation~$\Gamma_t$ of~$\Gamma$ under the curve shortening flow, the metrics~$F_t$ have no conjugate points, see Observation~\ref{obs:ncp}, and their Heber foliation is smooth (even for~$t=\infty$); see Observation~\ref{obs:foliation-deformation}. Applying Theorem~\ref{theo:conjugacy} as previously, we obtain that the Finsler metrics~$F_t$ have conjugate geodesic flows for~$t \in [0,\infty]$.


In order to conclude the proof of Theorem~\ref{theo:fin}, we need to show the point~\eqref{fin2}.


Fix a compact convex subset~$K$ of~$\R^2$. Let $\Lambda_K = [0,\infty] \times K \times K$ be the compact set where every point~$(t,x,y) \in [0,\infty] \times K \times K$ identifies with the geodesic arc~$c_t(x,y)$ of~$\bar{M}_t$ joining~$x$ to~$y$ (with $\bar{M}_\infty = \R^2$). Note that $c_t(x,y)$ lies in a curve of~$\Gamma_t$.

Consider the map $\varphi: \Lambda_K \times \Gamma_\diamond \to \R_+$ defined as
\[
\varphi(t,x,y,\gamma) = \frac{1}{4} \, \#\big(f_t \circ \bh(\gamma) \cap c_t(x,y)\big)
\]
for every $(t,x,y) \in \Lambda_K$ and $\gamma \in \Gamma_\diamond$. Observe that $f_t \circ \bh(\gamma)$ and~$c_t(x,y)$ intersect at most once as long as $\gamma$ is different from the line whose image $f_t \circ \bh(\gamma)$ under the curve shortening flow passes through~$x$ and~$y$. Thus, the map $\varphi : \Lambda_K \times \Gamma_\diamond \to \R_+$ is bounded by~$\frac{1}{4}$ almost everywhere.

Let $w$ be the maximal width of a geodesic of~$\bar{M}$; see Theorem~\ref{theo:hedlund}. The space~$\Gamma_\diamond(K_w)$ of lines in~$\Gamma_\diamond$ whose image under $\bh:\Gamma_\diamond \to \Gamma$ intersects the closed Euclidean $(w+1)$-neighborhood~$K_w$ of~$K \subseteq \R^2$ is compact. By the definition of $\Gamma_{\diamond}(K_{w})$, the image under $\bh:\Gamma_{\diamond}\to\Gamma$ of a line $\gamma\in\Gamma_{\diamond}\setminus\Gamma_{\diamond}(K_{w})$ is a geodesic of $\bar{M}$ at Euclidean distance greater than $w+1$ from $K$. It follows from the definition of $w$ that this image lies in a closed strip disjoint from~$K$ and so are its images~$f_t \circ \bh(\gamma)$ and~$f_\infty \circ \bh(\gamma)$ under the curve shortening flow. By Crofton's formula, see Theorem~\ref{theo:crofton}, we obtain
\begin{equation}\label{eq:integral}
d_{\bar{F}_t}(x,y) = \int_{\gamma\,\in\,\Gamma_\diamond(K_w)} \varphi(t,x,y,\gamma) \, d\lambda_\diamond
\end{equation}
for every $(t,x,y) \in \Lambda_K$.


Let $x_0,y_0 \in K$. By Theorem~\ref{theo:affine}, for almost every $\gamma \in \Gamma_\diamond(K_w)$ (actually, as long as the line~$f_\infty \circ \bh(\gamma)$ does not pass through~$x_0$ or~$y_0$), the map $\varphi(.,\gamma): \Lambda_K \to \R_+$ is continuous at~$(\infty,x_0,y_0)$. Indeed, for $(x,y)$ close to $(x_{0},y_{0})$ and for $t$ large enough, the arc $c_{t}(x,y)$ is $C^{k}$-close to the segment $c_{\infty}(x_{0},y_{0})=[x_{0},y_{0}]$. Similarly, for $t$ large enough, the curve $f_{t}\circ\bh(\gamma)$ is close to the line $f_{\infty}\circ\bh(\gamma)$. Thus, if the line $f_{\infty}\circ\bh(\gamma)$ transversely intersects the interior of the segment $[x_{0},y_{0}]$, then the curve $f_{t}\circ\bh(\gamma)$ intersects the arc $c_{t}(x,y)$ at a unique point. Similarly, if the line $f_{\infty}\circ\bh(\gamma)$ does not intersect the segment $[x_{0},y_{0}]$, then the curve $f_{t}\circ\bh(\gamma)$ does not intersect the arc $c_{t}(x,y)$ either. Therefore, the map $\varphi(\cdot,\gamma):\Lambda_{K}\to\R_{+}$ is constant (equal to~$0$ or~$1$) in a neighborhood of each point $(\infty,x_{0},y_{0})$, where $x_{0}$ and $y_{0}$ do not belong to the line $f_{\infty}\circ\bh(\gamma)$. It follows from Lebesgue's dominated convergence theorem that the integral~\eqref{eq:integral} converges to
\begin{align*}
\frac{1}{4} \, \int_{\gamma\,\in\,\Gamma_\diamond(K_w)} \#\big(f_\infty \circ \bh(\gamma) \cap [x_0,y_0]\big) \, d\lambda_\diamond & = \frac{1}{4} \, \int_{\gamma\,\in\,\Gamma_\diamond} \#(\gamma \cap [x_0,y_0]) \, d[(f_\infty \circ \bh)_* \lambda_\diamond] \\
& = d_{\bar{F}_\diamond}(x_0,y_0)
\end{align*}
as $(t,x,y) \to (\infty,x_0,y_0)$, where the last inequality comes from the measure invariance property of Proposition~\ref{prop:measure} and Crofton's formula. Hence, the distance~$d_{\bar{F}_t}$ uniformly converges to~$d_{\bar{F}_\infty}$ on any compact set as $t$ goes to infinity.

To show that the metric~$F_t$ converges to~$F_\diamond$ for the uniform convergence topology, it is enough to prove that the diameter of~$M_t$ is uniformly bounded. To this end, fix two (transverse) geodesic foliation~$\FF'_\diamond$ and~$\FF''_\diamond$ of~$M_\diamond$ by simple closed geodesics, see Proposition~\ref{prop:Hfoliation}.\eqref{geod1}, and denote by~$\ell'$ and~$\ell''$ their lengths. The corresponding geodesic foliations~$\FF'_t$ and~$\FF''_t$ of~$M_t$ have the same length (recall that $M_t$ has no conjugate points with the same marked length spectrum as~$M_\diamond$). Observe that any pair of points on~$M_t$ can be joined by a path formed of an arc lying a geodesic of~$\FF'_t$ followed by an arc lying a geodesic of~$\FF''_t$. This implies that the diameter of~$M_t$ is uniformly bounded by~$\ell'+\ell''$.
\end{proof}



%\forget
\begin{comment}
\begin{lemma}
The metric~$F_t$ converges to~$F_\diamond$ for the uniform convergence topology.
\end{lemma}

\begin{proof}
Let $x$, $y\in \bar{M_{t}}$ and $\gamma$ a geodesic in $\Gamma_{t}$ passing through $x$ and $y$. Let us look at all the curves in $\Gamma_{t}$ intersecting $\gamma$ and compare it with the flat Finsler metric $F_{\diamond}$ after deformation. We know that there exists $\lambda_{\diamond}$ and $\lambda_{t}$ smooth positive $\Z^{2}$-invariance measures on $\Gamma_{\diamond}$ and $\Gamma_{t}$ and a smooth $\Z^{2}$-equivariant map $f_{t}\circ\bh:\Gamma_{\diamond}\to\Gamma_{t}$ such that $\lambda_{t}=(f_{t}\circ\bh)_{*}\lambda_{\diamond}$.\\
For irrational directions, we know using Theorem~\ref{theo:unifwidth}, that for every geodesic passing through $x$ and $y$ of irrational direction $\theta$ of $\bar{M}_{t}$, so in particular for $\gamma$, and for every $\varepsilon>0$, there exists $t_{0}>0$ such that for every $t\geq t_{0}$, the width of $\gamma$ satisfies $\mbox{width}(\gamma)<\varepsilon$. Thus, by taking the limit, the image of $\gamma$ in $\Gamma_{\diamond}$ will also be of small width.\\
By definition, the distance between $x$ and $y$ can be measured using the number of points of all the curves in $\R^{2}$ intersecting $\gamma$, thus $d_{t}(x,y)=\frac{1}{4}\lambda_{t}(\Gamma_{t} / \langle \gamma \rangle)$, where $\langle \gamma \rangle$ denotes the subgroup of~$\pi_1(\T^2)=\Z^2$ generated by the homotopy class of the curve $\gamma$. Hence, by taking the limit, we will get the same number of intersections between the image of $\gamma$, which is contained in a strip of very small width, and the curves in $\Gamma_{\diamond}$. Thus, $d_{t}(x,y)\to d_{\diamond}(x,y)$ and the result follows.
\end{proof}
\end{comment}
%\forgotten


%\forget
\begin{comment}
Consider a Finsler metric $F$ without conjugate points on $\T^{2}$. Let $\rho$ be a smooth free $\R$-action on $U^{*}\T^{2}$ whose orbits are transverse to the contact structure given by the kernel of $\alpha$ on $U^{*}\T^{2}$ and project to embeddings of $\R$ into $\R^{2}$. The action $\rho$ induces a Legendrian action $\bar{\rho}$ defined as follows: Set the map
\[
\clR:U\T^{2}\rightarrow U^{*}\T^{2}
\]
sending every unit vector $v\in U_{x}\T^{2}$ to the unique unit co-vector $\xi\in U_{x}^{*}\T^{2}$ such that $\xi(v)=0$. The map $\clR$ and its restrictions $\clR_{x}:U_{x}\T^{2}\rightarrow U_{x}^{*}\T^{2}$ are diffeomorphisms. Thus, the free $\R$-action $\rho$ on $U^{*}\T^{2}$ induces a free $\R$-action $\bar{\rho}$ on $U^{*}\T^{2}$ defined as follows
\[
\bar{\rho}(t)=\Tau^{-1}\circ\rho(t)\circ\Tau
\]
for every $t\in\R$, where $\Tau$ is the composition map $\Tau:\LL\circ\clR^{-1}:U^{*}\T^{2}\rightarrow U^{*}\T^{2}$. Consider the fibration
\[
q_{\bar{\rho}}:U^{*}\T^{2}\rightarrow \Gamma_{\bar{\rho}}
\]
induced by the free $\R$-action $\bar{\rho}$, where $\Gamma_{\bar{\rho}}=U^{*}\T^{2}/\bar{\rho}$ and $\Gamma_{\rho}=U^{*}\T^{2}/\rho$. Since both actions $\rho$ and $\bar{\rho}$ are conjugate, then we have $\Gamma_{\bar{\rho}}\simeq \Gamma_{\rho}$.\\

We shall need the following results about the fibration $q_{\bar{\rho}}$

\begin{lemma}
The fibers of $q_{\bar{\rho}}:U^{*}\T^{2}\rightarrow\Gamma_{\bar{\rho}}$ are Legendrian with respect to the contact structure induced by $\alpha$ on $U^{*}\T^{2}$.
\end{lemma}

\begin{proof}
For every $\gamma\in\Gamma_{\bar{\rho}}$, let $\xi\in q_{\bar{\rho}}^{-1}(\gamma)$ and $X\in T_{\xi} q_{\bar{\rho}}^{-1}(\gamma)$. The vector $d\pi_{\xi}(X)$ is tangent to $\gamma$ at $\pi(\xi)$. By the definition of $q_{\bar{\rho}}$, the vectors that are tangent to $\gamma$ and based at $\pi(\xi)$ lie in the kernel of $\xi$. Hence, $\xi(d\pi_{\xi}(X))=0$, hence $\alpha_{\xi}(X)=0$, thus we get the result.
\end{proof}

\begin{lemma}
The product map $\phi_{\bar{\rho}}=\pi\times q_{\bar{\rho}}:U^{*}_{0}\T^{2}\rightarrow\T^{2}\times\Gamma_{\bar{\rho}}$ is an embedding.
\end{lemma}

\begin{proof}
Let us start by proving that $\phi_{\bar{\rho}}$ is an immersion. Consider $\xi\in U^{*}\T^{2}$ and $X\in T_{\xi} q_{\bar{\rho}}^{-1}(\gamma)$ such that $d\phi_{\bar{\rho}}(\xi)(X)=0$. Thus, the vector $X$ lies in the kernel of the differential of the fibration $q_{\bar{\rho}}$ at $\xi$. It follows that $V$ is tangent to the $\bar{\rho}$-orbit of $U^{*}\T^{2}$ at $\xi$. One deduces from the fact that $d\pi_{\xi}(V)=0$ that $V=0$. Thus, $d\phi_{\bar{\rho}}$ is injective and $\phi_{\bar{\rho}}$ is an immersion. \\
Let us prove that $\phi_{\bar{\rho}}$ is injective to conclude that it gives rise to an embedding. Let $\xi_{1}$, $\xi_{2}$ be in $U^{*}\T^{2}$ such that $\phi_{\bar{\rho}}(\xi_{1})=\phi_{\bar{\rho}}(\xi_{2})$, then $\pi(\xi_{1})=\pi(\xi_{2})$ and $q_{\bar{\rho}}(\xi_{1})=q_{\bar{\rho}}(\xi_{2})$. Thus, we get that $\xi_{1}$ and $\xi_{2}$ are based at the same point $x$ of $\T^{2}$. In addition, using the fact that the projections of the $\bar{\rho}$-orbits to $\T^{2}$ are embeddings of $\R$ into $\T^{2}$, then for every $\gamma\in\Gamma_{\bar{\rho}}$ such that $x\in\pi(q_{\bar{\rho}}^{-1}(\gamma))$, there exists a unique $\xi\in U^{*}_{x}\T^{2}$ such that $q_{\bar{\rho}}(\xi)=\gamma$. Hence, this implies that $\xi_{1}=\xi_{2}$.
\end{proof}


On the other hand, we can apply the results of~\cite{AB} stating that a non-vanishing area-form on $\Gamma_{\bar{\rho}}$ gives rise to a Finsler metric on $\T^{2}$ via the Crofton formula.

\begin{theorem}
There exists a unique Finsler metric $F_{\rho}$ on $\T^{2}$ satisfying the Crofton formula
\begin{equation}
\mbox{length}_{F_{\rho}}(c)=\frac{1}{4}\int_{\gamma\in\Gamma_{\bar{\rho}}}\#(\gamma\cap c)|\lambda_{\bar{\rho}}|
\end{equation}
for any smooth curve $c$ on $\T^{2}$.\\
In addition, the Finsler metric $F_{\rho}$ can be represented in the following way: for every $x\in\T^{2}$, there exists a unique non-vanishing one-form $\alpha_{x}$ on $U^{*}_{x}\T^{2}$ such that
\begin{equation}
F_{\rho}(x,v)=\int_{\xi\in U^{*}_{x}\T^{2}}|\xi(v)|\alpha_{x}
\end{equation}
for every $v\in T_{x}\T^{2}$, where the one-form $\alpha_{x}$, depends smoothly on $x$,
\end{theorem}
\begin{remark}
The one-form $\alpha$ varies continuously with $\rho$. Hence, using the expression of $F_{\rho}$, it follows that $F_{\rho}$ varies continuously with $\rho$ too.
\end{remark}
\end{comment}
%\forgotten



\bibliography{nakhle}
\end{document}
